% Lesson 15 — Vocabulary: Words and expressions related to buying and selling articles/foodstuff at a nearby store/market — Anglais 1ère A — Cours signé OnBuch+
% Programme officiel MINESEC. Compiler avec : tectonic ENA16-vocabulary-words-and-expressions-related-to-buying-and.tex
\newcommand{\DOCMATIERE}{Anglais}
\newcommand{\DOCNIVEAU}{1ère A}
\newcommand{\DOCMODULE}{Chapter 2 : Economic Life and Occupations}
\newcommand{\DOCLECON}{ENA16}
\newcommand{\DOCTITRE}{Lesson 15 — Vocabulary: Words and expressions related to buying and selling articles/foodstuff at a nearby store/market}
% OnBuch+ — préambule « cours signés » ENRICHI (compilé avec tectonic / XeLaTeX)
% Système visuel « Pop » d'OnBuch+ : fond crème (jamais blanc), bordures encre
% épaisses, ombres « dures » décalées, titres Archivo Black en capitales.
% Version MATHÉMATIQUES : graphes (pgfplots/tikz), boîtes pédagogiques
% (définition, propriété, méthode, attention, le savais-tu, exercice résolu,
% exercice, TP, bilan), page de garde et signature OnBuch+ ancrée en bas.
%
% Macros attendues dans main.tex AVANT \input{preamble} :
%   \DOCMATIERE  \DOCNIVEAU  \DOCMODULE  \DOCLECON (numéro)  \DOCTITRE
\documentclass[11pt]{article}

\usepackage{fontspec}
\usepackage[english,french]{babel} % français = langue principale ; l'anglais via \eng{..} / textebox
\usepackage[a4paper,margin=2cm,top=3.4cm,bottom=2.8cm,headheight=1.4cm,footskip=1.3cm]{geometry}
\usepackage{amsmath,amssymb,amsfonts}
\usepackage{mathtools} % \xmapsto, \coloneqq... souvent utilisés par les modèles
\usepackage{cancel} % simplifications barrées (bilans de forces)
% \abs{z} (module, valeur absolue) et \norm{u} (norme), utilisés par les modèles
\providecommand{\abs}{}\let\abs\relax\DeclarePairedDelimiter{\abs}{\lvert}{\rvert}
\providecommand{\norm}{}\let\norm\relax\DeclarePairedDelimiter{\norm}{\lVert}{\rVert}
\usepackage[table]{xcolor} % [table] : fournit aussi \rowcolors (alternance de lignes)
\usepackage{pagecolor}
\usepackage{tikz}
\usetikzlibrary{arrows.meta,positioning,calc,shapes.geometric,shapes.misc,shapes.arrows,decorations.pathmorphing,decorations.pathreplacing,decorations.markings,patterns,shadows,mindmap,trees,fit,backgrounds,matrix,intersections,angles,quotes}
\usepackage{pgfplots}
\usepackage[version=4]{mhchem} % équations nucléaires \ce{^{235}_{92}U}
\usepackage[european,siunitx]{circuitikz} % schémas électriques
\pgfplotsset{compat=1.18}
\usepgfplotslibrary{fillbetween,groupplots} % aires entre courbes (calcul intégral)
\usepackage{enumitem}
\usepackage{fancyhdr}
% [most] charge toutes les bibliothèques tcolorbox (listings, minted, raster,
% documentation...) et gonfle inutilement la mémoire TeX sur un document long
% (~300 boîtes) : on ne charge que ce qui est réellement utilisé.
\usepackage[skins,breakable]{tcolorbox}
\usepackage{graphicx}
\usepackage{colortbl}
\usepackage{booktabs}
\usepackage{tabularx}
\usepackage{diagbox} % case d'en-tête diagonale des tableaux à double entrée
% Colonnes extensibles centrée / gauche / droite (C, L, R), souvent utilisées par les modèles
\newcolumntype{C}{>{\centering\arraybackslash}X}
\newcolumntype{L}{>{\raggedright\arraybackslash}X}
\newcolumntype{R}{>{\raggedleft\arraybackslash}X}
\usepackage{array}
\usepackage{multicol}
\usepackage{multirow}
\usepackage{siunitx}
% parse-numbers=false : accepte aussi \SI{2,5 \times 10^{-3}}{...}
\sisetup{locale=FR,output-decimal-marker={,},group-minimum-digits=4,parse-numbers=false}
\DeclareSIUnit\ohm{\ensuremath{\symup{\Omega}}} % U+2126 absent de la police
\DeclareSIUnit\division{div} % réglages d'oscilloscope (ms/div, V/div)
\DeclareSIUnit\tr{tr} % tours (vitesse de rotation, ex. \SI{1500}{\tr\per\minute})
\DeclareSIUnit\molar{M} % concentration molaire (ex. \SI{3}{\molar}, \SI{1}{\micro\molar})
\DeclareSIUnit\an{an} % année (le modèle confond parfois avec \year, un compteur TeX) — ex. \SI{5}{\kilo\meter\per\an}
\DeclareSIUnit\FCFA{FCFA} % franc CFA (ex. \SI{500}{\FCFA\per\kilo\gram})
\DeclareSIUnit\degreeBrix{\textdegree Brix} % teneur en sucre d'un sirop/jus (réfractométrie)
\DeclareSIUnit\month{mois} % le modèle utilise parfois \month, un compteur TeX, comme unité de durée
\usepackage{qrcode}
% \degree : commande fréquemment utilisée par les modèles pour un angle ou une
% température en dehors de \SI{}{} (non fournie par les paquets chargés).
\providecommand{\degree}{^{\circ}}
% \qed (fin de démonstration), utilisé par les modèles sans amsthm
\providecommand{\qed}{\hfill$\square$}
% \barre : double barre des valeurs interdites dans un tableau de variations (tkz-tab)
\providecommand{\barre}{\ensuremath{\|}}
% environnement proof (amsthm non chargé) utilisé par les modèles
\ifdefined\proof\else\newenvironment{proof}[1][Démonstration]{\par\noindent\textit{#1.}\ }{\qed\par}\fi
\usepackage{lastpage}
\usepackage{hyperref}
\hypersetup{hidelinks}

% ============================================================
% Palette « Pop » OnBuch+ — mêmes hex que la classe P (plus_ui.dart)
% ============================================================
\definecolor{poporange}{HTML}{F59321}
\definecolor{popdark}{HTML}{E07A0C}
\definecolor{popcream}{HTML}{FFF6E8}
\definecolor{popink}{HTML}{1C1714}
\definecolor{poppurple}{HTML}{7A5AE0}
\definecolor{popgreen}{HTML}{2E9E6A}
\definecolor{poppink}{HTML}{E0457B}
\definecolor{popblue}{HTML}{2F7DE1}
\definecolor{popgold}{HTML}{C98A12}
\definecolor{popmuted}{HTML}{8A7F76}
\definecolor{popline}{HTML}{EADFD2}
\definecolor{popgray}{HTML}{8A7F76}
\colorlet{popgrey}{popgray}
% Alias tolérants (le modèle nomme parfois les couleurs différemment)
\colorlet{popwhite}{white}
\colorlet{popblack}{popink}
\colorlet{popred}{poppink}
\colorlet{popyellow}{popgold}
% Teintes claires (fonds de tableaux/encadrés) — le modèle nomme parfois
% les couleurs avec un suffixe L (ex. popblueL) sans les avoir définies.
\colorlet{poporangeL}{poporange!20}
\colorlet{popdarkL}{popdark!20}
\colorlet{poppurpleL}{poppurple!20}
\colorlet{popgreenL}{popgreen!20}
\colorlet{poppinkL}{poppink!20}
\colorlet{popblueL}{popblue!20}
\colorlet{popgoldL}{popgold!20}
\colorlet{popredL}{popred!20}
\colorlet{popgrayL}{popgray!20}
% Teintes claires pour fonds de tableaux / zones
\colorlet{popblueL}{popblue!10!white}
\colorlet{poporangeL}{poporange!14!white}
\colorlet{popgreenL}{popgreen!12!white}
\colorlet{poppurpleL}{poppurple!12!white}
\colorlet{poppinkL}{poppink!12!white}
\colorlet{popgoldL}{popgold!14!white}

\pagecolor{popcream}

% ============================================================
% Typographie
% ============================================================
\defaultfontfeatures{Ligatures=TeX}
\setmainfont{PlusJakartaSans-Medium.ttf}[Path=../../../fonts/,BoldFont=PlusJakartaSans-Bold.ttf]
% Symboles, API et emojis absents de la police principale : repli sur DejaVu Sans (caractères actifs)
\newfontface\symfont{DejaVuSans.ttf}[Path=../../../fonts/]
\makeatletter
\newcommand\symactive[1]{\catcode#1=13 \begingroup\lccode`\~=#1 \lowercase{\endgroup\def~}{{\symfont\char#1}}}
\newcommand\symblank[1]{\catcode#1=13 \begingroup\lccode`\~=#1 \lowercase{\endgroup\def~}{}}
\newcommand\symhyph[1]{\catcode#1=13 \begingroup\lccode`\~=#1 \lowercase{\endgroup\def~}{\mbox{-}}}
\newcount\symc
\newcommand\symrange[2]{\symc=#1 \@whilenum\symc<#2 \do{\expandafter\symactive\expandafter{\the\symc}\advance\symc by 1 }}
\newcommand\symblankrange[2]{\symc=#1 \@whilenum\symc<#2 \do{\expandafter\symblank\expandafter{\the\symc}\advance\symc by 1 }}
\makeatother
\symrange{"0250}{"02FF}\symactive{"2713}\symactive{"2714}\symactive{"2717}\symactive{"2718}\symactive{"2610}\symactive{"2611}\symactive{"2612}
\symhyph{"2011}\symhyph{"2E1A}\symblankrange{"1F300}{"1FAFF}\symblank{"FE0F}
% Maths en sans-serif (Fira Math) avec chiffres et lettres droites de Plus
% Jakarta, pour que les formules se fondent dans le texte.
\usepackage[math-style=ISO,bold-style=ISO]{unicode-math}
\setmathfont{FiraMath-Regular.otf}[Path=../../../fonts/]
\setmathfont{PlusJakartaSans-Medium.ttf}[Path=../../../fonts/,range={up/num,up/{Latin,latin},\mathup}]
\usepackage{icomma} % virgule décimale sans espace en mode math
% Caractères mathématiques Unicode absents des polices de texte (Plus Jakarta,
% Archivo Black) : rendus par leur équivalent mathématique, en texte comme en
% formule — sinon ils s'affichent en carré vide (ex. « Arithmétique dans ℤ »).
\usepackage{newunicodechar}
\newunicodechar{ℝ}{\ensuremath{\mathbb{R}}}
\newunicodechar{ℤ}{\ensuremath{\mathbb{Z}}}
\newunicodechar{ℂ}{\ensuremath{\mathbb{C}}}
\newunicodechar{ℕ}{\ensuremath{\mathbb{N}}}
\newunicodechar{ℚ}{\ensuremath{\mathbb{Q}}}
\newunicodechar{𝔻}{\ensuremath{\mathbb{D}}}
\newunicodechar{α}{\ensuremath{\alpha}}
\newunicodechar{β}{\ensuremath{\beta}}
\newunicodechar{γ}{\ensuremath{\gamma}}
\newunicodechar{δ}{\ensuremath{\delta}}
\newunicodechar{ε}{\ensuremath{\varepsilon}}
\newunicodechar{θ}{\ensuremath{\theta}}
\newunicodechar{λ}{\ensuremath{\lambda}}
\newunicodechar{μ}{\ensuremath{\mu}}
\newunicodechar{σ}{\ensuremath{\sigma}}
\newunicodechar{τ}{\ensuremath{\tau}}
\newunicodechar{φ}{\ensuremath{\varphi}}
\newunicodechar{ψ}{\ensuremath{\psi}}
\newunicodechar{ω}{\ensuremath{\omega}}
\newunicodechar{Σ}{\ensuremath{\Sigma}}
% majuscules produites par \MakeUppercase dans les titres (α-aminés -> Α)
\newunicodechar{Α}{\ensuremath{\alpha}}
\newunicodechar{Β}{\ensuremath{\beta}}
\newunicodechar{Γ}{\ensuremath{\Gamma}}
\newunicodechar{Δ}{\ensuremath{\Delta}}
\newunicodechar{Ε}{\ensuremath{\varepsilon}}
\newunicodechar{Θ}{\ensuremath{\Theta}}
\newunicodechar{Λ}{\ensuremath{\Lambda}}
\newunicodechar{Μ}{\ensuremath{\mu}}
\newunicodechar{Τ}{\ensuremath{\tau}}
\newunicodechar{Φ}{\ensuremath{\Phi}}
\newunicodechar{Ψ}{\ensuremath{\Psi}}
\newunicodechar{Ω}{\ensuremath{\Omega}}
\newunicodechar{↦}{\ensuremath{\mapsto}}
\newunicodechar{∈}{\ensuremath{\in}}
\newunicodechar{∉}{\ensuremath{\notin}}
\newunicodechar{∧}{\ensuremath{\wedge}}
\newunicodechar{⇔}{\ensuremath{\Leftrightarrow}}
\newunicodechar{⟺}{\ensuremath{\Longleftrightarrow}}
\newunicodechar{⇒}{\ensuremath{\Rightarrow}}
\newunicodechar{⟹}{\ensuremath{\Longrightarrow}}
\newunicodechar{∘}{\ensuremath{\circ}}
\newunicodechar{∪}{\ensuremath{\cup}}
\newunicodechar{∩}{\ensuremath{\cap}}
\newunicodechar{≡}{\ensuremath{\equiv}}
\newunicodechar{⊂}{\ensuremath{\subset}}
\newunicodechar{⊥}{\ensuremath{\perp}}
\newunicodechar{∃}{\ensuremath{\exists}}
\newunicodechar{∀}{\ensuremath{\forall}}
\newunicodechar{∛}{\ensuremath{\sqrt[3]{\,}}}
\newunicodechar{∜}{\ensuremath{\sqrt[4]{\,}}}
\newunicodechar{⃗}{\ensuremath{{}^{\to}}}
\newunicodechar{ⁿ}{\ifmmode^{n}\else\textsuperscript{\ensuremath{n}}\fi}
\newunicodechar{ˣ}{\ifmmode^{x}\else\textsuperscript{\ensuremath{x}}\fi}
\newunicodechar{ᵇ}{\ifmmode^{b}\else\textsuperscript{\ensuremath{b}}\fi}
\newunicodechar{ʳ}{\ifmmode^{r}\else\textsuperscript{\ensuremath{r}}\fi}
\newunicodechar{ᵘ}{\ifmmode^{u}\else\textsuperscript{\ensuremath{u}}\fi}
\newunicodechar{ʸ}{\ifmmode^{y}\else\textsuperscript{\ensuremath{y}}\fi}
\newunicodechar{ᵃ}{\ifmmode^{a}\else\textsuperscript{\ensuremath{a}}\fi}
\newunicodechar{ᵉ}{\ifmmode^{e}\else\textsuperscript{\ensuremath{e}}\fi}
\newunicodechar{ᵏ}{\ifmmode^{k}\else\textsuperscript{\ensuremath{k}}\fi}
\newunicodechar{ᵖ}{\ifmmode^{p}\else\textsuperscript{\ensuremath{p}}\fi}
\newunicodechar{ᴺ}{\ifmmode^{N}\else\textsuperscript{\ensuremath{N}}\fi}
\newunicodechar{ᶜ}{\ifmmode^{c}\else\textsuperscript{\ensuremath{c}}\fi}
\newunicodechar{ʰ}{\ifmmode^{h}\else\textsuperscript{\ensuremath{h}}\fi}
\newunicodechar{ᵠ}{\ifmmode^{q}\else\textsuperscript{\ensuremath{q}}\fi}
\newunicodechar{ₙ}{\ifmmode_{n}\else\textsubscript{\ensuremath{n}}\fi}
\newunicodechar{ₐ}{\ifmmode_{a}\else\textsubscript{\ensuremath{a}}\fi}
\newunicodechar{ᵢ}{\ifmmode_{i}\else\textsubscript{\ensuremath{i}}\fi}
\newunicodechar{ᵣ}{\ifmmode_{r}\else\textsubscript{\ensuremath{r}}\fi}
\newunicodechar{ₖ}{\ifmmode_{k}\else\textsubscript{\ensuremath{k}}\fi}
\newunicodechar{ᵦ}{\ifmmode_{\beta}\else\textsubscript{\ensuremath{\beta}}\fi}
\newunicodechar{ₓ}{\ifmmode_{x}\else\textsubscript{\ensuremath{x}}\fi}
\newfontfamily\popdisplay{ArchivoBlack-Regular.ttf}[Path=../../../fonts/]
\newfontfamily\poplogo{SpaceGrotesk-Bold.ttf}[Path=../../../fonts/]
\newfontfamily\popsemi{PlusJakartaSans-SemiBold.ttf}[Path=../../../fonts/]
\newfontfamily\popxbold{PlusJakartaSans-ExtraBold.ttf}[Path=../../../fonts/]
\newfontfamily\popxboldls{PlusJakartaSans-ExtraBold.ttf}[Path=../../../fonts/,LetterSpace=3.0]

\setlist[enumerate]{leftmargin=1.7em,itemsep=5pt,topsep=4pt}
\setlist[itemize]{leftmargin=1.7em,itemsep=4pt,topsep=4pt,label={\color{poporange}\textbullet}}
\setlength{\parindent}{0pt}
\setlength{\parskip}{5pt}
\renewcommand{\arraystretch}{1.25}

% pgfplots : style Pop par défaut
\pgfplotsset{
  every axis/.append style={axis line style={popink,line width=1pt},tick style={popink},
    grid style={popline,line width=0.5pt},label style={font=\small},tick label style={font=\footnotesize}},
  popcurve/.style={poporange,line width=1.8pt,no markers,smooth},
}
% Même style disponible dans un \draw TikZ simple (hors pgfplots)
\tikzset{popcurve/.style={poporange,line width=1.8pt}}
\tikzset{popgrid/.style={popline,very thin}}
\colorlet{popcurve}{poporange} % le nom du style est parfois utilisé comme couleur

% ============================================================
% Pill / squiggle
% ============================================================
\newcommand{\poppill}[3]{%
  \tikz[baseline=-0.55ex]\node[draw=popink,line width=1.2pt,rounded corners=2.6pt,%
    fill=#2,inner xsep=6.5pt,inner ysep=3pt]{%
    \popxboldls\fontsize{7}{7}\selectfont\color{#3}\MakeUppercase{#1}};%
}
\newcommand{\popsquiggle}[1]{%
  \tikz[baseline=-0.4ex]{
    \draw[#1,line width=2.6pt,line cap=round]
      plot [smooth] coordinates {(0,0.09) (0.34,0.01) (0.68,0.21) (1.02,0.01) (1.36,0.10)};
  }%
}
% Mot-clé surligné (marqueur orange sous le texte)
\newcommand{\cle}[1]{{\popxbold\color{popdark}#1}}

% ============================================================
% En-tête / pied
% ============================================================
\pagestyle{fancy}
\fancyhf{}
\renewcommand{\headrulewidth}{0pt}
\renewcommand{\footrulewidth}{0pt}
\lhead{\raisebox{-0.15ex}{\poplogo\fontsize{13}{13}\selectfont\color{popink}OnBuch\color{poporange}+}}
\rhead{\poppill{\DOCMATIERE\ \textbullet\ \DOCNIVEAU\ \textbullet\ Leçon \DOCLECON}{white}{popink}}
\lfoot{\popxbold\fontsize{7.6}{7.6}\selectfont\color{popmuted}\MakeUppercase{Cours signé OnBuch+}\ \textbullet\ {\popsemi\DOCTITRE}}
\rfoot{\tikz[baseline=-0.6ex]\node[rounded corners=8.5pt,fill=popink,minimum height=17pt,minimum width=17pt,inner xsep=5pt,inner sep=0]{\popxbold\fontsize{8}{8}\selectfont\color{popcream}\thepage\,/\,\pageref*{LastPage}};}

% ============================================================
% Titres
% ============================================================
\newcommand{\chaptitre}[2]{%
  \begin{tcolorbox}[enhanced,colback=poporange,colframe=popink,boxrule=2.6pt,arc=11pt,
      drop shadow=popink,left=15pt,right=15pt,top=12pt,bottom=11pt,before skip=3pt,after skip=18pt]
    \poppill{#2}{popink}{popcream}\par\vspace{9pt}
    {\raggedright\hyphenpenalty=10000\exhyphenpenalty=10000\popdisplay\fontsize{22}{24}\selectfont\color{popcream}\MakeUppercase{#1}\par}
  \end{tcolorbox}
}
% Section numérotée : pastille orange avec le numéro + titre Archivo
\newcounter{popsec}
\newcommand{\coursec}[1]{%
  \stepcounter{popsec}\setcounter{popsub}{0}%
  \par\vspace{16pt}\noindent
  \begin{minipage}{\linewidth}
  \tikz[baseline=-0.7ex]\node[draw=popink,line width=1.6pt,fill=poporange,rounded corners=4pt,minimum size=22pt,inner sep=0]{\popdisplay\fontsize{12}{12}\selectfont\color{popcream}\thepopsec};%
  \hspace{8pt}{\popdisplay\fontsize{15}{17}\selectfont\color{popink}\MakeUppercase{#1}}\par
  \vspace{4pt}\noindent{\color{poporange}\rule{36pt}{3.6pt}}
  \end{minipage}\par\nopagebreak\vspace{8pt}
}
\newcounter{popsub}
\newcommand{\courssub}[1]{%
  \stepcounter{popsub}%
  \par\vspace{9pt}\noindent{\popxbold\fontsize{11.5}{13}\selectfont\color{popink}{\color{poporange}\thepopsec.\thepopsub}\hspace{6pt}#1}\par\nopagebreak\vspace{4pt}
}

% ============================================================
% Boîtes pédagogiques — même recette : fond blanc, bordure épaisse colorée,
% ombre dure encre, badge plein. Titre optionnel [#1].
% ============================================================
\tcbset{popbase/.style={breakable,enhanced,colback=white,boxrule=2.3pt,arc=9pt,
  shadow={3pt}{-3pt}{0pt}{popink},left=14pt,right=14pt,top=11pt,bottom=11pt,before skip=11pt,after skip=15pt,
  fontupper=\popsemi\fontsize{10.6}{14.5}\selectfont\color{popink},
  fontlower=\popsemi\fontsize{10.6}{14.5}\selectfont\color{popink}}}
% Coche dessinée (le glyphe ✓ n'existe pas dans les polices du document)
\newcommand{\popcheck}[1]{\tikz[baseline=-0.5ex]\draw[#1,line width=1.6pt,line cap=round,line join=round](-0.13,0.0)--(-0.04,-0.09)--(0.13,0.1);}
\providecommand{\coche}{\popcheck{popgreen}}
\providecommand{\resultat}[1]{\par\smallskip\noindent\fcolorbox{popgreen}{popgreenL}{\parbox{\dimexpr\linewidth-2\fboxsep-2\fboxrule\relax}{\textbf{Résultat.} #1}}\par\smallskip}
\newcommand{\popboxhead}[3]{\poppill{#1}{#2}{white}\ifx\relax#3\relax\else\hspace{6pt}{\popxbold\fontsize{10.6}{12}\selectfont\color{popink}#3}\fi\par\vspace{7pt}}

\newtcolorbox{aretenir}[1][]{popbase,colframe=popgreen,before upper={\popboxhead{À retenir}{popgreen}{#1}}}
% Texte en anglais (document de lecture, dialogue, extrait) : cadre
% d'étude. Un titre optionnel (source, auteur) s'affiche dans l'en-tête.
\newtcolorbox{textebox}[1][]{popbase,colframe=popblue,before upper={\popboxhead{Text}{popblue}{#1}\begin{otherlanguage}{english}},after upper={\end{otherlanguage}}}
% Marques ✓ / ✗ (forme correcte / fautive) souvent utilisées par les modèles
\providecommand{\cross}{\textcolor{poppink}{\ensuremath{\times}}}
\providecommand{\xmark}{\cross}\providecommand{\crossmark}{\cross}\providecommand{\cmark}{\ensuremath{\checkmark}}
% Ligne de réponse / justification à compléter (inventée par les modèles)
\providecommand{\justification}[1]{\par\noindent\textit{Justification~:} \dotfill\par}
% \textipa (package tipa non chargé) : la police ne couvre pas l'API -> simple passage
\providecommand{\textipa}[1]{#1}
% Soulignement (ulem non chargé) : \uline, \ul
\providecommand{\uline}[1]{\underline{#1}}\providecommand{\ul}[1]{\underline{#1}}
% \glossaire{\item ...} inventée par les modèles : liste de glossaire
\providecommand{\glossaire}[1]{\begin{itemize}#1\end{itemize}}
% \alert (beamer) inventée par les modèles : texte mis en évidence
\providecommand{\alert}[1]{\textbf{\textcolor{poppink}{#1}}}
% variantes ASCII d'ordinaux écrites en macro
\providecommand{\iere}{\textsuperscript{re}}\providecommand{\ieme}{\textsuperscript{e}}\providecommand{\ere}{\textsuperscript{re}}\providecommand{\eme}{\textsuperscript{e}}\providecommand{\er}{\textsuperscript{er}}\providecommand{\ie}{\textsuperscript{e}}
% Ordinaux français écrits en macro par les modèles : 1\ière, 3\ième
\newcommand{\ière}{\textsuperscript{re}}\newcommand{\ième}{\textsuperscript{e}}
% \exemple (inventée par les modèles) : étiquette « Exemple : »
\providecommand{\exemple}{\textbf{Exemple~:}\ }
% \square utilisable aussi hors mode math (cases à cocher)
\AtBeginDocument{\let\squareMath\square\renewcommand{\square}{\ensuremath{\squareMath}}}
% Mot ou phrase en anglais à l'intérieur d'un paragraphe français
\newcommand{\eng}[1]{\ifmmode\text{\textit{#1}}\else\begingroup\selectlanguage{english}\itshape #1\endgroup\fi}
\let\esp\eng % alias toléré
% flèches d'intonation inventées par les modèles
\providecommand{\textsearrow}{}\renewcommand{\textsearrow}{\ensuremath{\searrow}}\providecommand{\textnearrow}{}\renewcommand{\textnearrow}{\ensuremath{\nearrow}}
% \fr{..} : mot français (inventé par les modèles)
\providecommand{\fr}[1]{\textit{#1}}
\newtcolorbox{exemplebox}[1][]{popbase,colframe=poporange,before upper={\popboxhead{Exemple}{poporange}{#1}}}
\newtcolorbox{definition}[1][]{popbase,colframe=popblue,before upper={\popboxhead{Définition}{popblue}{#1}}}
\newtcolorbox{propriete}[1][]{popbase,colframe=poppurple,before upper={\popboxhead{Propriété}{poppurple}{#1}}}
\newtcolorbox{methode}[1][]{popbase,colframe=popgold,before upper={\popboxhead{Méthode}{popgold}{#1}}}
\newtcolorbox{attention}[1][]{popbase,colframe=poppink,before upper={\popboxhead{Attention}{poppink}{#1}}}
\newtcolorbox{experience}[1][]{popbase,colframe=popblue,colback=popblueL,before upper={\popboxhead{Expérience}{popblue}{#1}}}
% « Le savais-tu ? » : carte violette pleine (contexte, culture, Cameroun)
\newtcolorbox{savaistu}[1][]{popbase,colback=poppurple,colframe=popink,
  fontupper=\popsemi\fontsize{10.4}{14.2}\selectfont\color{white},
  before upper={\poppill{Le savais-tu\,?}{white}{poppurple}\ifx\relax#1\relax\else\hspace{6pt}{\popxbold\color{white}#1}\fi\par\vspace{7pt}}}
% Exercice résolu : énoncé en haut, \tcblower puis la solution sur fond crème
\newcounter{popexo}\newcounter{popexr}
\newtcolorbox{exoresolu}[1][]{popbase,colframe=poporange,colbacklower=popcream,
  segmentation style={popink,line width=1.2pt,dashed},
  before upper={\stepcounter{popexr}\popboxhead{Exercice résolu \thepopexr}{poporange}{#1}},
  before lower={\poppill{Solution}{popink}{popcream}\par\vspace{6pt}}}
% Exercice (non résolu en place) — la correction est dans la partie Corrigés
\newtcolorbox{exercice}[1][]{popbase,colframe=popink,
  before upper={\stepcounter{popexo}\popboxhead{Exercice \thepopexo}{popink}{#1}}}
% Corrigé
\newtcolorbox{corrige}[1][]{popbase,colframe=popgreen,colback=popgreenL,
  before upper={\popboxhead{Corrigé}{popgreen}{#1}}}
% Objectifs / prérequis / bilan
\newtcolorbox{objectifs}{popbase,colframe=popink,colback=poporangeL,
  before upper={\popboxhead{Objectifs de la leçon}{popink}{}}}
\newtcolorbox{prerequis}{popbase,colframe=popink,colback=popblueL,
  before upper={\popboxhead{Prérequis}{popblue}{}}}
\newtcolorbox{bilan}{popbase,colframe=popink,colback=white,boxrule=2.8pt,
  before upper={\popboxhead{Fiche bilan}{poporange}{L'essentiel à connaître pour l'examen}}}
% Figure encadrée avec légende
\newtcolorbox{popfigure}[1][]{enhanced,breakable=false,colback=white,colframe=popline,boxrule=1.4pt,arc=8pt,
  left=10pt,right=10pt,top=10pt,bottom=8pt,before skip=10pt,after skip=12pt,halign=center,
  lower separated=false,halign lower=center,
  fontlower=\popsemi\fontsize{9}{11.5}\selectfont\color{popmuted},
  #1}
\newcommand{\legende}[1]{\par\vspace{6pt}{\popsemi\fontsize{9}{11.5}\selectfont\color{popmuted}#1}}

% ============================================================
% Signature OnBuch+
% ============================================================
\newcommand{\poplogoinline}[1]{{\poplogo\fontsize{#1}{#1}\selectfont\color{popink}OnBuch\color{poporange}+}}
\newcommand{\signatureonbuch}{%
  \begin{tcolorbox}[enhanced,colback=popink,colframe=popink,boxrule=2.2pt,arc=10pt,
      drop shadow=poporange,left=14pt,right=14pt,top=10pt,bottom=10pt,before skip=0pt,after skip=0pt]
    \tikz[baseline=-0.7ex]\node[circle,fill=popgreen,draw=popcream,line width=1.3pt,minimum size=22pt,inner sep=0]{\popcheck{white}};%
    \hspace{9pt}\begin{minipage}[c]{0.62\linewidth}
      {\popxbold\fontsize{12}{13}\selectfont\color{popcream}Signé\ \textbullet\ {\poplogo OnBuch\color{poporange}+}}\par\vspace{2pt}
      {\raggedright\popsemi\fontsize{8}{10.5}\selectfont\color{popline}\MakeUppercase{Équipe pédagogique OnBuch+}\\\MakeUppercase{Conforme au programme officiel MINESEC}\par}
    \end{minipage}\hfill
    {\poplogo\fontsize{22}{22}\selectfont\color{popcream}OnBuch\color{poporange}+}
  \end{tcolorbox}%
}

% ============================================================
% Page de garde
% #1 titre · #2 matière · #3 niveau · #4 module · #5 n° leçon · #6 durée ·
% #7 description / objectifs (texte ou itemize)
% ============================================================
\newcommand{\pagegarde}[7]{%
  \thispagestyle{empty}
  \newgeometry{a4paper,margin=1.7cm,top=1.6cm,bottom=1.4cm}%
  \begin{tikzpicture}[remember picture,overlay]
    \fill[poporange!18] ([xshift=-3.2cm,yshift=-3.6cm]current page.north east) circle (4.6cm);
    \fill[poppurple!14] ([xshift=2.4cm,yshift=7.6cm]current page.south west) circle (3.8cm);
  \end{tikzpicture}%
  {\poplogo\fontsize{30}{30}\selectfont\color{popink}OnBuch\color{poporange}+}\hfill
  \raisebox{0.6ex}{\poppill{Cours signé \textbullet\ Premium}{popink}{popcream}}\par
  \vspace{1pt}\popsquiggle{poppurple}\par
  \vspace{8pt}
  \begin{tcolorbox}[enhanced,colback=poporange,colframe=popink,boxrule=2.8pt,arc=13pt,
      drop shadow=popink,left=18pt,right=18pt,top=12pt,bottom=13pt,after skip=10pt]
    \poppill{#2 \textbullet\ #3}{popink}{popcream}\hspace{5pt}\poppill{#4}{white}{popink}\par\vspace{8pt}
    {\popdisplay\fontsize{14}{14}\selectfont\color{popink}LEÇON #5}\par\vspace{4pt}
    {\raggedright\hyphenpenalty=10000\exhyphenpenalty=10000\popdisplay\fontsize{25}{27}\selectfont\color{popcream}\MakeUppercase{#1}\par}
    \vspace{8pt}
    {\popxbold\fontsize{9.5}{9.5}\selectfont\color{popink}\MakeUppercase{Durée conseillée : #6}}
  \end{tcolorbox}
  \begin{tcolorbox}[enhanced,colback=white,colframe=popink,boxrule=2.2pt,arc=10pt,
      drop shadow=popink,left=16pt,right=16pt,top=10pt,bottom=10pt,after skip=8pt]
    \poppill{Dans cette leçon}{poppurple}{white}\par\vspace{5pt}
    \popsemi\fontsize{9.3}{12.6}\selectfont\color{popink}#7
  \end{tcolorbox}
  \noindent\begin{minipage}[t]{0.487\linewidth}
    \begin{tcolorbox}[enhanced,colback=popgreen,colframe=popink,boxrule=2.2pt,arc=10pt,
        drop shadow=popink,left=11pt,right=11pt,top=10pt,bottom=10pt,height=3.1cm,valign=center]
      \tcbox[colback=white,colframe=popink,boxrule=1.3pt,arc=5pt,boxsep=0pt,nobeforeafter,
        left=3pt,right=3pt,top=3pt,bottom=3pt,valign=center]{\qrcode[height=1.9cm]{https://play.google.com/store/apps/details?id=com.luvvix.onbuch&hl=fr}}%
      \hspace{8pt}\begin{minipage}[c]{0.5\linewidth}
        {\popdisplay\fontsize{11}{12.5}\selectfont\color{white}\MakeUppercase{Télécharge OnBuch}}\par\vspace{4pt}
        {\raggedright\popsemi\fontsize{8.4}{11}\selectfont\color{white}Gratuit \textbullet\ Google Play\\Tuteur IA, annales, concours, résultats\par}
      \end{minipage}
    \end{tcolorbox}
  \end{minipage}\hfill
  \begin{minipage}[t]{0.487\linewidth}
    \begin{tcolorbox}[enhanced,colback=poppurple,colframe=popink,boxrule=2.2pt,arc=10pt,
        drop shadow=popink,left=11pt,right=11pt,top=10pt,bottom=10pt,height=3.1cm,valign=center]
      \tikz[baseline=-0.7ex]\node[circle,fill=white,draw=popink,line width=1.3pt,minimum size=1.6cm,inner sep=0]{\popdisplay\fontsize{24}{24}\selectfont\color{poppurple}+};%
      \hspace{8pt}\begin{minipage}[c]{0.6\linewidth}
        {\popdisplay\fontsize{11}{12.5}\selectfont\color{white}\MakeUppercase{Passe à OnBuch+}}\par\vspace{4pt}
        {\raggedright\popsemi\fontsize{8.4}{11}\selectfont\color{white}Tous les cours signés, Tuteur IA illimité, fiches PDF — onglet OnBuch+ de l'app\par}
      \end{minipage}
    \end{tcolorbox}
  \end{minipage}
  \vspace{8pt}
  \popcontenu
  \vfill
  \signatureonbuch
  \restoregeometry
  \newpage
}

% Rangée « Ce cours contient » (page de garde)
\newcommand{\poptile}[3]{% #1 couleur #2 gros texte #3 légende
  \begin{tcolorbox}[enhanced,colback=#1,colframe=popink,boxrule=2pt,arc=9pt,shadow={2.5pt}{-2.5pt}{0pt}{popink},
      left=6pt,right=6pt,top=9pt,bottom=9pt,halign=center,valign=center,height=1.75cm,width=\linewidth,nobeforeafter]
    {\popdisplay\fontsize{12.5}{13}\selectfont\color{white}\MakeUppercase{#2}}\par\vspace{3pt}
    {\centering\popsemi\fontsize{7.8}{9.5}\selectfont\color{white}#3\par}
  \end{tcolorbox}}
\newcommand{\popcontenu}{%
  \par\noindent{\popxboldls\fontsize{7.5}{7.5}\selectfont\color{popmuted}CE COURS SIGNÉ CONTIENT}\par\vspace{7pt}
  \noindent\begin{minipage}[t]{0.235\linewidth}\poptile{popblue}{Cours}{complet, illustré, pas à pas}\end{minipage}\hfill
  \begin{minipage}[t]{0.235\linewidth}\poptile{poporange}{Méthodes}{et exercices résolus}\end{minipage}\hfill
  \begin{minipage}[t]{0.235\linewidth}\poptile{poppink}{Exercices}{type examen + corrigés}\end{minipage}\hfill
  \begin{minipage}[t]{0.235\linewidth}\poptile{popgreen}{Fiche bilan}{l'essentiel à retenir}\end{minipage}\par}

% Sommaire
\newtcolorbox{sommaire}{enhanced,colback=white,colframe=popink,boxrule=2.2pt,arc=10pt,
  drop shadow=popink,left=18pt,right=18pt,top=13pt,bottom=13pt,before skip=6pt,after skip=20pt,
  before upper={\poppill{Sommaire}{poppurple}{white}\par\vspace{9pt}\popsemi\fontsize{10.6}{14.5}\selectfont\color{popink}}}

% Signature de fin de cours — poussée tout en bas de la dernière page
\newcommand{\findecours}{%
  \par\vfill\nopagebreak
  \begin{center}\popsquiggle{poporange}\end{center}\vspace{4pt}
  \signatureonbuch
}
\AtBeginDocument{\def\llbracket{\mathopen{[\mkern-3.5mu[}}\def\rrbracket{\mathclose{]\mkern-3.5mu]}}} % intervalles d'entiers (glyphe ⟦ absent de la police)
% Glyphes absents de Fira Math : redéfinis à partir de glyphes disponibles
\AtBeginDocument{%
  \def\setminus{\mathbin{\backslash}}\let\smallsetminus\setminus
  \def\vdots{\mathord{\vbox{\baselineskip3.2pt\lineskiplimit0pt\kern4pt\hbox{.}\hbox{.}\hbox{.}}}}%
  \def\ddots{\mathinner{\mkern1mu\raise6.5pt\hbox{.}\mkern2mu\raise3.6pt\hbox{.}\mkern2mu\raise0.7pt\hbox{.}\mkern1mu}}%
  \def\triangle{\mathord{\scalebox{0.9}{$\symup{\Delta}$}}}%
  \def\nabla{\mathord{\raisebox{\depth}{\scalebox{1}[-1]{$\symup{\Delta}$}}}}%
  \def\checkmark{\popcheck{popdark}}%
  \def\star{\mathbin{\ast}}\def\bigstar{\mathord{\ast}}%
  \def\diamond{\mathbin{\scalebox{0.6}{$\Diamond$}}}%
  \def\bigcup{\mathop{\mathchoice{\raisebox{-0.3em}{\scalebox{1.7}{$\cup$}}}{\scalebox{1.25}{$\cup$}}{\cup}{\cup}}\displaylimits}%
  \def\bigcap{\mathop{\mathchoice{\raisebox{-0.3em}{\scalebox{1.7}{$\cap$}}}{\scalebox{1.25}{$\cap$}}{\cap}{\cap}}\displaylimits}%
  \def\lesseqqgtr{\mathrel{\lessgtr}}\def\bigoplus{\mathop{\oplus}\displaylimits}%
}
\providecommand{\ssi}{si et seulement si} % abréviation fréquente
\AtBeginDocument{\providecommand{\wideparen}{\overparen}} % arc de cercle

%% FIGFIX-BEGIN v3 — correctifs globaux des figures (docs/onbuch-generation/tools/figfix)
\usepackage{adjustbox}\usepackage{etoolbox}
% toute figure TikZ/pgfplots plus large que la ligne est réduite à la largeur disponible
\BeforeBeginEnvironment{tikzpicture}{\begin{adjustbox}{max width=\linewidth}}
\AfterEndEnvironment{tikzpicture}{\end{adjustbox}}
% légendes pgfplots lisibles par-dessus les courbes
\pgfplotsset{every axis/.append style={legend style={fill=white,fill opacity=0.92,text opacity=1,draw=black!15,inner sep=3pt}}}
% étiquettes d'arêtes (syntaxe edge["..."]) avec fond blanc
\tikzset{every edge quotes/.append style={fill=white,inner sep=1.2pt}}
% bulles de cartes mentales : texte à la ligne dans la bulle, bulle un peu plus grande
\tikzset{every concept/.append style={text width=2.0cm,align=center,minimum size=2.3cm,inner sep=2pt}}
%% FIGFIX-END
\begin{document}

\pagegarde{Lesson 15 — Vocabulary: Words and expressions related to buying and selling articles/foodstuff at a nearby store/market}{Anglais}{1ère A}{Chapter 2 : Economic Life and Occupations}{ENA16}{2 h}{Cette leçon de vocabulaire présente les mots, collocations et expressions nécessaires pour acheter et vendre des articles et des denrées alimentaires dans un magasin ou un marché. Elle insiste sur les faux amis, les familles de mots et l'emploi du lexique dans des dialogues authentiques.
\vspace{4pt}
\begin{multicols}{2}\small
\begin{itemize}
\item À la fin de cette leçon, je sais nommer les principaux articles et denrées vendus dans un magasin ou un marché (foodstuff, household items).
\item À la fin de cette leçon, je sais utiliser le vocabulaire de l'achat et de la vente (customer, seller, price, discount, receipt...).
\item À la fin de cette leçon, je sais employer les collocations courantes du commerce (to bargain for, to run out of, value for money...).
\item À la fin de cette leçon, je sais utiliser les expressions fonctionnelles pour demander un prix, négocier et payer.
\item À la fin de cette leçon, je sais éviter les faux amis et les calques du français (magasin ≠ magazine, monnaie ≠ money...).
\item À la fin de cette leçon, je sais former des mots de la même famille (buy/buyer, sell/sale/seller, cheap/cheaply...).
\end{itemize}\end{multicols}}

\begin{sommaire}
\begin{multicols}{2}
\begin{enumerate}
\item Places and people: where and with whom we buy and sell
\item Articles and foodstuff: what we buy
\item Prices, payment and bargaining
\item Collocations, idioms and functional expressions
\item False friends and traps for francophones
\item Word families, pronunciation and spelling
\item Using the vocabulary in context
\item Activité d'intégration
\item Méthodes et exercices résolus
\item Exercices
\item Corrigés des exercices
\item Fiche bilan
\end{enumerate}
\end{multicols}
\end{sommaire}

\begin{objectifs}
\begin{itemize}
\item À la fin de cette leçon, je sais nommer les principaux articles et denrées vendus dans un magasin ou un marché (foodstuff, household items).
\item À la fin de cette leçon, je sais utiliser le vocabulaire de l'achat et de la vente (customer, seller, price, discount, receipt...).
\item À la fin de cette leçon, je sais employer les collocations courantes du commerce (to bargain for, to run out of, value for money...).
\item À la fin de cette leçon, je sais utiliser les expressions fonctionnelles pour demander un prix, négocier et payer.
\item À la fin de cette leçon, je sais éviter les faux amis et les calques du français (magasin ≠ magazine, monnaie ≠ money...).
\item À la fin de cette leçon, je sais former des mots de la même famille (buy/buyer, sell/sale/seller, cheap/cheaply...).
\item À la fin de cette leçon, je sais prononcer et orthographier correctement les mots clés du thème.
\item À la fin de cette leçon, je sais réemployer ce lexique dans des phrases, des dialogues et un court texte.
\end{itemize}
\end{objectifs}

\begin{prerequis}
\begin{itemize}
\item Les noms de nourriture courants vus en Seconde (food and drinks).
\item Les nombres cardinaux et l'expression des prix (How much is/are...?).
\item Le présent simple et les questions en wh- et en do/does.
\item Les quantifieurs de base : some, any, much, many, a few, a little.
\item Les formules de politesse : Can I...? Could you...? I'd like...
\end{itemize}
\end{prerequis}

\newpage

\coursec{Places and people: where and with whom we buy and sell}

Every morning in Bamenda, Douala or Buea, thousands of Cameroonians walk to the market or the nearby shop to buy foodstuffs: plantains, tomatoes, rice, soap. But how do you ask for a discount in English? How do you complain about a bad product? Imagine you are at Mokolo Market in Yaoundé and a tourist asks you: \eng{“How much is a bunch of plantains?”} « Combien coûte un régime de plantains ? » — will you answer correctly? In this lesson, you will learn the exact words and expressions used to buy and sell articles and foodstuffs, so that you can shop confidently in any English-speaking environment, from a London supermarket to a Nigerian roadside stall.

\courssub{Where we buy: markets, shops and stalls}

\begin{textebox}{At Mokolo Market}
At Mokolo Market in Yaoundé, vendors sell fresh foodstuffs from small stalls, while shopkeepers in town sell packaged goods in their shops. On market days, the place is crowded with customers bargaining for tomatoes, plantains and dried fish. A wholesaler delivers bags of rice to retailers early in the morning. Hawkers walk along the roads selling groundnuts and water in plastic sachets. Supermarkets like Mahima or Casino offer fixed prices and trolleys, but many people prefer the open-air market for fresh produce and the possibility to negotiate.
\end{textebox}

\begin{definition}[Vocabulaire clé : Les lieux d'achat]
Voici les principaux lieux où l'on achète et vend, avec leurs nuances :
\begin{itemize}
    \item \cle{market} /ˈmɑːkɪt/ (nom) : marché (général), souvent en plein air. \eng{The market is busy on Saturdays.} « Le marché est animé les samedis. »
    \item \cle{open-air market} : marché en plein air (synonyme fréquent de \eng{market} en Afrique).
    \item \cle{stall} /stɔːl/ (nom) : étal, stand (emplacement d'un vendeur dans un marché). \eng{She has a stall selling spices.} « Elle tient un stand d'épices. »
    \item \cle{shop} /ʃɒp/ (nom, \textbf{GB}) : magasin, boutique. \eng{The corner shop sells bread.} « La supérette du coin vend du pain. »
    \item \cle{store} /stɔːr/ (nom, \textbf{US}) : magasin, grand magasin. \eng{A department store.} « Un grand magasin. » \textbf{Distinction :} Au Royaume-Uni, \eng{shop} est le terme courant ; aux États-Unis, on dit \eng{store}. Au Cameroun (anglais britannique), on utilise \eng{shop} pour une boutique et \eng{store} pour un grand magasin ou un entrepôt.
    \item \cle{supermarket} /ˈsuːpəmɑːkɪt/ : supermarché (libre-service, prix fixes).
    \item \cle{kiosk} /ˈkiːɒsk/ : kiosque (petite construction vendant journaux, boissons, recharge téléphonique).
    \item \cle{roadside vendor} /ˈrəʊdsaɪd ˈvɛndə(r)/ : vendeur au bord de la route.
    \item \cle{hawker} /ˈhɔːkə(r)/ : marchand ambulant, colporteur (qui crie sa marchandise en marchant).
\end{itemize}
\end{definition}

\begin{propriete}[Shop vs Store : le piège GB/US]
\begin{center}
\begin{tabularx}{\linewidth}{|X|X|X|}
\hline
\rowcolor{popblueL} Contexte & Terme principal & Exemple \\ \hline
Royaume-Uni (GB), Cameroun, Nigeria, Ghana & \textbf{shop} & \eng{I'm going to the shop.} (Je vais au magasin/boutique.) \\ \hline
États-Unis (US) & \textbf{store} & \eng{I'm going to the store.} (Je vais au magasin.) \\ \hline
Grand magasin (GB/US) & \textbf{department store} & \eng{Harrods is a famous department store.} \\ \hline
Entrepôt / Gros & \textbf{store} / \textbf{warehouse} & \eng{The goods are in the store.} (Les marchandises sont dans l'entrepôt/le stock.) \\ \hline
\end{tabularx}
\end{center}
\eng{Shop} peut aussi être verbe : \eng{to shop} (faire ses courses). \eng{Store} est presque toujours nom.
\end{propriete}

\begin{popfigure}
\begin{tikzpicture}[
    mindmap,
    concept color=popblue,
    level 1/.style={level distance=3.5cm, sibling angle=120},
    level 2/.style={level distance=2.5cm},
    every node/.style={font=\small, text=white},
    grow cyclic
]
\node[concept] {BUYING AND SELLING}
    child[concept color=popgreen] { node[concept] {Places}
        child { node[concept] {market / open-air market} }
        child { node[concept] {stall} }
        child { node[concept] {shop (GB) / store (US)} }
        child { node[concept] {supermarket} }
        child { node[concept] {kiosk} }
        child { node[concept] {roadside / hawker} }
    }
    child[concept color=poporange] { node[concept] {People}
        child { node[concept] {customer / consumer} }
        child { node[concept] {seller / vendor} }
        child { node[concept] {shopkeeper / shop assistant} }
        child { node[concept] {trader} }
        child { node[concept] {wholesaler} }
        child { node[concept] {retailer} }
    }
    child[concept color=poppink] { node[concept] {Actions}
        child { node[concept] {to shop / go shopping} }
        child { node[concept] {to trade} }
        child { node[concept] {to deal in} }
        child { node[concept] {to bargain / haggle} }
        child { node[concept] {to purchase / buy} }
        child { node[concept] {to sell / vend} }
    };
\end{tikzpicture}
\legende{Carte mentale : Lieux, acteurs et verbes de l'achat et de la vente.}
\end{popfigure}

\courssub{Who buys and sells: customers, vendors and the supply chain}

\begin{definition}[Vocabulaire clé : Les personnes]
La chaîne commerciale implique des acteurs précis. Attention aux nuances :
\begin{itemize}
    \item \cle{customer} /ˈkʌstəmə(r)/ : client (celui qui achète dans un magasin/marché). \eng{The customer is always right.} « Le client est roi. »
    \item \cle{consumer} /kənˈsjuːmə(r)/ : consommateur (terme économique, celui qui consomme le produit final). \eng{Consumers want lower prices.} « Les consommateurs veulent des prix bas. »
    \item \cle{seller} /ˈselə(r)/ : vendeur (terme général).
    \item \cle{vendor} /ˈvɛndə(r)/ : vendeur, vendeuse (souvent sur un marché, dans la rue, ou formel : \eng{street vendor}, \eng{food vendor}). \eng{A street vendor sells roasted corn.} « Un vendeur ambulant vend du maïs grillé. »
    \item \cle{shopkeeper} /ˈʃɒpkiːpə(r)/ : commerçant, propriétaire d'une boutique. \eng{The shopkeeper knows all his customers.} « Le commerçant connaît tous ses clients. »
    \item \cle{shop assistant} /ˈʃɒp əsɪstənt/ : vendeur(se) employé(e) dans une boutique. \eng{The shop assistant helped me find the right size.} « Le vendeur m'a aidé à trouver la bonne taille. »
    \item \cle{trader} /ˈtreɪdə(r)/ : commerçant, négociant (terme plus large, souvent pour celui qui achète et revend, ex: \eng{market trader}).
    \item \cle{wholesaler} /ˈhəʊlseɪlə(r)/ : grossiste (achète en gros aux producteurs/usines et revend aux détaillants).
    \item \cle{retailer} /ˈriːteɪlə(r)/ : détaillant (vend au détail au consommateur final).
\end{itemize}
\end{definition}

\begin{popfigure}
\begin{tikzpicture}[
    >=Stealth,
    node distance=1.5cm and 2.5cm,
    box/.style={draw=popdark, thick, rounded corners, fill=popcream, minimum width=2.2cm, minimum height=1cm, align=center, font=\small},
    arrow/.style={->, thick, popblue}
]
\node[box, align=center] (producer){Producer / Farmer\\(Producteur / Paysan)};
\node[box, right=of producer, align=center] (wholesaler){Wholesaler\\(Grossiste)};
\node[box, right=of wholesaler, align=center] (retailer){Retailer / Shopkeeper\\(Détaillant / Commerçant)};
\node[box, right=of retailer, align=center] (consumer){Consumer / Customer\\(Consommateur / Client)};

\draw[arrow] (producer) -- node[above, font=\tiny, text=popdark, align=center]{sells in bulk\\(vend en gros)} (wholesaler);
\draw[arrow] (wholesaler) -- node[above, font=\tiny, text=popdark, align=center]{supplies\\(approvisionne)} (retailer);
\draw[arrow] (retailer) -- node[above, font=\tiny, text=popdark, align=center]{sells retail\\(vend au détail)} (consumer);

% Feedback loop
\draw[arrow, poporange, bend left=30] (consumer) to node[above, font=\tiny, text=poporange] {demand / feedback} (producer);
\end{tikzpicture}
\legende{Chaîne de distribution : Wholesaler $\rightarrow$ Retailer $\rightarrow$ Consumer.}
\end{popfigure}

\begin{exemplebox}{Phrases contextualisées (Cameroun / Anglophone)}
\begin{itemize}
    \item \eng{The wholesaler at Nkwen market supplies plantains to retailers in Bamenda.} « Le grossiste du marché de Nkwen approvisionne les détaillants de Bamenda en plantains. »
    \item \eng{My mother is a regular customer at Madam Clara's shop.} « Ma mère est une cliente fidèle à la boutique de Madame Clara. »
    \item \eng{Street vendors sell puff-puff and beans near the school gate.} « Des vendeurs ambulants vendent des beignets et des haricots près de la grille de l'école. »
    \item \eng{The shop assistant wrapped the sugar in a blue paper bag.} « Le vendeur a emballé le sucre dans un sac en papier bleu. »
    \item \eng{As a consumer, you have the right to return expired goods.} « En tant que consommateur, vous avez le droit de retourner les produits périmés. »
\end{itemize}
\end{exemplebox}

\courssub{Action verbs: to shop, to trade, to deal in}

\begin{propriete}[Verbes associés aux lieux et aux personnes]
\begin{center}
\begin{tabularx}{\linewidth}{|l|X|X|}
\hline
\rowcolor{popblueL} Verbe & Sens / Emploi & Exemples traduits \\ \hline
\textbf{to shop} & Faire ses courses (achats courants). Intransitif. & \eng{We shop at the supermarket on Saturdays.} « Nous faisons nos courses au supermarché le samedi. » \\ \hline
\textbf{to go shopping} & Aller faire les courses (insiste sur le déplacement). & \eng{Let's go shopping for school uniforms.} « Allons acheter des uniformes scolaires. » \\ \hline
\textbf{to do the shopping} & Faire les courses (corvée ménagère, focus sur l'action accomplie). & \eng{Who does the shopping in your family?} « Qui fait les courses dans ta famille ? » \\ \hline
\textbf{to trade} & Commercier, faire du commerce (activité professionnelle) ; échanger. & \eng{They trade in cocoa and coffee.} « Ils commercent le cacao et le café. » \\ \hline
\textbf{to deal in} & Spécialisé dans la vente de (souvent pour un type de marchandises). & \eng{This shop deals in electronics.} « Ce magasin vend de l'électronique. » \\ \hline
\textbf{to bargain} / \textbf{to haggle} & Marchander, négocier le prix (courant au marché). & \eng{Don't bargain in a supermarket.} « Ne marchandez pas au supermarché. » \\ \hline
\end{tabularx}
\end{center}
\end{propriete}

\begin{attention}[Erreurs fréquentes des francophones]
\begin{enumerate}
    \item \textbf{Confondre \eng{shop} (nom) et \eng{to shop} (verbe).}
    \begin{itemize}
        \item \eng{The shop is open.} (Nom = la boutique).
        \item \eng{I shop here every day.} (Verbe = je fais mes courses ici).
    \end{itemize}
    \item \textbf{Utiliser \eng{store} pour une petite boutique au Cameroun.} $\rightarrow$ Préférez \eng{shop} (boutique) ou \eng{kiosk} (kiosque). \eng{Store} fait penser à un grand magasin ou un entrepôt.
    \item \textbf{Confondre \eng{customer} et \eng{consumer}.} \eng{Customer} = relation commerciale directe (je suis client \textbf{de} ce magasin). \eng{Consumer} = rôle économique global (je consomme du riz, de l'électricité).
    \item \textbf{Traduire "vendeur" toujours par \eng{seller}.} Au marché : \eng{vendor} ou \eng{trader}. En boutique : \eng{shopkeeper} (patron) ou \eng{shop assistant} (employé). \eng{Seller} est neutre/juridique (ex: \eng{the seller signed the contract}).
    \item \textbf{Oublier l'article avec \eng{go shopping} / \eng{do the shopping}.} Jamais \eng{go to shopping} ni \eng{do shopping} (sauf sens général abstrait).
\end{enumerate}
\end{attention}

\begin{savaistu}[Le "Provision Store" en Afrique de l'Ouest anglophone]
Au Nigeria, au Ghana, en Sierra Leone et au Liberia, une petite boutique de quartier vendant denrées alimentaires, boissons, savon, allumettes, etc. s'appelle couramment un \eng{provision store} (parfois raccourci en \eng{provisions}). L'expression vient de \eng{provisions} (vivres, provisions de bouche). Au Cameroun anglophone (Nord-Ouest, Sud-Ouest), on entend aussi \eng{provision shop} ou simplement \eng{shop}. Si vous voyagez à Lagos ou Accra, cherchez le panneau \eng{“Provision Store”} pour acheter du pain, du lait Peak ou du Milo !
\end{savaistu}

\courssub{Observation et pratique guidée}

\begin{exoresolu}{Identifier le lieu et l'acteur}
\textbf{Énoncé :} Lisez les phrases suivantes. Pour chacune, identifiez le \textbf{lieu} (market, shop, supermarket, stall, roadside) et l'\textbf{acteur principal} (customer, vendor, shopkeeper, hawker, wholesaler, retailer). Traduisez ensuite la phrase.

1. \eng{The \textbf{hawker} walked along the \textbf{roadside} selling roasted plantains.}
2. \eng{My father is a \textbf{wholesaler}; he supplies bags of rice to \textbf{retailers} in the \textbf{market}.}
3. \eng{The \textbf{shop assistant} in the \textbf{supermarket} helped the \textbf{customer} find the milk.}
4. \eng{Every Monday is \textbf{market day}; \textbf{traders} arrive early to set up their \textbf{stalls}.}

\textbf{Solution détaillée :}
\tcblower
\begin{enumerate}
    \item \textbf{Lieu :} \cle{roadside} (bord de route). \textbf{Acteur :} \cle{hawker} (marchand ambulant).
    \textit{Traduction :} « Le marchand ambulant marchait le long de la route en vendant des plantains grillés. »
    \item \textbf{Lieux :} \cle{market} (marché). \textbf{Acteurs :} \cle{wholesaler} (grossiste), \cle{retailers} (détaillants).
    \textit{Traduction :} « Mon père est grossiste ; il fournit des sacs de riz aux détaillants du marché. »
    \item \textbf{Lieu :} \cle{supermarket} (supermarché). \textbf{Acteurs :} \cle{shop assistant} (vendeur employé), \cle{customer} (client).
    \textit{Traduction :} « Le vendeur du supermarché a aidé le client à trouver le lait. »
    \item \textbf{Lieu :} \cle{market} (marché), \cle{stalls} (étals). \textbf{Acteurs :} \cle{traders} (commerçants). \textit{Note :} \eng{Market day} = jour de marché.
    \textit{Traduction :} « Chaque lundi est jour de marché ; les commerçants arrivent tôt pour installer leurs étals. »
\end{enumerate}
\end{exoresolu}

\begin{exercice}{Mots croisés lexicaux : Places \& People}
Complétez la grille avec les mots du vocabulaire (anglais). Les définitions sont en français.
\begin{center}
\small
\begin{tabular}{|c|c|c|c|c|c|c|c|c|c|c|c|}
\hline
\multicolumn{12}{|c|}{\textbf{Horizontal}} \\ \hline
1 & Vendeur ambulant (criant sa marchandise) & H & A & W & K & E & R &  &  &  &  \\ \hline
3 & Grossiste (achète en gros) & W & H & O & L & E & S & A & L & E & R \\ \hline
5 & Magasin (anglais britannique) & S & H & O & P &  &  &  &  &  &  \\ \hline
7 & Client (terme économique global) & C & O & N & S & U & M & E & R &  &  \\ \hline
\multicolumn{12}{|c|}{\textbf{Vertical}} \\ \hline
2 & Étal dans un marché & S & T & A & L & L &  &  &  &  &  \\ \hline
4 & Détaillant (vend au détail) & R & E & T & A & I & L & E & R &  &  \\ \hline
6 & Marchand ambulant (ne crie pas nécessairement) & V & E & N & D & O & R &  &  &  &  \\ \hline
8 & Propriétaire d'une boutique & S & H & O & P & K & E & E & P & E & R \\ \hline
\end{tabular}
\end{center}
\textit{Indice : Les cases vides indiquent la fin du mot. Écrivez les mots complets dans votre cahier.}
\end{exercice}

\begin{corrige}{Exercice : Mots croisés lexicaux}
\begin{itemize}
    \item 1 Horizontal : \textbf{HAWKER}
    \item 3 Horizontal : \textbf{WHOLESALER}
    \item 5 Horizontal : \textbf{SHOP}
    \item 7 Horizontal : \textbf{CONSUMER}
    \item 2 Vertical : \textbf{STALL}
    \item 4 Vertical : \textbf{RETAILER}
    \item 6 Vertical : \textbf{VENDOR}
    \item 8 Vertical : \textbf{SHOPKEEPER}
\end{itemize}
Vérifiez l'orthographe : double \textbf{L} dans \eng{stall}, \eng{sell} $\rightarrow$ \eng{seller} mais \eng{retailer} (un seul L), \eng{wholesaler} (un seul L). \eng{Shopkeeper} = \eng{shop} + \eng{keeper} (double P, double E).
\end{corrige}

\begin{methode}[Comment décrire une scène de marché (Speaking/Writing)]
Pour décrire un lieu d'achat et ses acteurs en anglais :
\begin{enumerate}
    \item \textbf{Situez :} \eng{At [Mokolo Market] / In a [supermarket] / Along the [roadside]...}
    \item \textbf{Nommez les lieux :} \eng{There are many [stalls] / a big [shop] / [kiosks]...}
    \item \textbf{Identifiez les gens :} \eng{[Vendors] sell... / [Shoppers] walk... / A [hawker] shouts...}
    \item \textbf{Utilisez les verbes d'action :} \eng{People [bargain] / [buy] / [carry] bags.}
    \item \textbf{Ajoutez l'atmosphère :} \eng{It is [crowded] / [noisy] / [colourful].}
\end{enumerate}
\textbf{Modèle :} \eng{“At the open-air market in Buea, traders set up their stalls early. Vendors sell fresh tomatoes and hawkers carry water sachets. Customers bargain loudly. It is a busy, colourful place.”}
\end{methode}

\begin{experience}[Jeu de rôle : "At the Market" — Binôme]
\textbf{Contexte :} Samedi matin, marché de Mfoundi (Yaoundé) ou Marché Central (Douala).
\textbf{Rôles :}
\begin{itemize}
    \item \textbf{Élève A (Customer) :} Vous cherchez des tomates, des oignons et du piment. Vous voulez marchander (\eng{bargain}).
    \item \textbf{Élève B (Vendor / Trader) :} Vous vendez des légumes frais. Vous connaissez vos prix.
\end{itemize}
\textbf{Contraintes linguistiques :} Utilisez au moins 5 mots du vocabulaire de cette section (\eng{stall, price, bargain, customer, vendor, fresh, kilo, bag, change, receipt}).
\textbf{Déroulement :} 3 min de préparation (notes autorisées), 3 min de jeu oral. Changez de rôle. Le professeur note l'usage correct du vocabulaire ciblé.
\end{experience}

\coursec{Articles and foodstuff: what we buy}

Dans la section précédente, nous avons identifié \textbf{où} l'on achète (\eng{market}, \eng{grocery store}, \eng{supermarket}, \eng{roadside stall}) et \textbf{avec qui} (\eng{seller}, \eng{vendor}, \eng{shopkeeper}, \eng{customer}, \eng{trader}). Nous allons maintenant nous concentrer sur \textbf{ce que} l'on achète : les denrées alimentaires (\eng{foodstuff}, \eng{food items}, \eng{groceries}, \eng{provisions}) et les articles ménagers (\eng{household items} / \eng{household goods}). Au Cameroun comme partout dans le monde anglophone, la précision du vocabulaire selon la nature du produit (comptable ou incomptable) et son conditionnement est indispensable pour faire ses courses sans malentendu.

\courssub{Les grandes catégories de denrées alimentaires (Foodstuff categories)}

Observons d'abord une scène de marché typique. Le texte suivant met en contexte le vocabulaire de base.

\begin{textebox}[Market scene at Mfoundi Market, Yaoundé]
Early on Saturday morning, the Mfoundi market is already busy. Mama Ngo Bell arranges her stall. She has a large \textbf{heap of tomatoes} beside a \textbf{bunch of ripe plantains}. In wooden crates, she displays \textbf{yams}, \textbf{cassava tubers} and \textbf{cocoyams}. On a wooden table, there are \textbf{bags of rice} (50\,kg and 25\,kg), \textbf{packets of cornflour} for \emph{fufu} and \textbf{cups of garri}. Near the scale, a \textbf{tin of tomato paste}, a \textbf{bottle of cooking oil} and a \textbf{packet of sugar} are waiting for customers. She also sells \textbf{smoked fish}, \textbf{dried fish} and \textbf{beef} kept in a cool box. A customer asks: “How much is a \textbf{heap of tomatoes}?” Mama Ngo Bell replies: “Five hundred francs.” The customer picks a \textbf{bunch of plantains} and asks for a \textbf{bottle of oil}.
\end{textebox}

\begin{definition}[Foodstuff]
\eng{Foodstuff} (prononcé « foud-steuf ») désigne une \textbf{substance utilisée comme aliment} (denrée alimentaire). C'est un nom \textbf{généralement incomptable} (uncountable) en anglais standard. On dit \eng{much foodstuff} (beaucoup de denrées), pas \eng{many foodstuffs}. Le pluriel \eng{foodstuffs} existe dans des contextes techniques ou administratifs (ex. \eng{foodstuffs regulations}), mais on l'évite dans une rédaction courante ou un dialogue de marché. Pour parler de « provisions » au sens de « vivres », on emploie plutôt \eng{provisions} ou \eng{groceries}.
\end{definition}

\begin{center}
\begin{tabularx}{\linewidth}{|X|X|X|}
\hline
\rowcolor{popblueL} \textbf{Category / Catégorie} & \textbf{Key vocabulary / Vocabulaire clé} & \textbf{Exemples contextuels / Contextual examples} \\
\hline
\textbf{Cereals / Céréales} & rice, maize (corn), cornflour / maize flour, millet, sorghum, wheat flour & \eng{A bag of rice}, \eng{a packet of cornflour}, \eng{some maize} \\
\hline
\textbf{Tubers / Tubercules} & plantains, yams, cassava, potatoes, sweet potatoes, cocoyams & \eng{A bunch of plantains}, \eng{a tuber of yam}, \eng{some cassava} \\
\hline
\textbf{Vegetables / Légumes} & tomatoes, onions, okra, carrots, cabbage, green beans, huckleberry (njamanjama) & \eng{A heap of tomatoes}, \eng{a net of onions}, \eng{some okra} \\
\hline
\textbf{Fruits / Fruits} & mangoes, pineapples, papayas, bananas, oranges, avocados, watermelons & \eng{A pile of mangoes}, \eng{a watermelon}, \eng{some bananas} \\
\hline
\textbf{Meat and Fish / Viande et Poisson} & beef, goat meat, bushmeat, chicken, dried fish, smoked fish, fresh fish & \eng{A piece of beef}, \eng{a smoked fish}, \eng{some dried fish} \\
\hline
\textbf{Drinks / Boissons} & water, palm wine, beer, soft drinks, fruit juice, tea, coffee, Milo & \eng{A bottle of water}, \eng{a can of beer}, \eng{a packet of tea} \\
\hline
\end{tabularx}
\end{center}

\begin{attention}[Orthographe : pluriels en \eng{-o}]
Les noms en \eng{-o} prennent souvent \eng{-es} au pluriel : \eng{tomato} $\rightarrow$ \eng{tomatoes}, \eng{potato} $\rightarrow$ \eng{potatoes}, \eng{mango} $\rightarrow$ \eng{mangoes} (ou \eng{mangos}), \eng{hero} $\rightarrow$ \eng{heroes}. Mais : \eng{avocado} $\rightarrow$ \eng{avocados}, \eng{kilo} $\rightarrow$ \eng{kilos}, \eng{photo} $\rightarrow$ \eng{photos}. Attention aussi à \eng{leaf} $\rightarrow$ \eng{leaves} (\eng{bitterleaf} reste incomptable).
\end{attention}

\courssub{Articles ménagers (Household items / goods)}

Au-delà de l'alimentaire, le « \eng{nearby store} » ou le \eng{provisions store} vend des biens de consommation courante.

\begin{center}
\begin{tabularx}{\linewidth}{|X|X|}
\hline
\rowcolor{popgreenL} \textbf{English} & \textbf{Français} \\
\hline
soap (bar soap, laundry soap, liquid soap) & savon (savon en pain, savon de lessive, savon liquide) \\
matches / a box of matches & allumettes / une boîte d'allumettes \\
candles & bougies \\
cooking oil (groundnut oil, palm oil, vegetable oil) & huile de cuisine (arachide, palme, végétale) \\
salt (table salt, coarse salt) & sel (sel de table, gros sel) \\
sugar (cube sugar, granulated sugar) & sucre (en morceaux, en poudre) \\
detergent / washing powder & détergent / poudre à laver \\
bleach & eau de Javel \\
toothpaste, toothbrush & dentifrice, brosse à dents \\
\hline
\end{tabularx}
\end{center}

\courssub{Quantités et conditionnements : la structure \eng{« a/an + quantifier + of + noun »}}

C'est le point grammatical central du vocabulaire d'achat. En français, on dit souvent « un sac de riz », « une bouteille d'huile » sans article devant le contenu. En anglais, la structure est figée :

\begin{propriete}[Noms de conditionnement et structure \eng{of}]
Pour exprimer une quantité d'un nom \textbf{incomptable} (\eng{rice}, \eng{oil}, \eng{sugar}, \eng{flour}, \eng{garri}, \eng{water}) ou d'un \textbf{pluriel indéfini} (\eng{tomatoes}, \eng{plantains}, \eng{matches}), on utilise un \textbf{nom de conditionnement} (container / quantity word) suivi de la préposition \eng{of}.
\begin{center}
\eng{a/an + quantity word + of + noun}
\end{center}
Le nom de conditionnement s'accorde au singulier ou au pluriel (\eng{a bag}, \eng{two bags}) ; le nom qui suit \eng{of} ne change jamais (\eng{two bags of rice}, jamais \eng{*two bags of rices}).
\end{propriete}

\begin{center}
\begin{tabularx}{\linewidth}{|X|X|X|}
\hline
\rowcolor{poporangeL} \textbf{Quantity word / Conditionnement} & \textbf{Typical collocations / Collocations types} & \textbf{Traduction} \\
\hline
a bag / bags of & rice, maize, flour, cement, charcoal & un sac de / des sacs de \\
a packet / packets of & sugar, salt, tea, coffee, biscuits, cornflour, washing powder & un paquet de / des paquets de \\
a bottle / bottles of & cooking oil, water, palm wine, beer, soft drink, bleach & une bouteille de / des bouteilles de \\
a tin / tins of & tomato paste, milk, sardines, corned beef, paint & une boîte (métallique) de / des boîtes de \\
a bunch / bunches of & plantains, bananas, grapes, carrots, spinach & un régime / une grappe / une botte de \\
a heap / heaps of & tomatoes, onions, potatoes, groundnuts, sand & un tas / une pile de \\
a cup / cups of & garri, flour, sugar, tea (mesure) & une tasse / un bol de \\
a piece / pieces of & meat, fish, soap, cake, furniture, advice & un morceau / un bout de \\
a bar / bars of & soap, chocolate, gold & un pain / une barre de \\
a box / boxes of & matches, chocolates, cereal, tissues & une boîte (carton) de \\
a sachet / sachets of & tomato paste (small), milk powder, spices & un sachet de \\
\hline
\end{tabularx}
\end{center}

\begin{exemplebox}[Exemples traduits]
\eng{She bought two bags of rice.} = Elle a acheté deux sacs de riz. \\
\eng{I need a bottle of palm oil.} = J'ai besoin d'une bouteille d'huile de palme. \\
\eng{Give me a heap of tomatoes, please.} = Donnez-moi un tas de tomates, s'il vous plaît. \\
\eng{He drinks a cup of garri every morning.} = Il boit un bol de gari chaque matin. \\
\eng{We need three tins of tomato paste for the stew.} = Il nous faut trois boîtes de concentré de tomates pour le ragoût.
\end{exemplebox}

\courssub{Noms comptables et incomptables (Countable vs Uncountable) — Piège majeur}

La distinction détermine l'usage de \eng{a/an}, \eng{many/much}, \eng{few/little}, \eng{some/any}. Le français a tendance à tout compter (« des riz », « des huiles ») ; l'anglais standard ne le permet pas.

\begin{popfigure}
\begin{tikzpicture}[
    node distance=0.8cm and 1.2cm,
    box/.style={draw=popblue, thick, rounded corners, fill=popblueL, minimum width=3.2cm, minimum height=1cm, align=center, font=\small},
    arrow/.style={-Stealth, thick, popdark},
    title/.style={font=\bfseries\large, text=popblue, align=center}
]
% Titles
\node[title] (titleC) at (0,4.5) {COUNTABLE (Comptables)};
\node[title] (titleU) at (7,4.5) {UNCOUNTABLE (Incomptables)};

% Countable column
\node[box, align=center] (c1) at (0,3.2){Singulier : \eng{a tomato}\\\eng{an onion}\\\eng{a plantain}};
\node[box, align=center] (c2) at (0,2.0){Pluriel : \eng{tomatoes}\\\eng{onions}\\\eng{plantains}};
\node[box, align=center] (c3) at (0,0.8){Quantificateurs :\\\eng{many}, \eng{a few}, \eng{few}\\\eng{How many...?}};
\node[box, align=center] (c4) at (0,-0.4){Verbe : \eng{is} (sg) / \eng{are} (pl)\\\eng{This tomato is...}\\\eng{These tomatoes are...}};

% Uncountable column
\node[box, align=center] (u1) at (7,3.2){Pas de \eng{a/an}\\\eng{rice}, \eng{oil}, \eng{sugar}\\\eng{garri}, \eng{water}, \eng{flour}};
\node[box, align=center] (u2) at (7,2.0){Pas de pluriel en \eng{-s}\\\eng{*rices}, \eng{*oils}, \eng{*sugars}};
\node[box, align=center] (u3) at (7,0.8){Quantificateurs :\\\eng{much}, \eng{a little}, \eng{little}\\\eng{How much...?}};
\node[box, align=center] (u4) at (7,-0.4){Verbe toujours singulier :\\\eng{is}, \eng{has}\\\eng{This rice is good.}};

% Arrows for packaging (bridge)
\node[box, fill=poporangeL, draw=poporange, text width=5cm] (bridge) at (3.5,-1.8) {PONT : Conditionnement (\eng{of})\\\eng{a bag \textbf{of} rice}\\\eng{a bottle \textbf{of} oil}\\\eng{a heap \textbf{of} tomatoes}};

\draw[arrow] (c1) -- (c2);
\draw[arrow] (c2) -- (c3);
\draw[arrow] (c3) -- (c4);
\draw[arrow] (u1) -- (u2);
\draw[arrow] (u2) -- (u3);
\draw[arrow] (u3) -- (u4);
\draw[arrow, dashed] (c4.south) |- (bridge.west);
\draw[arrow, dashed] (u4.south) |- (bridge.east);
\end{tikzpicture}
\legende{Noms comptables vs incomptables : articles et quantificateurs}
\end{popfigure}

\begin{attention}[Pièges francophones : Comptable / Incomptable]
\begin{itemize}
    \item \eng{Rice}, \eng{garri}, \eng{flour}, \eng{cornflour}, \eng{sugar}, \eng{salt}, \eng{oil}, \eng{water}, \eng{soap}, \eng{advice}, \eng{furniture}, \eng{luggage}, \eng{information}, \eng{news} sont \textbf{incomptables}. \eng{*A rice}, \eng{*two garris}, \eng{*three soaps}, \eng{*many furnitures} sont \textbf{incorrects}.
    \item Pour les compter : \eng{a grain of rice}, \eng{a cup of garri}, \eng{a bag of flour}, \eng{a bar of soap}, \eng{a piece of furniture}, \eng{a piece of advice}.
    \item \textbf{Vegetables / Fruits} : \eng{tomato}, \eng{onion}, \eng{plantain}, \eng{yam}, \eng{potato}, \eng{mango}, \eng{orange} sont \textbf{comptables} (\eng{a tomato}, \eng{three tomatoes}). Mais \eng{okra}, \eng{spinach}, \eng{lettuce} et souvent \eng{cabbage} (traité comme un tout) sont \textbf{incomptables} (\eng{some okra}, \eng{a head of cabbage}).
    \item \textbf{Fish} : \eng{fish} (la chair, l'aliment) est incomptable (\eng{some fish}) ; \eng{a fish} = un poisson entier (l'animal). \eng{Fishes} = espèces de poissons (biologie).
    \item \textbf{Money} est incomptable : \eng{much money}, \eng{some money}, jamais \eng{*a money} ni \eng{*moneys} (sauf droit/finance : \eng{public monies}).
    \item \textbf{Hair} : incomptable (\eng{She has long hair.}). \eng{A hair} = un cheveu unique.
\end{itemize}
\end{attention}

\courssub{Adjectifs de qualité (Quality adjectives)}

Pour négocier ou décrire un produit, ces adjectifs sont quotidiens.

\begin{center}
\begin{tabularx}{\linewidth}{|X|X|X|}
\hline
\rowcolor{poppurpleL} \textbf{Adjective} & \textbf{Meaning / Sens} & \textbf{Exemple / Example} \\
\hline
fresh & frais (récolté/pêché récemment) & \eng{Fresh fish arrives every morning.} \\
ripe & mûr (prêt à manger) & \eng{These plantains are ripe.} \\
unripe / green & pas mûr, vert & \eng{The mangoes are still unripe.} \\
rotten & pourri (impropre à la consommation) & \eng{Throw away the rotten tomatoes.} \\
stale & rassis, pas frais (pain, poisson séché) & \eng{The bread is stale.} \\
cheap & bon marché, pas cher & \eng{Tomatoes are cheap this season.} \\
expensive / dear & cher & \eng{Beef is becoming expensive.} \\
genuine / authentic & authentique, vrai & \eng{Is this genuine honey?} \\
fake / counterfeit & faux, contrefait & \eng{Beware of fake engine oil.} \\
adulterated & frelaté, mélangé (ex. huile, miel) & \eng{Adulterated palm oil is dangerous.} \\
\hline
\end{tabularx}
\end{center}

\begin{exemplebox}[Adjectifs en contexte — Traduits]
\eng{These plantains are not ripe yet.} = Ces plantains ne sont pas encore mûrs. \\
\eng{The fish is stale; don't buy it.} = Le poisson n'est pas frais ; ne l'achetez pas. \\
\eng{This rice is adulterated with stones.} = Ce riz est frelaté (mélangé à des cailloux). \\
\eng{A heap of fresh tomatoes costs 500 francs.} = Un tas de tomates fraîches coûte 500 francs. \\
\eng{The smoked fish is genuine, not fake.} = Le poisson fumé est authentique, pas faux.
\end{exemplebox}

\begin{attention}[Faux amis à éviter]
\begin{itemize}
    \item \eng{Corn} en anglais britannique = céréales en général ; en anglais américain et camerounais = \textbf{maïs}. Pour éviter l'ambiguïté, utilisez \eng{maize}.
    \item \eng{Actually} = en fait, réellement (et non « actuellement » = \eng{currently}, \eng{at present}) : \eng{Tomatoes are actually cheap this season.} = En fait, les tomates sont bon marché cette saison.
    \item \eng{To demand} = exiger (et non « demander » = \eng{to ask for}) : \eng{The customer asked for a bag of rice.} = Le client a demandé un sac de riz.
    \item \eng{A tin} = une boîte de conserve métallique ; \eng{a can} est l'équivalent américain (et pour les boissons : \eng{a can of beer}).
\end{itemize}
\end{attention}

\courssub{Mise en pratique : dialogue d'achat}

\begin{textebox}[Dialogue at the provisions store]
\textbf{Customer:} Good morning. I need some provisions for the week.\\
\textbf{Seller:} Good morning. What exactly are you looking for?\\
\textbf{Customer:} First, \textbf{a bag of rice} (50\,kg) and \textbf{two packets of cornflour}.\\
\textbf{Seller:} The rice is 18,000 francs; cornflour is 1,200 francs a packet. Anything else?\\
\textbf{Customer:} Yes, \textbf{a bottle of groundnut oil}, \textbf{a packet of cube sugar} and \textbf{a tin of tomato paste}.\\
\textbf{Seller:} The oil is 3,500, the sugar 800, the tomato paste 600. Do you need \textbf{soap} or \textbf{matches}?\\
\textbf{Customer:} Two \textbf{bars of laundry soap} and a \textbf{box of matches}.\\
\textbf{Seller:} That comes to 25,900 francs. Cash or Mobile Money?\\
\textbf{Customer:} Mobile Money. Here is my phone.\\
\textbf{Seller:} Transaction confirmed. Thank you, come again!
\end{textebox}

\begin{savaistu}[Le mot « garri » dans l'anglais ouest-africain]
\eng{Garri} (ou \eng{gari}) désigne la farine de manioc fermentée et grillée, aliment de base au Nigeria, au Ghana, dans les régions anglophones du Cameroun, en Sierra Leone et au Liberia. Le mot figure dans les dictionnaires d'anglais ouest-africain et dans les grands dictionnaires de référence. Vous pouvez écrire : \eng{“I drink garri with groundnuts and sugar.”} (Je bois du gari avec des arachides et du sucre). C'est un nom \textbf{incomptable} : \eng{some garri}, \eng{a cup of garri}, jamais \eng{*a garri} ni \eng{*garris}.
\end{savaistu}

\begin{exoresolu}[Compléter avec le conditionnement correct + traduction]
\textbf{Énoncé :} Complétez chaque phrase avec le mot de conditionnement approprié (\eng{bag, bottle, tin, bunch, heap, cup, packet, bar, box, piece}). Traduisez ensuite.
\begin{enumerate}
    \item Mama sells a \_\_\_\_\_\_ of ripe plantains for 1,000 francs.
    \item Please buy a \_\_\_\_\_\_ of cooking oil (5 litres).
    \item We need three \_\_\_\_\_\_ of tomato paste for the party.
    \item Add two \_\_\_\_\_\_ of garri to the water.
    \item The customer chose a \_\_\_\_\_\_ of laundry soap.
    \item There is a large \_\_\_\_\_\_ of tomatoes near the scale.
    \item Give me a \_\_\_\_\_\_ of sugar, please.
    \item He lit the candle with a match from a \_\_\_\_\_\_ of matches.
    \item She cut a \_\_\_\_\_\_ of beef for the stew.
    \item They carried two \_\_\_\_\_\_ of rice on the motorbike.
\end{enumerate}

\tcblower
\textbf{Solution détaillée :}
\begin{enumerate}
    \item \textbf{bunch} — \eng{Mama sells a bunch of ripe plantains for 1,000 francs.} = Maman vend un régime de plantains mûrs pour 1 000 francs. (\eng{plantains} = comptable pluriel ; conditionnement \eng{bunch}.)
    \item \textbf{bottle} — \eng{Please buy a bottle of cooking oil (5 litres).} = Achète une bouteille d'huile de cuisine (5 litres), s'il te plaît. (\eng{oil} = incomptable, liquide $\rightarrow$ \eng{bottle}.)
    \item \textbf{tins} — \eng{We need three tins of tomato paste for the party.} = Nous avons besoin de trois boîtes de concentré de tomates pour la fête. (\eng{tomato paste} = incomptable, boîte métallique $\rightarrow$ \eng{tin}.)
    \item \textbf{cups} — \eng{Add two cups of garri to the water.} = Ajoute deux bols de gari dans l'eau. (\eng{garri} = incomptable, mesure $\rightarrow$ \eng{cup}.)
    \item \textbf{bar} — \eng{The customer chose a bar of laundry soap.} = Le client a choisi un pain de savon de lessive. (\eng{soap} = incomptable, forme solide rectangulaire $\rightarrow$ \eng{bar}.)
    \item \textbf{heap} — \eng{There is a large heap of tomatoes near the scale.} = Il y a un gros tas de tomates près de la balance. (\eng{tomatoes} = comptable pluriel, tas $\rightarrow$ \eng{heap}.)
    \item \textbf{packet} — \eng{Give me a packet of sugar, please.} = Donne-moi un paquet de sucre, s'il te plaît. (\eng{sugar} = incomptable, emballage souple $\rightarrow$ \eng{packet}.)
    \item \textbf{box} — \eng{He lit the candle with a match from a box of matches.} = Il a allumé la bougie avec une allumette prise dans une boîte d'allumettes. (\eng{matches} = pluriel, petit contenant en carton $\rightarrow$ \eng{box}.)
    \item \textbf{piece} — \eng{She cut a piece of beef for the stew.} = Elle a coupé un morceau de bœuf pour le ragoût. (\eng{beef} = incomptable, portion $\rightarrow$ \eng{piece}.)
    \item \textbf{bags} — \eng{They carried two bags of rice on the motorbike.} = Ils ont transporté deux sacs de riz sur la moto. (\eng{rice} = incomptable, grand sac $\rightarrow$ \eng{bag}.)
\end{enumerate}
\end{exoresolu}

\begin{exercice}[Entraînement personnel — Vocabulaire achat/vente]
\begin{enumerate}
    \item Classez les mots suivants en \textbf{Countable (C)} ou \textbf{Uncountable (U)} : \eng{okra, yam, salt, onion, furniture, advice, mango, garri, smoke, fish (la chair), matches, money}.
    \item Choisissez le bon quantificateur : \eng{much / many} \_\_\_\_\_\_ plantains? \quad \eng{much / many} \_\_\_\_\_\_ rice? \quad \eng{a little / a few} \_\_\_\_\_\_ sugar? \quad \eng{a little / a few} \_\_\_\_\_\_ tomatoes?
    \item Traduisez en anglais :
    \begin{enumerate}
        \item Un sac de maïs de 50 kg.
        \item Trois boîtes de sardines.
        \item Un tas d'okra (okra est incomptable $\rightarrow$ \eng{a heap of okra}).
        \item Du poisson fumé frais (attention à la place de \eng{fresh}).
        \item Ce savon est frelaté.
    \end{enumerate}
    \item Rédaction courte (40--50 mots) : vous êtes au marché. Écrivez une liste de courses pour préparer le \emph{ndolé} et le \emph{bobolo} pour 6 personnes. Utilisez au moins 6 structures \eng{quantity word + of}.
\end{enumerate}
\end{exercice}

\begin{corrige}[Exercice n°1 -- Entraînement personnel]
\begin{enumerate}
    \item \textbf{C} : yam, onion, mango, matches (pluriel). \textbf{U} : okra, salt, furniture, advice, garri, smoke, fish (la chair), money. \textit{Note :} \eng{fish} est comptable seulement pour l'animal entier (\eng{a fish}).
    \item \eng{Many plantains} (C pluriel) — \eng{Much rice} (U) — \eng{A little sugar} (U) — \eng{A few tomatoes} (C pluriel).
    \item
    \begin{enumerate}
        \item A 50-kg bag of maize.
        \item Three tins of sardines.
        \item A heap of okra.
        \item Fresh smoked fish.
        \item This soap is adulterated.
    \end{enumerate}
    \item \textit{Exemple de liste :} \eng{A bundle of bitterleaf, a tuber of cassava (for the bobolo), a packet of groundnut paste, a tin of crayfish, a bottle of palm oil, a heap of onions, a piece of smoked fish, a box of matches and a bar of soap.} (9 structures avec \eng{of}.)
\end{enumerate}
\end{corrige}

\coursec{Prices, payment and bargaining}

Après avoir identifié \emph{ce que} l'on achète (articles et denrées alimentaires) et \emph{où} on l'achète (boutique, marché, supermarché), nous devons maintenant maîtriser le langage de la transaction elle-même : demander le prix, négocier, payer et gérer les imprévus (rupture de stock, hausse des prix). Au Cameroun comme ailleurs, l'achat au marché ou en boutique suit un rituel codifié que l'anglais restitue par un vocabulaire précis et des structures verbales qu'il ne faut pas calquer sur le français.

\courssub{Asking for and stating prices : le vocabulaire du prix}

Dans la section précédente, nous avons vu les noms des produits. Voici maintenant les mots pour parler de leur valeur marchande. Observons d'abord un court échange typique entre un client et une vendeuse au Marché Central de Yaoundé.

\begin{textebox}[Dialogue at the market]
— \eng{“Good morning, Madam. How much is a bag of rice?”}
— \eng{“Good morning. This one is 25,000 francs. The other brand is 23,000.”}
— \eng{“That's too expensive! Prices have gone up again. Can you give me a discount?”}
— \eng{“I can reduce it to 24,000. That is my last price.”}
— \eng{“Okay, it's a bargain. I'll take it. How much for two kilos of tomatoes?”}
— \eng{“They are 1,500 francs a kilo. So 3,000 for two.”}
— \eng{“Here is 5,000 francs.”}
— \eng{“Thank you. Here is your change and your receipt.”}
\end{textebox}

\begin{definition}[Les noms du prix et de la valeur]
\begin{itemize}
    \item \cle{price} « pra-iss » : le prix (étiquette, montant demandé). \eng{The price of fuel has increased.} « Le prix du carburant a augmenté. »
    \item \cle{cost} « kosst » : le coût, la dépense (souvent pour l'acheteur ou la production). \eng{The cost of living is high in Douala.} « Le coût de la vie est élevé à Douala. »
    \item \cle{charge} « tcha-rdj » : frais, tarif (service, livraison, entrée). \eng{There is no charge for delivery.} « La livraison est gratuite (sans frais). »
    \item \cle{fee} « fi » : honoraires, droits (inscription, visa, avocat). \eng{School fees are due in September.} « Les frais de scolarité sont dus en septembre. »
    \item \cle{rate} « reït » : taux, tarif (change, horaire, assurance). \eng{The exchange rate is 650 F for 1 USD.} « Le taux de change est de 650 F pour 1 USD. »
    \item \cle{cheap} « tchip » : bon marché, peu cher. \eng{Second-hand clothes are cheap at Marché B.} « Les friperies sont bon marché au Marché B. »
    \item \cle{expensive} « ik-spen-siv » : cher. \eng{Imported apples are expensive.} « Les pommes importées sont chères. »
    \item \cle{affordable} « a-for-da-beul » : abordable, à la portée du budget. \eng{Local food is more affordable.} « La nourriture locale est plus abordable. »
    \item \cle{costly} « kosst-li » : coûteux (souvent pour un projet, une erreur, un article de luxe). \eng{A costly mistake.} « Une erreur coûteuse. »
\end{itemize}
\end{definition}

\begin{center}
\begin{tabularx}{\linewidth}{|X|X|X|}
\hline
\rowcolor{popblueL} \textbf{Français} & \textbf{English} & \textbf{Remarque / Collocation} \\
\hline
Le prix (étiquette) & \eng{price} & \eng{asking price, half price, price tag} \\
Le coût / La dépense & \eng{cost} & \eng{cost of living, at no cost, cover the cost} \\
Les frais (service) & \eng{charge} & \eng{service charge, delivery charge, free of charge} \\
Les honoraires / droits & \eng{fee} & \eng{tuition fee, entrance fee, lawyer's fee} \\
Le taux (change, intérêt) & \eng{rate} & \eng{exchange rate, interest rate, hourly rate} \\
Bon marché & \eng{cheap} & \eng{dirt cheap (donné), good value for money (bon rapport qualité-prix)} \\
Cher & \cle{expensive} & \eng{a bit pricey (familier), cost an arm and a leg (coûter les yeux de la tête)} \\
Abordable & \cle{affordable} & \eng{reasonable, within budget} \\
\hline
\end{tabularx}
\end{center}

\courssub{Les verbes de la transaction : cost, spend, pay, charge...}

C'est ici que les francophones commettent le plus d'erreurs par traduction littérale. Observons les sujets de ces verbes dans le dialogue ci-dessus : \eng{“This one \textbf{is} 25,000”} (le sujet est la chose), \eng{“I \textbf{'ll take} it”} (le sujet est la personne), \eng{“Prices \textbf{have gone up}”} (sujet = prices).

\begin{propriete}[Cost vs Spend : la règle d'or des sujets]
\begin{itemize}
    \item \textbf{\cle{to cost}} : le sujet est \textbf{la chose} (le prix est une propriété de l'objet). \eng{The bag of rice \textbf{costs} 25,000 F.} (Le sac de riz coûte 25 000 F.) $\rightarrow$ \emph{Jamais : ``I cost 25,000 F''}.
    \item \textbf{\cle{to spend}} : le sujet est \textbf{la personne} (c'est elle qui effectue la dépense). \eng{I \textbf{spent} 25,000 F \textbf{on} rice.} (J'ai dépensé 25 000 F pour du riz.) $\rightarrow$ \emph{Jamais : ``The rice spent me...''}.
    \item \textbf{\cle{to pay}} : sujet = personne. \eng{She \textbf{pays} \textbf{by} mobile money.} Prépositions : \eng{pay \textbf{for} something} (payer qqch), \eng{pay someone} (payer qqn), \eng{pay \textbf{in} cash / \textbf{by} card}.
    \item \textbf{\cle{to charge}} : sujet = vendeur / prestataire. \eng{The shopkeeper \textbf{charges} 500 F for delivery.} (Le commerçant facture 500 F la livraison.)
    \item \textbf{\cle{to sell}} : sujet = vendeur. \eng{They \textbf{sell} plantains \textbf{at} 200 F \textbf{each}.} (Ils vendent les plantains à 200 F pièce.) $\rightarrow$ \eng{sell \textbf{for}} (prix final) : \eng{He sold his car \textbf{for} 2 million.}
    \item \textbf{\cle{to buy}} : sujet = acheteur. \eng{We \textbf{bought} fish \textbf{for} 3,000 F.}
    \item \textbf{\cle{to afford}} : sujet = personne. Sens : ``avoir les moyens de''. Toujours précédé de \eng{can/could/be able to}. \eng{I \textbf{can't afford} a new phone right now.}
\end{itemize}
\end{propriete}

\begin{center}
\begin{tabular}{|l|l|l|l|}
\hline
\rowcolor{poporangeL} \textbf{Verbe} & \textbf{Prétérit} & \textbf{Participe passé} & \textbf{Exemple traduit} \\
\hline
\eng{cost} & \eng{cost} & \eng{cost} & \eng{It cost 1,000 F yesterday.} (Il a coûté 1 000 F hier.) \\
\eng{spend} & \eng{spent} & \eng{spent} & \eng{We spent hours bargaining.} (Nous avons passé des heures à marchander.) \\
\eng{pay} & \eng{paid} & \eng{paid} & \eng{He paid the bill.} (Il a payé l'addition/la facture.) \\
\eng{sell} & \eng{sold} & \eng{sold} & \eng{She sold all her stock.} (Elle a vendu tout son stock.) \\
\eng{buy} & \eng{bought} & \eng{bought} & \eng{They bought yams.} (Ils ont acheté des ignames.) \\
\eng{charge} & \eng{charged} & \eng{charged} & \eng{They charged extra.} (Ils ont facturé un supplément.) \\
\hline
\end{tabular}
\end{center}

\begin{attention}[Pièges majeurs pour francophones]
\begin{enumerate}
    \item \textbf{``How many is it?''} $\rightarrow$ \textbf{FAUX}. Pour demander un prix, on utilise \cle{How much} (combien \emph{d'argent}), car \eng{price} / \eng{money} sont indénombrables. \eng{``\textbf{How much} is this shirt?''} / \eng{``\textbf{How much} does it cost?''}
    \item \textbf{``It costs me 5,000 F''} $\rightarrow$ \textbf{FAUX} (sauf si vous êtes l'objet, ex: ``This mistake cost me my job''). Pour dire ``Je l'ai payé 5 000 F'', on dit : \eng{``I \textbf{paid} 5,000 F \textbf{for} it''} ou \eng{``I \textbf{bought} it \textbf{for} 5,000 F''}.
    \item \textbf{``I spend 5,000 F for rice''} $\rightarrow$ \textbf{FAUX}. Préposition : \eng{spend money \textbf{on} something}. \eng{``I spent 5,000 F \textbf{on} rice.''}
    \item \textbf{``The price is cheap/expensive''} $\rightarrow$ \textbf{Maladroit}. Le prix n'est pas cher ou bon marché, il est \eng{high} (élevé) ou \eng{low} (bas). C'est l'\emph{article} qui est \eng{cheap} ou \eng{expensive}. \eng{``The price is \textbf{high} but the phone is \textbf{affordable}.''}
    \item \textbf{``I afford it''} $\rightarrow$ \textbf{FAUX}. \eng{Afford} s'emploie quasi exclusivement avec un modal : \eng{can/could afford}, \eng{be able to afford}.
\end{enumerate}
\end{attention}

\courssub{Bargaining and negotiating : marchander et négocier}

Au Cameroun, au Nigeria ou au Ghana, marchander (\cle{to bargain} / \cle{to haggle}) est une pratique courante sur les marchés, moins en supermarché où les prix sont fixes (\cle{fixed prices}). Voici le lexique et la méthode.

\begin{definition}[Vocabulaire de la négociation]
\begin{itemize}
    \item \cle{to bargain} (US) / \cle{to haggle} (UK) : marchander, négocier le prix. \eng{Don't haggle in a supermarket.}
    \item \cle{to negotiate} : négocier (plus formel, contrats, salaires, gros volumes).
    \item \cle{to reduce the price} / \cle{to lower the price} : baisser le prix.
    \item \cle{to give a discount} / \cle{to offer a discount} : faire une remise. \eng{A 10\% discount.}
    \item \cle{a bargain} (n.) : une bonne affaire. \eng{It's a real bargain!}
    \item \cle{a fair price} / \cle{a reasonable price} : un prix juste / correct.
    \item \cle{last price} / \cle{final offer} : prix final, dernière offre. \eng{That's my last price, take it or leave it.}
    \item \cle{It's a deal!} / \cle{You've got a deal!} : C'est entendu / Marché conclu !
    \item \cle{to rip someone off} (argot) : arnaquer, faire payer trop cher. \eng{I got ripped off!}
\end{itemize}
\end{definition}

\begin{methode}[Comment négocier un prix en anglais (4 étapes)]
\begin{enumerate}
    \item \textbf{Demander le prix} : \eng{``How much is this?''} / \eng{``What's the price of...?''} / \eng{``How much are you selling this for?''}
    \item \textbf{Réagir (choc / refus)} : \eng{``That's too expensive!''} / \eng{``It costs an arm and a leg!''} / \eng{``That's more than I expected.''} / \eng{``Prices have gone up!''}
    \item \textbf{Proposer / Demander une remise} : \eng{``Can you reduce it?''} / \eng{``Can you give me a discount?''} / \eng{``Will you take 20,000?''} / \eng{``What's your last price?''}
    \item \textbf{Conclure ou partir} :
    \begin{itemize}
        \item Accord : \eng{``Okay, it's a deal.''} / \eng{``Sold!''} / \eng{``I'll take it.''}
        \item Désaccord : \eng{``No, thank you. I'll leave it.''} / \eng{``I can't afford that.''}
    \end{itemize}
\end{enumerate}
\end{methode}

\begin{exemplebox}[Exemples de négociation traduits]
\eng{``Can you do it for 15,000?''} = « Tu peux me le faire à 15 000 ? » (Proposition directe)
\eng{``If I buy three, will you give me a discount?''} = « Si j'en prends trois, tu me fais une remise ? » (Conditionnel)
\eng{``That's a bit steep for me.''} = « C'est un peu raide / trop cher pour moi. » (Familier, poli)
\eng{``Meet me halfway?''} = « On partage la différence ? » (Idiome : \eng{meet halfway})
\eng{``No discount today, madam. Fixed price.''} = « Pas de remise aujourd'hui, madame. Prix fixe. »
\end{exemplebox}

\courssub{Payment : paiement, monnaie, reçus et dettes}

Une fois le prix convenu (\cle{agreed price}), vient l'étape du paiement. Au Cameroun, le \cle{cash} (espèces) et le \cle{Mobile Money} (MoMo, Orange Money, MTN MoMo) sont rois.

\begin{definition}[Moyens de paiement et documents]
\begin{itemize}
    \item \cle{to pay in cash} : payer en espèces / cash. \eng{Do you accept cash only?}
    \item \cle{to pay by mobile money} / \cle{to pay via MoMo} : payer par Mobile Money. \eng{Can I pay by Mobile Money?}
    \item \cle{to pay by card} / \cle{by bank transfer} : payer par carte / virement.
    \item \cle{change} (n. indénombrable) : \textbf{la monnaie rendue} (pas \eng{money}!). \eng{``Here is your \textbf{change}.''} $\neq$ \eng{money} (l'argent en général).
    \item \cle{receipt} « ri-sit » : le reçu, le ticket de caisse (le \emph{p} ne se prononce pas). \eng{Keep the receipt for the warranty.}
    \item \cle{bill} / \cle{invoice} : la facture, l'addition (restaurant). \eng{Can I have the bill, please?}
    \item \cle{debt} « det » : la dette. \eng{He is in debt.}
    \item \cle{credit} (n.) : le crédit (acheter à crédit). \eng{To buy on credit.}
    \item \cle{to owe} (verbe irrégulier : owe, owed, owed) « o » : devoir de l'argent (à qqn). \eng{I \textbf{owe} him 5,000 F.} (Je lui dois 5 000 F.) $\rightarrow$ Attention : \eng{I owe \textbf{him} money}, pas \eng{to him}.
\end{itemize}
\end{definition}

\begin{exemplebox}[Paiement et monnaie : exemples traduits]
\eng{``The total comes to 12,500 francs.''} = « Le total fait 12 500 francs. »
\eng{``Here is 20,000. Keep the change.''} = « Voici 20 000. Gardez la monnaie. » (Pourboire)
\eng{``I don't have enough change. Do you have smaller notes?''} = « Je n'ai pas assez de monnaie. Vous avez des petites coupures ? »
\eng{``Can you send me the MoMo number?''} = « Tu peux m'envoyer le numéro MoMo ? »
\eng{``He still owes me for the yams.''} = « Il me doit encore pour les ignames. »
\end{exemplebox}

\courssub{Stock issues and price changes : ruptures et hausses de prix}

Dans un contexte d'inflation ou de saisonnalité (saison des pluies, fêtes de fin d'année), le vocabulaire de la rareté et de la variation des prix est indispensable.

\begin{propriete}[Expressions de rupture et de variation]
\begin{itemize}
    \item \cle{to run out of stock} / \cle{to be out of stock} : être en rupture de stock. \eng{We have \textbf{run out of} frozen chicken.}
    \item \cle{shortage} (n.) : pénurie. \eng{There is a \textbf{shortage} of flour.}
    \item \cle{scarcity} (n.) : rareté (plus formel). \eng{Fuel \textbf{scarcity} affects transport costs.}
    \item \cle{prices have gone up} / \cle{prices have risen} / \cle{prices have increased} : les prix ont augmenté. \eng{``Prices \textbf{have gone up} because of \textbf{scarcity}.''} (Les prix ont augmenté à cause de la pénurie.)
    \item \cle{prices have dropped} / \cle{fallen} / \cle{decreased} : les prix ont baissé.
    \item \cle{price hike} (n.) : hausse brutale des prix. \eng{A sudden \textbf{price hike} on cement.}
    \item \cle{inflation} (n.) : inflation. \eng{High \textbf{inflation} erodes purchasing power.}
\end{itemize}
\end{propriete}

\begin{savaistu}[Mobile Money au Cameroun : un vocabulaire spécifique]
Au Cameroun, le Mobile Money (MoMo) a généré son propre jargon anglais/camfranglais :
\begin{itemize}
    \item \eng{``Send me the money.''} / \eng{``Transfer to this number.''}
    \item \eng{``I've sent it. Check your balance.''}
    \item \eng{``Transaction successful / failed.''}
    \item \eng{``Agent / Merchant''} : le dépositaire / le commerçant agréé.
    \item \eng{``Cash out''} : retirer de l'argent (chez l'agent).
    \item \eng{``Cash in''} : déposer de l'argent.
\end{itemize}
Contrairement à la France où l'on dit souvent ``faire un virement'', ici on dit couramment \eng{``Do a transfer''} ou \eng{``Send money''}.
\end{savaistu}

\begin{popfigure}
\begin{tikzpicture}[
    node distance=1.2cm and 1.5cm,
    every node/.style={font=\small, align=center},
    box/.style={draw=popblue, thick, fill=popblueL, rounded corners, minimum width=2.8cm, minimum height=1.2cm},
    arrow/.style={-Stealth, thick, popdark}
]
\node[box, align=center] (ask){Ask the price\\ \eng{``How much is it?''}};
\node[box, right=of ask, align=center] (bargain){Bargain / Negotiate\\ \eng{``Can you reduce it?''}};
\node[box, right=of bargain, align=center] (agree){Agree on price\\ \eng{``It's a deal!''}};
\node[box, right=of agree, align=center] (pay){Pay\\ \eng{Cash / MoMo / Card}};
\node[box, right=of pay, align=center] (receive){Receive change\\ \& Receipt\\ \eng{``Here is your change.''}};

\draw[arrow] (ask) -- (bargain);
\draw[arrow] (bargain) -- (agree);
\draw[arrow] (agree) -- (pay);
\draw[arrow] (pay) -- (receive);
\end{tikzpicture}
\legende{Organigramme d'une transaction commerciale type : de la demande de prix à la remise du reçu.}
\end{popfigure}

\courssub{Exercice résolu : Mettre en pratique}

\begin{exoresolu}[Compléter le dialogue avec les mots appropriés]
\textbf{Énoncé :} Complétez l'échange entre un client (C) et un vendeur (V) au Marché Mfoundi. Utilisez : \eng{change, receipt, costs, expensive, discount, bargain, owe, spend, pay, price}.

V: Good morning. What can I get you?
C: \eng{How much} \_\_\_\_\_\_ a kilo of onions?
V: It \_\_\_\_\_\_ 800 francs.
C: 800? That's too \_\_\_\_\_\_! Last week the \_\_\_\_\_\_ was 600.
V: Prices have gone up, sir. Shortage in the North.
C: Can you give me a \_\_\_\_\_\_? I'll take 5 kilos.
V: Okay, 3,500 for five. That's a \_\_\_\_\_\_.
C: Deal. I'll \_\_\_\_\_\_ by Mobile Money. \textit{(Transaction done)}
V: Thank you. Here is your \_\_\_\_\_\_ and your \_\_\_\_\_\_.
C: Thanks. By the way, my brother still \_\_\_\_\_\_ you 2,000 F from last month.
V: Tell him to come \_\_\_\_\_\_ it soon!

\tcblower
\textbf{Solution détaillée :}
\begin{enumerate}
    \item \eng{is} (verbe \eng{be} au présent, 3e pers. sg., sujet \eng{a kilo}). \textit{Alternative possible avec deux trous :} \eng{does ... cost}.
    \item \eng{costs} (verbe \eng{cost}, sujet \eng{It} = la chose, 3e pers. sg.).
    \item \eng{expensive} (adjectif après \eng{too}).
    \item \eng{price} (nom, sujet de \eng{was}).
    \item \eng{discount} (collocation \eng{give a discount}).
    \item \eng{bargain} (nom, \eng{That's a bargain}).
    \item \eng{pay} (infinitif après \eng{'ll} / \eng{will}).
    \item \eng{change} (monnaie rendue, indénombrable).
    \item \eng{receipt} (ticket de caisse).
    \item \eng{owes} (verbe \eng{to owe}, présent, 3e pers. sg., sujet = \eng{brother}).
    \item \eng{pay} (infinitif après \eng{to} : \eng{come pay it} / \eng{come and pay it} / \eng{come to pay it}).
\end{enumerate}
\textbf{Traduction du dialogue complet :}
V: Bonjour. Que puis-je vous servir ?
C: Combien coûte le kilo d'oignons ?
V: Il coûte 800 francs.
C: 800 ? C'est trop cher ! La semaine dernière le prix était 600.
V: Les prix ont monté, monsieur. Pénurie dans le Nord.
C: Vous pouvez me faire une remise ? J'en prends 5 kilos.
V: D'accord, 3 500 pour cinq. C'est une affaire.
C: Marché conclu. Je paie par Mobile Money.
V: Merci. Voici votre monnaie et votre reçu.
C: Merci. Au fait, mon frère vous doit encore 2 000 F du mois dernier.
V: Dites-lui de venir payer ça bientôt !
\end{exoresolu}

\courssub{Exercices d'entraînement}

\begin{exercice}[Entraînement : Prix, Paiement, Négociation]
\textbf{A. Choisissez le bon mot (cost / spend / pay / charge / owe). Conjuguez si nécessaire.}
\begin{enumerate}
    \item This designer bag \_\_\_\_\_\_ a fortune. I \_\_\_\_\_\_ not \_\_\_\_\_\_ that much on a handbag.
    \item The mechanic \_\_\_\_\_\_ me 15,000 F to repair the bendskin. I \_\_\_\_\_\_ him by card.
    \item How much did you \_\_\_\_\_\_ on fuel last month? — I \_\_\_\_\_\_ about 50,000 F.
    \item My neighbour \_\_\_\_\_\_ me 10,000 F since January. He still hasn't \_\_\_\_\_\_ me back.
\end{enumerate}

\textbf{B. Corrigez les erreurs dans les phrases suivantes (grammaire, vocabulaire, préposition).}
\begin{enumerate}
    \item How many is this bunch of plantains?
    \item The price is cheap today.
    \item I cost 2,000 F for this t-shirt.
    \item Can you make a discount? I haven't many money.
    \item Here is your money (le vendeur rend 500 F au client).
    \item I afford not this laptop.
    \item The tomatoes are out of stock since Monday.
\end{enumerate}

\textbf{C. Expression écrite :} Imaginez une scène au supermarché \eng{``Casino''} ou \eng{``Mahima''} à Douala. Vous voulez acheter du riz, de l'huile et du savon. Le total est 12 500 F. Vous payez avec un billet de 20 000 F. Écrivez le dialogue (8-10 répliques) en utilisant : \eng{How much...}, \eng{total}, \eng{pay by cash}, \eng{change}, \eng{receipt}, \eng{Have a nice day}.
\end{exercice}

\begin{corrige}[Exercice : Prix, Paiement, Négociation]
\textbf{A. Verbes conjugués :}
\begin{enumerate}
    \item \eng{costs} (sujet \eng{bag}) / \eng{don't spend} (sujet \eng{I}, présent simple négatif) \textit{ou} \eng{can't afford to spend} (si trois trous).
    \item \eng{charged} (sujet \eng{mechanic}, prétérit) / \eng{paid} (sujet \eng{I}, prétérit).
    \item \eng{spend} (base verbale après \eng{did}) / \eng{spent} (prétérit de \eng{spend}).
    \item \eng{owes} (présent, 3e pers. sg., verbe d'état) / \eng{paid} (participe passé après \eng{hasn't}).
\end{enumerate}

\textbf{B. Corrections :}
\begin{enumerate}
    \item \eng{How \textbf{much} is this bunch...?} (\eng{much} pour prix/argent indénombrable).
    \item \eng{The price is \textbf{low} / \textbf{reasonable} today.} (Ou : \eng{It's \textbf{cheap} today.} $\rightarrow$ l'article est bon marché, pas le prix).
    \item \eng{\textbf{It costs} 2,000 F.} / \eng{\textbf{I paid} 2,000 F \textbf{for} this t-shirt.} / \eng{\textbf{I bought} this t-shirt \textbf{for} 2,000 F.}
    \item \eng{Can you \textbf{give me a discount}? I \textbf{don't have much} money.} (\eng{give a discount}, \eng{much} pour indénombrable \eng{money}).
    \item \eng{Here is your \textbf{change}.} (\eng{change} = monnaie rendue).
    \item \eng{\textbf{I can't afford} this laptop.} (\eng{afford} + modal \eng{can}).
    \item \eng{The tomatoes \textbf{have been} out of stock \textbf{since} Monday.} (Present perfect \eng{have been} pour durée depuis le passé).
\end{enumerate}

\textbf{C. Proposition de dialogue (modèle) :}
\eng{Customer:} Good afternoon. \eng{How much} is this 20kg bag of rice?
\eng{Cashier:} It's 25,000 F. The oil is 6,000 and the soap 1,500.
\eng{Customer:} I'll take one of each. What is the \eng{total}?
\eng{Cashier:} The \eng{total} comes to 32,500 F.
\eng{Customer:} Okay. I'll \eng{pay by cash}. Here is 40,000 F.
\eng{Cashier:} Thank you. Here is your \eng{change} (7,500 F) and your \eng{receipt}.
\eng{Customer:} Thank you. \eng{Have a nice day!}
\eng{Cashier:} You too, sir. Goodbye!
\end{corrige}

\coursec{Collocations, idioms and functional expressions}

\courssub{Transition depuis « Prices, payment and bargaining »}

Dans la section précédente, nous avons étudié le vocabulaire isolé des prix (\cle{price}, \cle{cost}, \cle{change}, \cle{receipt}) et des moyens de paiement (\cle{cash}, \cle{credit card}, \cle{instalments}). Cependant, en situation réelle de communication — au marché Mokolo à Yaoundé, dans un \eng{supermarket} à Douala ou dans une \eng{corner shop} à Buea — les mots ne s'emploient pas seuls. Ils s'associent en \textbf{collocations} (associations lexicales fréquentes et attendues), en \textbf{expressions idiomatiques} (sens global différent de la somme des mots) et en \textbf{formules fonctionnelles} (rituels de politesse et de transaction). Maîtriser ces blocs prêts à l'emploi, c'est passer du « vocabulaire de dictionnaire » à l'anglais naturel, fluide et socialement approprié — exactement ce qu'exigent les épreuves de compréhension et d'expression de la classe de Première.

\begin{definition}[Collocation, idiome, formule fonctionnelle]
\begin{itemize}
\item Une \textbf{collocation} est une association habituelle de mots : \eng{to raise prices} (et non \emph{to rise prices}), \eng{a high price} (et non \emph{a big price}).
\item Un \textbf{idiome} est une expression figée dont le sens est imagé : \eng{to cost a fortune} = coûter très cher.
\item Une \textbf{formule fonctionnelle} est une phrase toute faite remplissant une fonction sociale : \eng{Can I help you?} = accueillir un client.
\end{itemize}
\end{definition}

\courssub{Collocations avec \eng{buy} et \eng{sell} : au-delà du sens de base}

\begin{textebox}{Dialogue au marché de Nkolbisson — Observation des collocations}
— Good morning, Mama Ngo. I'd like to buy in bulk today. Ten bags of rice and twenty litres of oil.

— Ah, my regular customer! Since you buy in bulk, I can give you a discount. I usually sell at a profit, but for you I'll reduce my margin.

— Thank you. By the way, do you still have that brand of spaghetti? Last week you were sold out.

— Yes, the new stock arrived yesterday. But hurry, they sell out fast. Oh, and if you're short of cash, you can buy on credit and pay me next market day.

— No need, I have enough. Let's agree on a price.
\end{textebox}

\begin{propriete}[Collocations fréquentes avec \eng{buy} et \eng{sell}]
\begin{center}
\begin{tabularx}{\linewidth}{|X|X|X|}
\hline
\rowcolor{popblueL} \textbf{Collocation} & \textbf{Sens (FR)} & \textbf{Exemple + Traduction} \\
\hline
\eng{to buy in bulk} & acheter en gros, en grande quantité & \eng{We buy maize in bulk during the harvest season.} (Nous achetons le maïs en gros pendant la récolte.) \\
\hline
\eng{to sell at a profit} & vendre avec bénéfice & \eng{He sells his cocoa at a profit.} (Il vend son cacao avec bénéfice.) \\
\hline
\eng{to sell at a loss} & vendre à perte & \eng{The shopkeeper sold the damaged goods at a loss.} (Le commerçant a vendu la marchandise abîmée à perte.) \\
\hline
\eng{to sell out} (intransitif) & être épuisé, tout vendre & \eng{The tickets sold out in an hour.} (Les billets ont été épuisés en une heure.) \\
\hline
\eng{to buy on credit} & acheter à crédit & \eng{Many farmers buy fertiliser on credit.} (Beaucoup d'agriculteurs achètent l'engrais à crédit.) \\
\hline
\eng{to buy for cash} & acheter au comptant & \eng{I bought this radio for cash.} (J'ai acheté cette radio au comptant.) \\
\hline
\end{tabularx}
\end{center}
Rappel de conjugaison : \eng{buy -- bought -- bought} ; \eng{sell -- sold -- sold} ; \eng{cost -- cost -- cost}.
\end{propriete}

\begin{attention}[Piège majeur : \eng{to sell out} n'est pas « vendre dehors »]
\eng{To sell out} est ici un verbe \textbf{intransitif} (pas de COD après) qui signifie « être entièrement vendu, être en rupture de stock ». Le sujet est la \emph{marchandise}, pas le vendeur.
\begin{itemize}
\item \eng{The bread is sold out.} = Le pain est épuisé (il n'y en a plus).
\item \eng{We sold out the bread.} $\leftarrow$ \textbf{INCORRECT}. On dit : \eng{We have sold out of bread.} ou \eng{We have sold all the bread.}
\end{itemize}
Ne pas confondre avec le sens figuré transitif : \eng{He sold out to the big corporation.} = « Il a trahi ses principes / vendu son âme à la grande société. »
\end{attention}

\courssub{Collocations avec \eng{price}, \eng{cost} et \eng{value} : parler du prix comme un pro}

\begin{exemplebox}[Observation : le prix dans la presse et la conversation]
\begin{itemize}
\item \eng{The government has fixed a minimum price for cocoa.} (Le gouvernement a fixé un prix minimum pour le cacao.)
\item \eng{Supermarkets often raise their prices before Christmas.} (Les supermarchés augmentent souvent leurs prix avant Noël.)
\item \eng{This phone offers excellent value for money.} (Ce téléphone offre un excellent rapport qualité-prix.)
\item \eng{It costs a fortune to import these goods.} (Cela coûte une fortune d'importer ces marchandises.)
\end{itemize}
\end{exemplebox}

\begin{propriete}[Collocations clés : \eng{price}, \eng{cost}, \eng{value}]
\begin{center}
\begin{tabularx}{\linewidth}{|X|X|X|}
\hline
\rowcolor{popblueL} \textbf{Collocation} & \textbf{Sens / Synonyme} & \textbf{Remarque} \\
\hline
\eng{to fix / set a price} & fixer un prix & Prix décidé par l'État ou par le vendeur. \\
\hline
\eng{to raise / increase prices} & augmenter les prix & Antonymes : \eng{to lower / cut / reduce prices}. \\
\hline
\eng{prices rise / go up} & les prix augmentent & \eng{rise -- rose -- risen} : verbe \textbf{intransitif}. \\
\hline
\eng{high / low / reasonable prices} & prix élevés / bas / raisonnables & \eng{affordable prices} = prix abordables. \\
\hline
\eng{value for money} (n.) & bon rapport qualité-prix & Expression invariable : \eng{good value for money}. \\
\hline
\eng{to cost a fortune} & coûter très cher & Familier. Synonyme : \eng{to cost an arm and a leg}. \\
\hline
\eng{to cost a pretty penny} & coûter bonbon, coûter cher & Registre plaisant, un peu vieilli. \\
\hline
\end{tabularx}
\end{center}
\end{propriete}

\begin{attention}[\eng{raise} ou \eng{rise} ? Le piège du transitif / intransitif]
\begin{itemize}
\item \eng{\textbf{To raise}} (raise -- raised -- raised) est \textbf{transitif} : il a un COD. \eng{The shopkeeper has raised his prices.} (Le commerçant a augmenté ses prix.)
\item \eng{\textbf{To rise}} (rise -- rose -- risen) est \textbf{intransitif} : pas de COD. \eng{Prices have risen.} (Les prix ont augmenté.)
\item Jamais \eng{prices have raised} $\leftarrow$ \textbf{INCORRECT}. Dites : \eng{prices have risen} ou \eng{prices have been raised}.
\end{itemize}
\end{attention}

\begin{definition}[Value for money]
\eng{Value for money} (nom) : \textbf{good quality compared to the price paid} (bonne qualité par rapport au prix payé). C'est l'équivalent exact de « \textbf{rapport qualité-prix} ». On ne dit jamais \eng{quality-price ratio} (calque du français).
\begin{itemize}
\item \eng{This local fabric is great value for money.} (Ce tissu local a un excellent rapport qualité-prix.)
\item \eng{Is it good value for money?} (Est-ce que cela vaut son prix ?)
\end{itemize}
\end{definition}

\courssub{Collocations avec \eng{shop} et \eng{shopping} : l'acte d'acheter}

\begin{popfigure}
\begin{tikzpicture}[
    mindmap,
    concept color=popblue,
    level 1 concept/.append style={level distance=4.2cm, sibling angle=90},
    every node/.style={font=\small},
    grow cyclic
]
\node[concept] {shop / shopping}
    child[concept color=popgreen] { node[concept, align=center] {Verbes\\ \eng{go shopping}\\ \eng{do the shopping}\\ \eng{shop around}} }
    child[concept color=poporange] { node[concept, align=center] {Noms composés\\ \eng{shopping list}\\ \eng{shopping basket}\\ \eng{shopping centre}\\ \eng{shopping spree}} }
    child[concept color=poppurple] { node[concept, align=center] {Personnes et lieux\\ \eng{shopkeeper}\\ \eng{shop assistant}\\ \eng{shop window}\\ \eng{corner shop}} }
    child[concept color=poppink] { node[concept, align=center] {Expressions\\ \eng{window-shopping}\\ \eng{shop till you drop}\\ \eng{one-stop shop}} };
\end{tikzpicture}
\legende{Carte mentale du champ lexical \eng{shop / shopping}}
\end{popfigure}

\begin{propriete}[Distinction cruciale : \eng{go shopping} ou \eng{do the shopping} ?]
\begin{itemize}
\item \eng{\textbf{Go shopping}} = faire les magasins (loisir, achats non alimentaires le plus souvent). \eng{Let's go shopping in Douala this Saturday.} (Allons faire les magasins à Douala samedi.)
\item \eng{\textbf{Do the shopping}} = faire les courses (achats courants, alimentaires). \eng{I do the shopping every morning at the market.} (Je fais les courses au marché chaque matin.)
\item \eng{\textbf{Window-shopping}} = faire du lèche-vitrine, regarder sans acheter. Voir la boîte \emph{Le saviez-vous ?} ci-dessous.
\end{itemize}
\end{propriete}

\begin{savaistu}[Window shopping : une tradition du samedi]
Au Royaume-Uni et aux États-Unis, \eng{window shopping} (ou \eng{to go window-shopping}) est une activité de loisir très populaire le week-end. On se promène dans les rues commerçantes (la \eng{High Street} en Grande-Bretagne) ou les centres commerciaux (\eng{shopping malls}) pour regarder les vitrines (\eng{shop windows}) et comparer les tendances, sans intention d'acheter immédiatement.
\begin{itemize}
\item \eng{We spent the afternoon window-shopping in Oxford Street.} (Nous avons passé l'après-midi à faire du lèche-vitrine dans Oxford Street.)
\item Orthographe : le nom s'écrit \eng{window shopping} (sans trait d'union) ; le verbe s'écrit \eng{to window-shop} (avec trait d'union). \eng{Shopping} reste invariable : jamais de \eng{-s}.
\end{itemize}
\end{savaistu}

\courssub{Expressions idiomatiques du commerce : l'anglais « couleur locale »}

\begin{exemplebox}[Idiomes courants : sens imagé et traduction]
\begin{center}
\begin{tabularx}{\linewidth}{|X|X|X|}
\hline
\rowcolor{poporangeL} \textbf{Idiome} & \textbf{Image} & \textbf{Sens réel + Exemple traduit} \\
\hline
\eng{a rip-off} & « un arrachage » & \textbf{Une arnaque, un prix exorbitant.} \eng{5,000 francs for a pen? That's a rip-off!} (5 000 F pour un stylo ? C'est de l'arnaque !) \\
\hline
\eng{a good deal} / \eng{a bargain} & « une bonne affaire » & \textbf{Un achat avantageux.} \eng{I got a good deal on this laptop.} (J'ai fait une bonne affaire avec ce portable.) \\
\hline
\eng{dirt cheap} & « bon marché comme la terre » & \textbf{Vraiment pas cher, donné.} \eng{These mangoes are dirt cheap in season.} (Ces mangues sont données en saison.) \\
\hline
\eng{to cost a fortune} & « coûter une fortune » & \textbf{Coûter extrêmement cher.} \eng{Imported cheese costs a fortune here.} (Le fromage importé coûte une fortune ici.) \\
\hline
\eng{It's a deal!} & « C'est un marché ! » & \textbf{Marché conclu !} (accord sur le prix). \\
\hline
\eng{to shop around} & « faire le tour des magasins » & \textbf{Comparer les prix avant d'acheter.} \eng{Don't buy the first one; shop around first.} (N'achète pas le premier ; compare d'abord.) \\
\hline
\end{tabularx}
\end{center}
\end{exemplebox}

\begin{attention}[Faux amis et calques à proscrire]
\begin{itemize}
\item \eng{\textbf{Deal}} : en anglais standard, \eng{a business deal} = un accord commercial, \eng{a good deal} = une bonne affaire. Ne pas employer \eng{deal} dans le sens argotique français de « contrat louche ».
\item \eng{\textbf{Cheap}} = « bon marché » (prix bas). Ce n'est \textbf{pas} un synonyme de « de mauvaise qualité ». Pour « de mauvaise qualité », dites \eng{poor quality}, \eng{shoddy} ou \eng{flimsy}. Jamais \eng{cheap quality}.
\item \eng{\textbf{Expensive}} (cher) $\neq$ \eng{\textbf{expansive}} (vaste, étendu ; expansif en parlant d'une personne). Erreur classique d'orthographe et de prononciation : « èk-span-siv » contre « èk-span-siv » — attention à l'orthographe à l'écrit.
\item \eng{\textbf{To cost}} (verbe irrégulier : cost -- cost -- cost) $\neq$ \eng{\textbf{cost}} (nom = le coût). \eng{The cost is high.} (Le coût est élevé.) / \eng{It costs a lot.} (Cela coûte cher.)
\item \eng{\textbf{Economic}} (économique, relatif à l'économie) $\neq$ \eng{\textbf{economical}} (économe, qui fait faire des économies). \eng{An economical car} = une voiture économe en carburant.
\end{itemize}
\end{attention}

\courssub{Formules fonctionnelles : le rituel vendeur / client}

La transaction suit un \textbf{script social} codifié. Maîtriser ces formules, c'est être poli, efficace et naturel — une compétence directement évaluée à l'oral comme à l'écrit (dialogues à compléter, mises en situation).

\begin{popfigure}
\begin{tikzpicture}[
    every node/.style={font=\small, align=center},
    seller/.style={rectangle, draw=popblue, thick, fill=popblueL, rounded corners, minimum width=6.2cm, minimum height=1.1cm},
    customer/.style={rectangle, draw=popgreen, thick, fill=popgreenL, rounded corners, minimum width=6.2cm, minimum height=1.1cm}
]
\node[font=\bfseries, text=popblue] at (-3.6, 0.6) {THE SELLER SAYS};
\node[font=\bfseries, text=popgreen] at (3.6, 0.6) {THE CUSTOMER SAYS};

\node[seller, align=center] at (-3.6, -0.6){\eng{“Can I help you?”}\\ \eng{“What can I do for you?”}};
\node[customer, align=center] at (3.6, -0.6){\eng{“I'm just looking, thanks.”}\\ \eng{“I'd like [item], please.”}};

\node[seller, align=center] at (-3.6, -2.0){\eng{“Here you are.”}\\ \eng{“Anything else?”}};
\node[customer, align=center] at (3.6, -2.0){\eng{“Do you have [item]?”}\\ \eng{“Have you got [item]?”}};

\node[seller, align=center] at (-3.6, -3.4){\eng{“That comes to [price].”}\\ \eng{“Would you like a bag?”}};
\node[customer, align=center] at (3.6, -3.4){\eng{“How much is it / are they?”}\\ \eng{“I'll take it / them.”}};

\node[seller, align=center] at (-3.6, -4.8){\eng{“We are out of stock.”}\\ \eng{“Sorry, we don't have any left.”}};
\node[customer, align=center] at (3.6, -4.8){\eng{“Can I have a discount?”}\\ \eng{“Is that your final price?”}};

\node[seller, align=center] at (-3.6, -6.2){\eng{“Cash or card?”}\\ \eng{“Your receipt is in the bag.”}};
\node[customer, align=center] at (3.6, -6.2){\eng{“Cash, please.” / “By card.”}\\ \eng{“Keep the change.”}};
\end{tikzpicture}
\legende{Échange type : formules du vendeur (bleu) et du client (vert)}
\end{popfigure}

\begin{propriete}[Analyse des formules : forme, emploi, registre]
\begin{center}
\begin{tabularx}{\linewidth}{|l|X|X|X|}
\hline
\rowcolor{popblueL} \textbf{Locuteur} & \textbf{Formule} & \textbf{Fonction} & \textbf{Registre / Note} \\
\hline
Vendeur & \eng{Can I help you?} & Accueillir, ouvrir l'échange & Standard, poli. \\
\hline
Vendeur & \eng{We are out of stock.} & Signaler une rupture de stock & \eng{Out of stock} = locution adjectivale. \\
\hline
Vendeur & \eng{That comes to 5,000 francs.} & Annoncer le total & Plus naturel que \eng{The total is...}. \\
\hline
Client & \eng{I'm just looking, thanks.} & Refuser l'aide poliment & Indispensable pour rester courtois. \\
\hline
Client & \eng{I'd like... / I'll have...} & Commander, exprimer un choix & \eng{I'd like} = \eng{I would like} (poli). \eng{I want} = trop direct. \\
\hline
Client & \eng{Do you have...? / Have you got...?} & Demander la disponibilité & \eng{Have you got} = britannique, oral. \\
\hline
Client & \eng{I'll take it / them.} & Valider l'achat & \eng{Will} : décision prise au moment de parler. \\
\hline
Client & \eng{Can I have a discount?} & Négocier (au marché) & En grand magasin : \eng{Is there any discount?} \\
\hline
\end{tabularx}
\end{center}
\end{propriete}

\courssub{Exercice résolu : de l'observation à la production}

\textbf{Contexte :} au supermarché de Bastos (Yaoundé), Mme Mballa fait ses courses hebdomadaires. Lisez d'abord le dialogue à trous, puis essayez de le compléter avant de consulter la solution.

\begin{textebox}{Dialogue à trous — At the supermarket}
\textbf{Assistant:} Good morning, madam. \_\_\_\_\_\_\_\_\_\_\_\_ ?

\textbf{Mrs Mballa:} \_\_\_\_\_\_\_\_\_\_\_\_, I have my \_\_\_\_\_\_\_\_\_\_\_\_. I need rice, oil and tomato paste.

\textbf{Assistant:} The rice is in aisle 3. \_\_\_\_\_\_\_\_\_\_\_\_ ?

\textbf{Mrs Mballa:} Yes, \_\_\_\_\_\_\_\_\_\_\_\_ a 25-kilo bag of “Mama Africa” rice, please.

\textbf{Assistant:} Certainly. \_\_\_\_\_\_\_\_\_\_\_\_ ?

\textbf{Mrs Mballa:} Let me check... Oh, the tomato paste is \_\_\_\_\_\_\_\_\_\_\_\_ ! I'll take the “Gino” brand instead.

\textbf{Assistant:} Good choice. That \_\_\_\_\_\_\_\_\_\_\_\_ 12,500 francs. \_\_\_\_\_\_\_\_\_\_\_\_ ?

\textbf{Mrs Mballa:} Cash. Here is 15,000. \_\_\_\_\_\_\_\_\_\_\_\_.

\textbf{Assistant:} Here is your change and your receipt. \_\_\_\_\_\_\_\_\_\_\_\_ !

\textbf{Mrs Mballa:} Thank you. Goodbye.
\end{textebox}

\begin{exoresolu}[Corrigé du dialogue à trous]
Complétez le dialogue ci-dessus avec les collocations et formules appropriées de la leçon.
\tcblower
\textbf{Solution détaillée et justifications :}
\begin{enumerate}
\item \eng{Can I help you?} (ou \eng{What can I do for you?}) — formule d'accueil standard du vendeur.
\item \eng{No, thank you, I'm just looking} — politesse : le client signale qu'il n'a pas besoin d'aide immédiate. (Accepter aussi \eng{I'm just looking, thanks}.)
\item \eng{shopping list} — collocation fixe : \eng{shopping list} (liste de courses). Jamais \eng{list of shopping}.
\item \eng{Anything else?} — relance du vendeur après un premier article. (Ici : après avoir indiqué le rayon, le vendeur propose son aide pour la suite.)
\item \eng{I'd like} — formule de commande polie (\eng{I would like}). Éviter \eng{I want}, trop direct.
\item \eng{sold out} — collocation intransitive : le produit est épuisé. Sujet singulier \eng{the tomato paste} $\rightarrow$ \eng{is sold out}.
\item \eng{comes to} — verbe à particule \eng{come to} = « s'élever à, faire un total de ». Très courant pour annoncer un total.
\item \eng{Cash or card?} — formule elliptique standard pour demander le mode de paiement.
\item \eng{Keep the change.} — formule pour dire « gardez la monnaie ».
\item \eng{Have a nice day!} (ou \eng{Thank you, goodbye.}) — formule de clôture.
\end{enumerate}
\end{exoresolu}

\begin{aretenir}[Check-list des points clés de la section]
\begin{itemize}
\item \textbf{Collocations \eng{buy/sell}} : \eng{buy in bulk}, \eng{sell at a profit / at a loss}, \eng{sell out} (intransitif = être épuisé), \eng{buy on credit}, \eng{buy for cash}.
\item \textbf{Collocations \eng{price}} : \eng{fix / set a price}, \eng{raise prices} (transitif) vs \eng{prices rise} (intransitif), \eng{value for money} (rapport qualité-prix).
\item \textbf{Collocations \eng{shop}} : \eng{go shopping} (loisir) vs \eng{do the shopping} (courses), \eng{shopping list}, \eng{window-shopping}.
\item \textbf{Idiomes} : \eng{a rip-off} (arnaque), \eng{dirt cheap} (donné), \eng{cost a fortune} (très cher), \eng{It's a deal!} (marché conclu), \eng{shop around} (comparer).
\item \textbf{Formules vendeur} : \eng{Can I help you?}, \eng{Anything else?}, \eng{We are out of stock}, \eng{That comes to...}, \eng{Cash or card?}
\item \textbf{Formules client} : \eng{I'm just looking}, \eng{I'd like...}, \eng{Do you have...?}, \eng{I'll take it}, \eng{Keep the change}.
\item \textbf{Pièges} : \eng{sell out} $\neq$ « vendre dehors » ; \eng{cheap} = bon marché (jamais \eng{cheap quality}) ; \eng{I'd like} et non \eng{I want} ; \eng{prices rise} et non \eng{prices raise}.
\end{itemize}
\end{aretenir}

\courssub{Mise en pratique guidée : exercices d'application}

\begin{exercice}[Entraînement 1 : Collocations et choix du mot juste]
\textbf{Complétez les phrases avec la collocation correcte (forme verbale adaptée).}
\begin{enumerate}
\item Since the harvest was bad, the price of plantains has \_\_\_\_\_\_\_\_\_\_\_\_ in all the markets of Yaoundé.
\item My uncle \_\_\_\_\_\_\_\_\_\_\_\_ his old car \_\_\_\_\_\_\_\_\_\_\_\_ ; he needed money urgently.
\item Did you \_\_\_\_\_\_\_\_\_\_\_\_ at that new mall? No, everything there \_\_\_\_\_\_\_\_\_\_\_\_.
\item This local wine is excellent \_\_\_\_\_\_\_\_\_\_\_\_ ; it tastes great and costs only 1,000 francs.
\item We need to \_\_\_\_\_\_\_\_\_\_\_\_ before deciding which laptop to buy.
\end{enumerate}
\end{exercice}

\begin{exercice}[Entraînement 2 : Reformulation (français $\rightarrow$ anglais naturel)]
\textbf{Traduisez en utilisant les idiomes et formules de la leçon.}
\begin{enumerate}
\item « Ces chaussures sont une arnaque : elles coûtent une fortune pour cette qualité. »
\item « Je vais faire les courses au marché ; j'ai ma liste de courses. »
\item « Le pain est épuisé, il n'y en a plus. »
\item « Marché conclu ! Je vous l'achète à ce prix. »
\item « Ce sac est vraiment donné, il ne coûte presque rien. »
\end{enumerate}
\end{exercice}

\begin{experience}[Entraînement 3 : Jeu de rôle — « Au marché de Mfoundi »]
\textbf{Mise en situation : achat / vente (travail en binômes).}

\textbf{Rôles :}
\begin{itemize}
\item \textbf{Élève A (vendeur) :} vous tenez un stand de vivres frais (tomates, oignons, piments, bananes plantains). Vos prix sont fixes mais vous acceptez une petite réduction pour les gros achats (\eng{buy in bulk}). Vous êtes en rupture de gingembre (\eng{out of stock}).
\item \textbf{Élève B (client) :} vous devez acheter 10 kg de tomates, 5 kg d'oignons, du gingembre (mais il n'y en a plus) et 2 régimes de bananes plantains. Vous voulez négocier (\eng{ask for a discount}) et vous payez en espèces (\eng{cash}).
\end{itemize}
\textbf{Contraintes linguistiques :} utilisez \textbf{au moins 5 collocations ou idiomes} et \textbf{4 formules fonctionnelles} de la leçon.

\textbf{Déroulement :} 3 min de préparation $\rightarrow$ 3 min de jeu $\rightarrow$ 2 min de feedback (correction par les pairs puis par le professeur).
\end{experience}

\begin{corrige}[Exercices 1 et 2]
\textbf{Exercice 1 :}
\begin{enumerate}
\item \eng{risen} (ou \eng{gone up}) — \eng{prices have raised} est incorrect : \eng{raise} est transitif ; on dit \eng{prices have risen} ou \eng{prices have been raised}.
\item \eng{sold ... at a loss} — collocation \eng{sell at a loss} (vendre à perte).
\item \eng{go shopping} ; \eng{costs a fortune} — \eng{go shopping} pour le loisir au centre commercial ; idiome \eng{cost a fortune} (sujet singulier \eng{everything} $\rightarrow$ \eng{costs}).
\item \eng{value for money} — expression nominale invariable.
\item \eng{shop around} — verbe à particule intransitif : comparer les prix.
\end{enumerate}

\textbf{Exercice 2 (propositions de traduction) :}
\begin{enumerate}
\item \eng{These shoes are a rip-off: they cost a fortune for this quality.}
\item \eng{I'm going to do the shopping at the market; I have my shopping list.}
\item \eng{The bread is sold out; there is none left.} (ou \eng{There is no bread left.})
\item \eng{It's a deal! I'll buy it from you at that price.}
\item \eng{This bag is dirt cheap; it costs almost nothing.}
\end{enumerate}
\end{corrige}

\coursec{False friends and traps for francophones}

After mastering the collocations, idioms and functional expressions that make your interactions at the market or in a shop sound natural, we must now address a silent enemy of the francophone learner: \cle{faux amis} (false friends) and \cle{calques} (literal translations from French structure). These traps are particularly dangerous in a commercial context because they often lead to misunderstandings about prices, items, or intentions. A single wrong word can change “I bought a newspaper” into “I bought a shop”, or “How much does it cost?” into the ungrammatical “How much costs it?”. This section trains you to spot these errors, understand why they happen, and replace them with correct, idiomatic English.

\courssub{Shop, store vs. magazine : le piège de l'orthographe identique}

Commençons par le faux ami le plus classique du vocabulaire commercial. En français, \eng{magasin} désigne un lieu de vente. En anglais, \eng{magazine} existe mais signifie “revue, magazine (périodique)”. Le lieu de vente se dit \eng{shop} (boutique, petit commerce, terme britannique) ou \eng{store} (grand magasin, supermarché, terme plus américain).

\begin{textebox}[Dialogue at the newsagent's]
\textbf{Customer:} Excuse me, where is the nearest \eng{shop}? I need to buy a \eng{magazine} about football.\par
\textbf{Passer-by:} There is a \eng{newsagent's} on the corner. They sell newspapers and \eng{magazines}.\par
\textbf{Customer:} Thank you. Is it a big \eng{store}?\par
\textbf{Passer-by:} No, just a small \eng{shop}. But they have everything.
\end{textebox}

\begin{definition}[Glossaire]
\textbf{newsagent's} (nom) : marchand de journaux (boutique) ; \textbf{corner} (nom) : coin de rue ; \textbf{nearest} (adj., superlatif) : le plus proche.
\end{definition}

\begin{attention}[Magasin ≠ Magazine]
N'écrivez jamais \eng{“I went to the magazine to buy bread.”} (Je suis allé à la revue acheter du pain).
\eng{Shop} (UK) / \eng{Store} (US) = Magasin, boutique.\par
\eng{Magazine} = Revue, magazine (ex. \eng{Vogue}, \eng{Time}, \eng{Jeune Afrique}).
\par\medskip
\textbf{Exemple correct :} \eng{I bought a magazine at the shop.} (J'ai acheté une revue au magasin.)
\end{attention}

\begin{exemplebox}[Autres pièges du même type]
\begin{itemize}
    \item \eng{Library} ≠ \eng{Librairie}. \eng{Library} = Bibliothèque (prêt). \eng{Bookshop} (UK) / \eng{Bookstore} (US) = Librairie (vente).
    \item \eng{Supermarket} = Supermarché. \eng{Hypermarket} = Hypermarché.
    \item \eng{Market} = Marché (plein air, halles). Ne pas confondre avec \eng{Marketing} (mercatique).
\end{itemize}
\end{exemplebox}

\courssub{Money, change, currency : trois mots pour « monnaie »}

Le mot français “monnaie” est polysémique : il désigne l'argent en général, la petite monnaie rendue après un achat, ou le système monétaire d'un pays. L'anglais utilise trois mots distincts.

\begin{propriete}[Money / Change / Currency]
\begin{center}
\begin{tabularx}{\linewidth}{|X|X|X|}
\hline
\rowcolor{popblueL} \textbf{Mot anglais} & \textbf{Sens} & \textbf{Équivalent français contextuel} \\
\hline
\eng{Money} & Argent (général, billets, pièces, compte bancaire) & Argent, sous, fonds \\
\hline
\eng{Change} & Monnaie rendue (pièces/billets donnés en retour) & Monnaie (rendue), appoint \\
\hline
\eng{Currency} & Monnaie (système monétaire national, devise) & Monnaie, devise, unité monétaire \\
\hline
\end{tabularx}
\end{center}
\end{propriete}

\begin{exemplebox}[Exemples traduits]
\begin{itemize}
    \item \eng{“How much \textbf{money} do you have?”} = “Combien d'\textbf{argent} as-tu ?”
    \item \eng{“Keep the \textbf{change}.”} = “Garde la \textbf{monnaie}.” (Pourboire)
    \item \eng{“The local \textbf{currency} is the CFA franc.”} = “La \textbf{monnaie} / la \textbf{devise} locale est le franc CFA.”
    \item \eng{“She spent all her \textbf{money} at the market.”} = “Elle a dépensé tout son \textbf{argent} au marché.”
    \item \eng{“I don't have any small \textbf{change} for the taxi.”} = “Je n'ai pas de \textbf{petite monnaie} pour le taxi.”
\end{itemize}
\end{exemplebox}

\begin{attention}[Piège fréquent]
Ne dites pas \eng{“I have no money”} si vous voulez dire “Je n'ai pas de monnaie (pièces)”. Dites \eng{“I have no change”} ou \eng{“I have no coins”}.
\end{attention}

\courssub{To buy vs. to achieve : le piège phonétique}

C'est une erreur d'audition et d'orthographe très courante chez les francophones. \eng{To achieve} [\eng{ə-ˈtʃiːv}] sonne vaguement comme “achever” ou “acheter” à l'oreille non entraînée, mais son sens est radicalement différent.

\begin{definition}[Distinction stricte]
\begin{itemize}
    \item \cle{To buy} [baɪ] (irrégulier : buy, bought, bought) = Acheter. \eng{“I \textbf{bought} a kilo of plantains.”}
    \item \cle{To achieve} [əˈtʃiːv] (régulier) = Accomplir, réaliser, atteindre (un objectif). \eng{“She \textbf{achieved} her goal.”}
\end{itemize}
\end{definition}

\begin{attention}[Erreur type au Bac]
\eng{“I want to \textbf{achieve} a new phone.”} $\rightarrow$ INCORRECT (Je veux accomplir un nouveau téléphone).
\eng{“I want to \textbf{buy} a new phone.”} $\rightarrow$ CORRECT.
\end{attention}

\courssub{To spend vs. to dispense : le piège de la transparence morphologique}

\eng{To spend} ressemble à “dépenser”. \eng{To dispense} ressemble à “dispenser” (distribuer, exempter). Ce sont de faux amis partiels.

\begin{propriete}[Spend vs. Dispense]
\begin{center}
\begin{tabular}{|l|l|l|}
\hline
\rowcolor{popblueL} \textbf{Verbe} & \textbf{Sens principal} & \textbf{Exemple} \\
\hline
\eng{To spend} (spent, spent) & 1. Dépenser (argent). 2. Passer (temps). & \eng{He \textbf{spent} 5000 francs.} / \eng{We \textbf{spent} the afternoon shopping.} \\
\hline
\eng{To dispense} & Distribuer (médicaments, conseils), exempter. & \eng{The pharmacist \textbf{dispenses} medicine.} / \eng{Students are \textbf{dispensed from} sport.} \\
\hline
\end{tabular}
\end{center}
\end{propriete}

\begin{exemplebox}[Contextes commerciaux]
\begin{itemize}
    \item \eng{“Don't \textbf{spend} all your pocket money on sweets.”} (Ne dépense pas tout ton argent de poche en bonbons.)
    \item \eng{“The ATM \textbf{dispenses} cash.”} (Le distributeur \textbf{délivre/distribue} de l'argent.)
    \item \eng{“This card \textbf{dispenses with} the need for cash.”} (Cette carte \textbf{dispense de} / supprime le besoin d'espèces.) — \emph{registre formel}.
\end{itemize}
\end{exemplebox}

\courssub{To order vs. to command : le piège de l'autorité}

Au restaurant, au marché ou en ligne, on “commande” un plat ou un article. Le verbe anglais n'est pas \eng{to command} (donner un ordre militaire), mais \eng{to order}.

\begin{definition}[Order vs. Command]
\begin{itemize}
    \item \cle{To order} : Commander (repas, marchandises, en ligne). \eng{“Are you ready to \textbf{order}?”} (Prêt à commander ?). \eng{“I \textbf{ordered} a book on Amazon.”}
    \item \cle{To command} : Commander (donner un ordre), exiger, dominer. \eng{“The officer \textbf{commanded} the soldiers to stop.”}
\end{itemize}
\end{definition}

\begin{attention}[Situation au restaurant / marché]
Ne dites jamais au serveur : \eng{“I \textbf{command} the fish.”} (Il croira que vous êtes un officier donnant un ordre au poisson !).
Dites : \eng{“I would like to \textbf{order} the fish.”} ou \eng{“I'll \textbf{have} the fish.”}
\end{attention}

\courssub{Cheap, economic, economical : le triptyque de « économique »}

L'adjectif français “économique” couvre trois réalités distinctes en anglais. C'est une source majeure d'erreurs en expression écrite (composition, lettre formelle).

\begin{propriete}[Cheap / Economic / Economical]
\begin{center}
\begin{tabularx}{\linewidth}{|X|X|X|}
\hline
\rowcolor{popblueL} \textbf{Mot} & \textbf{Sens} & \textbf{Exemple traduit} \\
\hline
\eng{Cheap} [tʃiːp] & Bon marché, peu cher (prix bas). Peut avoir une connotation péjorative “de mauvaise qualité”. & \eng{“These tomatoes are \textbf{cheap} today.”} (Ces tomates sont \textbf{bon marché} aujourd'hui.) \\
\hline
\eng{Economic} [ˌiːkəˈnɒmɪk] & Relatif à l'économie (science, système, politique). Neutre. & \eng{“The \textbf{economic} growth of Cameroon.”} (La croissance \textbf{économique} du Cameroun.) \\
\hline
\eng{Economical} [ˌiːkəˈnɒmɪkl] & Qui fait économiser (argent, carburant, temps), rentable, économe. Positif. & \eng{“This car is very \textbf{economical} on fuel.”} (Cette voiture est très \textbf{économe} en carburant.) \\
\hline
\end{tabularx}
\end{center}
\end{propriete}

\begin{exemplebox}[Nuances dans un texte commercial]
\eng{“Buying in bulk at the \textbf{wholesale market} is more \textbf{economical} than buying retail. However, during the \textbf{economic} crisis, even \textbf{cheap} goods become unaffordable for many families.”}
\par\medskip
Traduction : « Acheter en gros au \textbf{marché de gros} est plus \textbf{économe/rentable} qu'acheter au détail. Cependant, pendant la crise \textbf{économique}, même les biens \textbf{bon marché} deviennent inabordables pour de nombreuses familles. »
\end{exemplebox}

\courssub{Calques structurels : questions sur le prix et expressions du marché}

Au-delà des mots isolés, la structure de la phrase piège le francophone. Deux calques majeurs reviennent systématiquement dans les copies d'élèves.

\courssub{« Combien coûte-t-il ? » $\neq$ \eng{“How much costs it?”}}

En français, le verbe suit le sujet dans la question formelle (inversion) ou on utilise “est-ce que”. En anglais, la question avec \eng{How much} (combien) exige l'auxiliaire \eng{do/does} (présent simple) ou \eng{did} (prétérit).

\begin{propriete}[Règle de la question \eng{How much}]
\begin{center}
\begin{tabular}{|l|l|}
\hline
\rowcolor{popblueL} \textbf{Structure} & \textbf{Exemple} \\
\hline
\eng{How much} + \textbf{do/does/did} + \textbf{sujet} + \textbf{verbe (base)} ? & \eng{How much \textbf{does} it \textbf{cost}?} \\
\hline
\eng{How much} + \textbf{is/are} + \textbf{sujet} ? (pour le prix affiché) & \eng{How much \textbf{is} this shirt?} \\
\hline
\end{tabular}
\end{center}
\end{propriete}

\begin{attention}[Erreurs interdites]
\begin{itemize}
    \item $\times$ \eng{How much costs it?} (Calque français “Combien coûte-t-il ?”)
    \item $\times$ \eng{How much it costs?} (Ordre affirmatif, pas interrogatif direct).
    \item $\checkmark$ \eng{How much \textbf{does} it \textbf{cost}?}
    \item $\checkmark$ \eng{How much \textbf{is} it?} (Court, très courant).
    \item $\checkmark$ \eng{How much \textbf{did} you \textbf{pay} for it?} (Prétérit).
\end{itemize}
\end{attention}

\courssub{« Faire le marché » $\neq$ \eng{“To make the market”}}

\eng{To make the market} n'existe pas en anglais standard (sauf sens très technique de “créer un marché financier”). Pour dire “faire les courses / aller au marché”, on utilise des verbes de déplacement ou d'action générale.

\begin{definition}[Aller au marché / Faire les courses]
\begin{itemize}
    \item \cle{To go to the market} / \cle{To go marketing} (surtout Afrique de l'Ouest / Cameroun) : Aller au marché.
    \item \cle{To do the shopping} / \cle{To go shopping} : Faire les courses (général, supermarché, centre-ville).
    \item \cle{To buy} / \cle{To get} : Acheter (action précise).
\end{itemize}
\end{definition}

\begin{exemplebox}[Dialogue camerounais typique]
\textbf{A:} \eng{“Where are you going this morning?”}\par
\textbf{B:} \eng{“I'm \textbf{going to the market} / I'm \textbf{going marketing}. I need to \textbf{buy} plantains and groundnuts.”}\par
\textbf{A:} \eng{“Okay. \textbf{Do the shopping} for the house too, please.”}
\end{exemplebox}

\begin{methode}[Comment ne plus calquer le français sur l'anglais]
\begin{enumerate}
    \item \textbf{Identifiez l'intention} en français (ex. : “Je veux dire ‘faire le marché’”).
    \item \textbf{Ne traduisez pas mot à mot} (\eng{make} + \eng{the} + \eng{market}).
    \item \textbf{Cherchez le verbe/nom anglais dans un dictionnaire anglais-anglais} (Oxford, Cambridge, Longman) ou une corpus de phrases (Reverso Context, Linguee) en tapant l'idée clé : \eng{“go to market”}, \eng{“buy food”}.
    \item \textbf{Vérifiez la collocation} : \eng{go to the market} (oui), \eng{make the market} (non).
    \item \textbf{Mémorisez le bloc (chunk)} : \eng{“go marketing”}, \eng{“do the shopping”}, \eng{“how much does it cost?”}.
\end{enumerate}
\end{methode}

\begin{aretenir}[Les 5 pièges à retenir pour le Bac]
\begin{itemize}
    \item \textbf{Magazine} = Revue ; \textbf{Shop/Store} = Magasin.
    \item \textbf{Money} = Argent ; \textbf{Change} = Monnaie rendue ; \textbf{Currency} = Devise.
    \item \textbf{Buy} (bought) = Acheter ; \textbf{Achieve} = Accomplir.
    \item \textbf{Order} = Commander (repas) ; \textbf{Command} = Ordonner.
    \item \textbf{Cheap} = Pas cher ; \textbf{Economic} = Économique (science) ; \textbf{Economical} = Économe (rentable).
    \item \textbf{How much does it cost?} (pas \eng{costs it}) ; \textbf{Go to the market} (pas \eng{make the market}).
\end{itemize}
\end{aretenir}

\begin{popfigure}
\begin{tikzpicture}[
    node distance=0.8cm and 1.2cm,
    every node/.style={font=\small},
    fr/.style={rectangle, draw=popblue, thick, fill=popblueL, rounded corners, minimum width=3.5cm, align=center},
    fa/.style={rectangle, draw=poporange, thick, fill=poporangeL, rounded corners, minimum width=3.5cm, align=center},
    co/.style={rectangle, draw=popgreen, thick, fill=popgreenL, rounded corners, minimum width=3.5cm, align=center},
    arr/.style={-Stealth, thick, popdark}
]
% En-têtes
\node[fr] (h1) at (0,0) {\textbf{Français (Idée)}};
\node[fa] (h2) at (5,0) {\textbf{Faux ami / Calque incorrect}};
\node[co] (h3) at (10,0) {\textbf{Mot / Structure correcte}};

% Ligne 1
\node[fr] (f1) at (0,-1.2) {Magasin (lieu)};
\node[fa] (a1) at (5,-1.2) {\eng{magazine}};
\node[co] (c1) at (10,-1.2) {\eng{shop / store}};

% Ligne 2
\node[fr] (f2) at (0,-2.4) {Monnaie (rendue)};
\node[fa] (a2) at (5,-2.4) {\eng{money}};
\node[co] (c2) at (10,-2.4) {\eng{change}};

% Ligne 3
\node[fr] (f3) at (0,-3.6) {Monnaie (système)};
\node[fa] (a3) at (5,-3.6) {\eng{money}};
\node[co] (c3) at (10,-3.6) {\eng{currency}};

% Ligne 4
\node[fr] (f4) at (0,-4.8) {Acheter};
\node[fa] (a4) at (5,-4.8) {\eng{achieve}};
\node[co] (c4) at (10,-4.8) {\eng{buy}};

% Ligne 5
\node[fr] (f5) at (0,-6.0) {Dépenser (argent)};
\node[fa] (a5) at (5,-6.0) {\eng{dispense}};
\node[co] (c5) at (10,-6.0) {\eng{spend}};

% Ligne 6
\node[fr] (f6) at (0,-7.2) {Commander (repas)};
\node[fa] (a6) at (5,-7.2) {\eng{command}};
\node[co] (c6) at (10,-7.2) {\eng{order}};

% Ligne 7
\node[fr] (f7) at (0,-8.4) {Bon marché (prix)};
\node[fa] (a7) at (5,-8.4) {\eng{economic}};
\node[co] (c7) at (10,-8.4) {\eng{cheap}};

% Ligne 8
\node[fr] (f8) at (0,-9.6) {Économe (rentable)};
\node[fa] (a8) at (5,-9.6) {\eng{economic}};
\node[co] (c8) at (10,-9.6) {\eng{economical}};

% Ligne 9
\node[fr] (f9) at (0,-10.8) {Combien coûte-t-il ?};
\node[fa] (a9) at (5,-10.8) {\eng{How much costs it?}};
\node[co] (c9) at (10,-10.8) {\eng{How much does it cost?}};

% Ligne 10
\node[fr] (f10) at (0,-12.0) {Faire le marché};
\node[fa] (a10) at (5,-12.0) {\eng{Make the market}};
\node[co] (c10) at (10,-12.0) {\eng{Go to the market / Do the shopping}};

% Flèches de correction
\draw[arr, poporange] (a1) -- (c1) node[midway, above, sloped, font=\tiny, text=popdark] {CORRECTION};
\draw[arr, poporange] (a2) -- (c2) node[midway, above, sloped, font=\tiny, text=popdark] {CORRECTION};
\draw[arr, poporange] (a3) -- (c3) node[midway, above, sloped, font=\tiny, text=popdark] {CORRECTION};
\draw[arr, poporange] (a4) -- (c4) node[midway, above, sloped, font=\tiny, text=popdark] {CORRECTION};
\draw[arr, poporange] (a5) -- (c5) node[midway, above, sloped, font=\tiny, text=popdark] {CORRECTION};
\draw[arr, poporange] (a6) -- (c6) node[midway, above, sloped, font=\tiny, text=popdark] {CORRECTION};
\draw[arr, poporange] (a7) -- (c7) node[midway, above, sloped, font=\tiny, text=popdark] {CORRECTION};
\draw[arr, poporange] (a8) -- (c8) node[midway, above, sloped, font=\tiny, text=popdark] {CORRECTION};
\draw[arr, poporange] (a9) -- (c9) node[midway, above, sloped, font=\tiny, text=popdark] {CORRECTION};
\draw[arr, poporange] (a10) -- (c10) node[midway, above, sloped, font=\tiny, text=popdark] {CORRECTION};
\end{tikzpicture}
\legende{Tableau récapitulatif : les 10 pièges majeurs du vocabulaire « achat/vente » pour francophones.}
\end{popfigure}

\begin{savaistu}[Actually : le faux ami qui coûte des points au Bac]
Le mot \eng{actually} [ˈæktʃuəli] ne signifie \textbf{jamais} “actuellement” (ce qui se dit \eng{currently}, \eng{at the moment}, \eng{nowadays}).
\eng{Actually} veut dire \textbf{“en fait”, “en réalité”, “pour dire la vérité”}. Il sert à corriger une idée reçue ou à introduire une surprise.
\par\medskip
\eng{“People think imported rice is better. \textbf{Actually}, local rice is more nutritious.”}
(Les gens pensent que le riz importé est meilleur. \textbf{En fait}, le riz local est plus nutritif.)
\par\medskip
\textbf{Prononciation :} « èk-chou-eu-li » (accent sur \textbf{chou}). Ne pas dire « ac-tuel-le-ment ».
\end{savaistu}

\begin{exoresolu}[Correction d'erreurs de faux amis et calques]
Corrigez les erreurs dans les phrases suivantes (vocabulaire, grammaire, structure). Justifiez chaque correction.
\begin{enumerate}
    \item I want to \textbf{achieve} a new pair of shoes at the \textbf{magazine}.
    \item How much \textbf{costs} this bag of rice? I have no \textbf{money} (meaning small coins).
    \item The \textbf{economic} situation is hard, so I look for \textbf{economical} clothes.
    \item The waiter asked if I was ready to \textbf{command}.
    \item She \textbf{dispended} all her salary at the market yesterday.
    \item We go to \textbf{make the market} every Saturday morning.
\end{enumerate}

\tcblower
\textbf{Solution détaillée :}
\begin{enumerate}
    \item \eng{I want to \textbf{buy} a new pair of shoes at the \textbf{shop/store}.}
    \textbf{Justification :} \eng{Achieve} = accomplir ; \eng{buy} = acheter. \eng{Magazine} = revue ; \eng{shop/store} = magasin.
    \item \eng{How much \textbf{does} this bag of rice \textbf{cost}? I have no \textbf{change}.}
    \textbf{Justification :} Question \eng{How much} $\rightarrow$ auxiliaire \eng{does} + base verbale \eng{cost}. \eng{Money} = argent général ; \eng{change} = monnaie rendue/pièces.
    \item \eng{The \textbf{economic} situation is hard, so I look for \textbf{cheap} clothes.}
    \textbf{Justification :} \eng{Economic} = relatif à l'économie (correct ici pour “situation”). \eng{Economical} = économe/rentable (qualité d'un objet). Pour “pas cher”, on utilise \eng{cheap} (ou \eng{affordable} / \eng{inexpensive}).
    \item \eng{The waiter asked if I was ready to \textbf{order}.}
    \textbf{Justification :} \eng{Command} = donner un ordre. Au restaurant, on \eng{order}.
    \item \eng{She \textbf{spent} all her salary at the market yesterday.}
    \textbf{Justification :} \eng{Dispense} = distribuer/exempter. \eng{Spend} (spent, spent) = dépenser.
    \item \eng{We \textbf{go to the market} / \textbf{do the shopping} every Saturday morning.}
    \textbf{Justification :} \eng{Make the market} n'existe pas. \eng{Go to the market} (aller au marché) ou \eng{do the shopping} (faire les courses).
\end{enumerate}
\end{exoresolu}

\begin{exercice}[Entraînement : Chasse aux faux amis]
\textbf{A. Traduisez en anglais en évitant tous les pièges de cette leçon.}
\begin{enumerate}
    \item J'ai acheté une revue au supermarché.
    \item Gardez la monnaie, s'il vous plaît.
    \item Combien coûte ce kilo de tomates ? (Deux versions possibles).
    \item Ce téléphone n'est pas cher, il est économe en batterie.
    \item Nous allons au marché samedi pour faire les courses.
    \item Le serveur a demandé si nous voulions commander.
    \item La situation économique du pays s'améliore.
    \item Elle a dépensé tout son argent en vêtements.
\end{enumerate}

\textbf{B. Cochez la phrase correcte dans chaque paire.}
\begin{enumerate}
    \item a) How much costs the bread? \quad b) How much does the bread cost?
    \item a) I bought a magazine at the shop. \quad b) I bought a shop at the magazine.
    \item a) This car is very economic. \quad b) This car is very economical. (Contexte : elle consomme peu).
    \item a) I have no money for the taxi (meaning coins). \quad b) I have no change for the taxi.
    \item a) Are you ready to command? \quad b) Are you ready to order?
\end{enumerate}
\end{exercice}


\begin{corrige}[Exercice n]
\textbf{A. Traductions}
\begin{enumerate}
    \item \eng{I bought a \textbf{magazine} at the \textbf{supermarket}.}
    \item \eng{\textbf{Keep the change}, please.}
    \item \eng{\textbf{How much does this kilo of tomatoes cost?} / \textbf{How much is this kilo of tomatoes?}}
    \item \eng{This phone is \textbf{cheap} / \textbf{inexpensive}; it is \textbf{economical} on battery.}
    \item \eng{We are \textbf{going to the market} / \textbf{going marketing} on Saturday to \textbf{do the shopping}.}
    \item \eng{The waiter asked if we wanted to \textbf{order}.}
    \item \eng{The \textbf{economic} situation of the country is improving.}
    \item \eng{She \textbf{spent} all her \textbf{money} on clothes.}
\end{enumerate}

\textbf{B. Choix correct}
\begin{enumerate}
    \item \textbf{b)} (\eng{How much} + auxiliaire \eng{does}).
    \item \textbf{a)} (\eng{magazine} = revue, \eng{shop} = magasin).
    \item \textbf{b)} (\eng{economical} = économe en carburant/batterie).
    \item \textbf{b)} (\eng{change} = petite monnaie/pièces).
    \item \textbf{b)} (\eng{order} = commander un repas).
\end{enumerate}
\end{corrige}

\coursec{Word families, pronunciation and spelling}

\courssub{Transition from false friends to word building}

After identifying the \cle{false friends} and translation traps that often mislead francophone learners in the previous section, we now turn to a more constructive skill: \textbf{building vocabulary systematically}. Mastering a lexical field does not mean memorising isolated words; it means understanding how words relate to one another through \cle{word families} (derivation), how they sound (\cle{pronunciation}), and how they are written (\cle{spelling}). This section will equip you with the tools to expand your commercial vocabulary from a few key roots.

\begin{textebox}{A market scene in Buea Town}
\eng{Mrs Epote arrives at the Buea Main Market early on a Saturday morning. She is a regular \textbf{customer} at Mama Sarah's foodstuff stall. Today, she wants to \textbf{buy} a basin of \textbf{plantains} and some \textbf{groundnuts}. She approaches the \textbf{seller} and asks, “How much \textbf{do} you \textbf{sell} the plantains \textbf{for}?” Mama Sarah replies, “Five thousand francs per basin, but since you are a loyal \textbf{buyer}, I can give you a \textbf{discount}. They are \textbf{on sale} today.” Mrs Epote \textbf{bargains} a little, \textbf{pays} the agreed \textbf{price}, and collects her \textbf{receipt}. She leaves \textbf{satisfied} with her \textbf{purchase}, happy that she did not \textbf{spend} too much.}
\medskip

\noindent\textbf{Glossaire :}
\begin{itemize}
    \item \eng{stall} (n) : étal, stand
    \item \eng{plantains} (n pl) : plantains (bananes plantains)
    \item \eng{groundnuts} (n pl) : arachides
    \item \eng{loyal} (adj) : fidèle
    \item \eng{discount} (n) : remise, réduction
    \item \eng{on sale} : en promotion, en solde
    \item \eng{bargains} (v) : marchande
    \item \eng{receipt} (n) : reçu, ticket de caisse
    \item \eng{purchase} (n) : achat
\end{itemize}
\end{textebox}

\courssub{The \eng{BUY} family: irregular verb and derivatives}

Observe the text above: \eng{buy} (verb), \eng{buyer} (noun), \eng{buying} (noun/gerund), \eng{bought} (past forms). English builds families by adding suffixes to a root verb.

\begin{propriete}[Famille lexicale de BUY]
\begin{center}
\begin{tabularx}{\linewidth}{|X|X|X|X|}
\hline
\rowcolor{popblueL}
\textbf{Forme} & \textbf{Mot} & \textbf{Nature} & \textbf{Traduction / Note} \\
\hline
Verbe de base & \eng{buy} & v. & acheter \\
\hline
Prétérit / Participe passé & \eng{bought} & v. & \textbf{irrégulier} (ne jamais écrire *\eng{buyed}) \\
\hline
Agent (personne) & \eng{buyer} & n. & acheteur, client (celui qui achète) \\
\hline
Action / Nom verbal & \eng{buying} & n. & achat (action d'acheter) \\
\hline
Nom d'action (formel) & \eng{purchase} & n. & achat (registre soutenu) \\
\hline
Expression utile & \eng{a buying spree} & loc. & une frénésie d'achats \\
\hline
\end{tabularx}
\end{center}

\medskip
\noindent \textbf{Règle de dérivation :} Verbe + \cle{-er} $\rightarrow$ personne (agent). Verbe + \cle{-ing} $\rightarrow$ action, nom verbal.
\end{propriete}

\begin{exemplebox}[Exemples traduits — Famille BUY]
\eng{My father is the main \textbf{buyer} for the household.} « Mon père est le principal acheteur du ménage. »

\eng{\textbf{Buying} in bulk saves money.} « Acheter en gros fait économiser de l'argent. »

\eng{She \textbf{bought} a new wrapper yesterday.} « Elle a acheté un nouveau pagne hier. »

\eng{Have you \textbf{bought} the tickets for the match?} « As-tu acheté les billets pour le match ? »
\end{exemplebox}

\courssub{The \eng{SELL} family: verb, noun and agent}

The root \eng{sell} generates the agent \eng{seller} and the noun \eng{sale} (note the vowel change). Note the spelling of the past forms.

\begin{propriete}[Famille lexicale de SELL]
\begin{center}
\begin{tabularx}{\linewidth}{|X|X|X|X|}
\hline
\rowcolor{popblueL}
\textbf{Forme} & \textbf{Mot} & \textbf{Nature} & \textbf{Traduction / Note} \\
\hline
Verbe de base & \eng{sell} & v. & vendre \\
\hline
Prétérit / Participe passé & \eng{sold} & v. & \textbf{irrégulier} (ne jamais écrire *\eng{selled}) \\
\hline
Agent (personne) & \eng{seller} & n. & vendeur, vendeuse \\
\hline
Nom d'action / Événement & \eng{sale} & n. & vente ; solde (\eng{on sale} = en promotion) \\
\hline
Expression figée & \eng{for sale} & loc. & à vendre (sur un panneau) \\
\hline
Nom composé & \eng{goods on sale} & n. & articles en solde \\
\hline
\end{tabularx}
\end{center}

\medskip
\noindent \textbf{Attention orthographe :} \eng{sell} (verbe, double \eng{l}, voyelle courte « è ») et \eng{sale} (nom, un seul \eng{l}, voyelle longue « é-i »). Le double \eng{ll} de \eng{seller} protège la voyelle courte devant le suffixe \eng{-er}.
\end{propriete}

\begin{popfigure}
\begin{tikzpicture}[
    mindmap,
    grow cyclic,
    every node/.style={concept, align=center},
    root concept/.append style={concept color=popblue, minimum size=2.4cm, text=white, font=\bfseries\large},
    level 1 concept/.append style={concept color=popblue!70, sibling angle=90, minimum size=1.9cm, text=white, font=\bfseries\small},
    level 2 concept/.append style={concept color=popblue!40, sibling angle=60, minimum size=1.5cm, text=white, font=\small}
]
\node[root concept, align=center]{SELL\\ \footnotesize « sèl »}
    child[clockwise from=60] { node[level 1 concept, align=center]{SOLD\\ \footnotesize « solde »\\ \footnotesize prétérit / p.p.} }
    child[clockwise from=150] { node[level 1 concept, align=center]{SELLER\\ \footnotesize « sè-leur »\\ \footnotesize personne} }
    child[clockwise from=240] { node[level 1 concept, align=center]{SALE\\ \footnotesize « sé-il »\\ \footnotesize nom}
        child { node[level 2 concept, align=center]{FOR SALE\\ \footnotesize à vendre} }
        child { node[level 2 concept, align=center]{ON SALE\\ \footnotesize en solde} }
    }
    child[clockwise from=330] { node[level 1 concept, align=center]{SELLING\\ \footnotesize « sè-ling »\\ \footnotesize action} };
\end{tikzpicture}
\legende{Carte mentale de la famille SELL (prononciation approximée en lettres françaises)}
\end{popfigure}

\begin{exemplebox}[Exemples traduits — Famille SELL]
\eng{This car is \textbf{for sale}.} « Cette voiture est à vendre. »

\eng{The \textbf{seller} agreed to lower the price.} « Le vendeur a accepté de baisser le prix. »

\eng{They \textbf{sold} all their cocoa beans.} « Ils ont vendu toutes leurs fèves de cacao. »

\eng{I bought these mangoes \textbf{on sale}.} « J'ai acheté ces mangues en promotion. »

\eng{The \textbf{sale} of the land was finalised.} « La vente du terrain a été finalisée. »
\end{exemplebox}

\courssub{The \eng{PAY} family: payment, payable and the irregular past}

\begin{propriete}[Famille lexicale de PAY]
\begin{center}
\begin{tabularx}{\linewidth}{|X|X|X|X|}
\hline
\rowcolor{popblueL}
\textbf{Forme} & \textbf{Mot} & \textbf{Nature} & \textbf{Traduction / Note} \\
\hline
Verbe de base & \eng{pay} & v. & payer \\
\hline
Prétérit / Participe passé & \eng{paid} & v. & \textbf{irrégulier} : \eng{y} $\rightarrow$ \eng{i} + \eng{d} \\
\hline
Nom d'action & \eng{payment} & n. & paiement, versement \\
\hline
Adjectif & \eng{payable} & adj & payable (à payer) \\
\hline
Phrasal verb & \eng{pay for} & v. & payer (quelque chose) \\
\hline
Expression utile & \eng{pay in cash / by mobile money} & loc. & payer en espèces / par Mobile Money \\
\hline
\end{tabularx}
\end{center}
\end{propriete}

\begin{attention}[Piège orthographique : PAY $\rightarrow$ PAID]
Le verbe \eng{pay} se termine par \eng{-ay}. Au prétérit et au participe passé, le \eng{y} devient \eng{i} devant \eng{d} : \eng{paid} (prononcé « pé-id », comme \eng{said} avec un \eng{p}).
\medskip

\noindent *\eng{Payed} est une erreur très fréquente (calque sur \eng{played}, \eng{stayed}). La forme \eng{payed} n'existe que dans un sens technique très rare (marine : laisser filer un câble) ; elle ne s'emploie jamais pour l'argent.
\medskip

\noindent \textbf{Correct :} \eng{I \textbf{paid} 5,000 francs.}

\noindent \textbf{Incorrect :} *\eng{I payed 5,000 francs.}
\medskip

\noindent \textbf{Autre piège :} on paie \textbf{pour} quelque chose : \eng{pay \textbf{for} the goods} (« payer les marchandises »). Ne dites pas *\eng{pay the goods} quand vous voulez dire « payer les articles » ; \eng{pay the goods} signifierait « payer les marchandises elles-mêmes », ce qui est absurde. En revanche on dit bien \eng{pay the bill} (payer la facture) et \eng{pay the seller} (payer le vendeur).
\end{attention}

\begin{exemplebox}[Exemples traduits — Famille PAY]
\eng{You can \textbf{pay} by mobile money.} « Vous pouvez payer par Mobile Money. »

\eng{The \textbf{payment} is due tomorrow.} « Le paiement est dû pour demain. »

\eng{Is the bill \textbf{payable} at the counter?} « La facture est-elle payable au guichet ? »

\eng{He \textbf{paid for} the drinks.} « Il a payé les boissons. »
\end{exemplebox}

\courssub{The \eng{CHEAP} / \eng{EXPENSIVE} families: adjectives, comparatives, adverbs}

These are the core adjectives for evaluating prices. Their morphology follows standard rules, but pronunciation holds traps.

\begin{propriete}[Familles CHEAP et EXPENSIVE]
\begin{center}
\begin{tabularx}{\linewidth}{|X|X|X|X|}
\hline
\rowcolor{popblueL}
\textbf{Adjectif} & \textbf{Comparatif} & \textbf{Superlatif} & \textbf{Adverbe} \\
\hline
\eng{cheap} (« tchipe ») & \eng{cheaper} & \eng{the cheapest} & \eng{cheaply} (« tchip-li ») \\
\hline
\eng{expensive} (« iks-pèn-siv ») & \eng{more expensive} & \eng{the most expensive} & \eng{expensively} \\
\hline
\end{tabularx}
\end{center}

\medskip
\noindent \textbf{Rappel de grammaire :} adjectif court (une syllabe) $\rightarrow$ \eng{-er / -est} ; adjectif long (trois syllabes et plus) $\rightarrow$ \eng{more / the most}. Ne jamais mélanger : *\eng{more cheaper} est incorrect.
\end{propriete}

\begin{attention}[Prononciation : CHEAP vs SHEEP — le piège « tch » / « ch »]
C'est l'erreur n°1 des francophones.
\begin{itemize}
    \item \eng{CHEAP} « \textbf{tch}ipe » : commence par le son « \textbf{tch} » (comme dans « \textbf{tch}èque »). La langue touche le palais puis se relâche.
    \item \eng{SHEEP} « \textbf{ch}ipe » : commence par le son « \textbf{ch} » (comme dans « \textbf{ch}aussure »). La langue ne touche pas le palais ; l'air passe en continu.
\end{itemize}
\medskip

\noindent \textbf{Astuce :} \eng{CH} en anglais se prononce souvent « tch » (\eng{church}, \eng{cheap}, \eng{choose}) ; \eng{SH} se prononce « ch » comme en français (\eng{shop}, \eng{sheep}, \eng{shoe}).
\medskip

\noindent Phrase d'entraînement volontairement absurde : \eng{Don't buy cheap shoes, buy sheep!} « N'achetez pas des chaussures bon marché, achetez des moutons ! »
\end{attention}

\begin{exemplebox}[Exemples traduits — Cheap / Expensive]
\eng{These tomatoes are \textbf{cheap} today.} « Ces tomates sont bon marché aujourd'hui. »

\eng{This shop is \textbf{cheaper} than the supermarket.} « Ce magasin est moins cher que le supermarché. »

\eng{It is \textbf{the cheapest} market in Yaoundé.} « C'est le marché le moins cher de Yaoundé. »

\eng{Imported rice is \textbf{expensive}.} « Le riz importé est cher. »

\eng{She dresses \textbf{expensively}.} « Elle s'habille de façon coûteuse. »
\end{exemplebox}

\courssub{The \eng{TRADE} family: noun/verb homograph and agent}

\begin{propriete}[Famille lexicale de TRADE]
\begin{center}
\begin{tabularx}{\linewidth}{|X|X|X|X|}
\hline
\rowcolor{popblueL}
\textbf{Forme} & \textbf{Mot} & \textbf{Nature} & \textbf{Traduction / Note} \\
\hline
Verbe / Nom (homographe) & \eng{trade} & v. / n. & commercer, échanger / commerce, échange \\
\hline
Agent & \eng{trader} & n. & commerçant, négociant \\
\hline
Action / Secteur & \eng{trading} & n. & commerce, négoce \\
\hline
Composé usuel & \eng{fair trade} & n. & commerce équitable \\
\hline
Composé usuel & \eng{trade fair} & n. & foire commerciale \\
\hline
\end{tabularx}
\end{center}

\medskip
\noindent \textbf{Remarque :} \eng{trade} est un verbe régulier : \eng{trade – traded – traded}. Le \eng{-e} final disparaît devant \eng{-er} et \eng{-ing} : \eng{trader}, \eng{trading} (jamais *\eng{tradeer}, *\eng{tradeing}).
\end{propriete}

\begin{exemplebox}[Exemples traduits — Famille TRADE]
\eng{They \textbf{trade} in cocoa and coffee.} « Ils font le commerce du cacao et du café. »

\eng{International \textbf{trade} boosts the economy.} « Le commerce international stimule l'économie. »

\eng{My uncle is a \textbf{trader} at the Central Market.} « Mon oncle est commerçant au Marché Central. »

\eng{\textbf{Trading} hours are 8 a.m. to 6 p.m.} « Les heures d'ouverture sont de 8 h à 18 h. »
\end{exemplebox}

\courssub{Critical pronunciation points for commercial vocabulary}

Three high-frequency words have a counter-intuitive pronunciation for francophones.

\begin{propriete}[Prononciation clé : bargain, receipt, customer]
\begin{center}
\begin{tabularx}{\linewidth}{|X|X|X|}
\hline
\rowcolor{popblueL}
\textbf{Mot} & \textbf{Approximation française} & \textbf{Piège à éviter} \\
\hline
\eng{bargain} & « \textbf{ba}-guine » (accent sur la 1\iere{} syllabe) & Ne pas dire « bar-gaine » en deux syllabes dures ; le \eng{-ain} final est faible, comme dans \eng{captain}. \\
\hline
\eng{receipt} & « ri-\textbf{site} » & Le \textbf{p est muet} ! Le \eng{-ei-} se prononce « i » long (comme dans \eng{receive}). \\
\hline
\eng{customer} & « \textbf{keu}s-to-meur » & Le \eng{cu-} se prononce « keu » bref (pas « kiu » ni « kou ») ; accent sur la 1\iere{} syllabe. \\
\hline
\end{tabularx}
\end{center}
\end{propriete}

\begin{savaistu}[Pourquoi \eng{receipt} a-t-il un \eng{p} muet ?]
Le mot vient du latin \emph{recepta} (« reçue »), passé par l'ancien français \emph{receite}. Au XVI\ieme{} siècle, les érudits anglais ont réintroduit le \eng{p} latin pour marquer l'étymologie — comme dans \eng{debt} (latin \emph{debitum}) et \eng{doubt} (latin \emph{dubitare}) — mais la prononciation n'a jamais suivi l'orthographe savante.
\medskip

\noindent \textbf{Règle mnémotechnique pour l'orthographe :} “\eng{i before e, except after c}” (« i avant e, sauf après c ») $\rightarrow$ \eng{re\textbf{cei}pt}, \eng{re\textbf{cei}ve}, \eng{\textbf{cei}ling}. Cette règle aide à écrire, pas à prononcer le \eng{p} !
\end{savaistu}

\courssub{Spelling focus: double consonants and single letters}

English spelling often doubles the final consonant after a short stressed vowel, to protect the short vowel sound.

\begin{propriete}[Règles d'orthographe : doublement et lettre unique]
\begin{center}
\begin{tabularx}{\linewidth}{|X|X|X|}
\hline
\rowcolor{popblueL}
\textbf{Mot} & \textbf{Règle / Explication} & \textbf{Erreur fréquente} \\
\hline
\eng{shopping} & \eng{shop} : une syllabe, voyelle courte + une consonne finale $\rightarrow$ on double le \eng{p} devant \eng{-ing}. & *\eng{shoping} \\
\hline
\eng{seller} & \eng{sell} : voyelle courte + \eng{ll} déjà double $\rightarrow$ on garde \eng{ll} devant \eng{-er}. & *\eng{seler} \\
\hline
\eng{sale} & Nom à voyelle longue « é-i » $\rightarrow$ un seul \eng{l}. & *\eng{salle} (mot français !), confusion avec \eng{sell} \\
\hline
\eng{payment} & \eng{pay} se termine par \eng{-y} $\rightarrow$ pas de doublement devant \eng{-ment}. & *\eng{payyment} \\
\hline
\eng{bargaining} & \eng{bargain} : l'accent tombe sur la 1\iere{} syllabe $\rightarrow$ pas de doublement du \eng{n}. & *\eng{bargainning} \\
\hline
\end{tabularx}
\end{center}
\end{propriete}

\begin{aretenir}[Règle générale du doublement (« règle 1-1-1 »)]
Pour un verbe d'\textbf{une} syllabe, se terminant par \textbf{une} voyelle courte + \textbf{une} consonne : on double la consonne finale devant un suffixe commençant par une voyelle (\eng{-ing}, \eng{-er}, \eng{-ed}, \eng{-est}).

\eng{shop} $\rightarrow$ \eng{shopping}, \eng{shopper} \quad | \quad \eng{stop} $\rightarrow$ \eng{stopping}, \eng{stopped} \quad | \quad \eng{big} $\rightarrow$ \eng{bigger}, \eng{the biggest}.
\medskip

\noindent \textbf{Exceptions :} pas de doublement après \eng{-w}, \eng{-x}, \eng{-y} (\eng{play} $\rightarrow$ \eng{playing} ; \eng{mix} $\rightarrow$ \eng{mixing}), ni quand l'accent ne porte pas sur la dernière syllabe (\eng{bargain} $\rightarrow$ \eng{bargaining}).
\end{aretenir}

\courssub{Method: how to memorise a word family effectively}

\begin{methode}[Mémoriser une famille de mots en 4 étapes]
\begin{enumerate}
    \item \textbf{Identifiez le verbe racine} (souvent le mot le plus court : \eng{buy}, \eng{sell}, \eng{pay}, \eng{trade}). Apprenez sa prononciation et ses formes irrégulières (prétérit et participe passé).
    \item \textbf{Construisez le nom d'agent} : ajoutez \eng{-er} $\rightarrow$ \eng{buyer}, \eng{seller}, \eng{trader}. Vérifiez le doublement de la consonne (\eng{seller}) et la chute du \eng{-e} final (\eng{trader}).
    \item \textbf{Construisez le nom d'action} : ajoutez \eng{-ing} (\eng{buying}, \eng{selling}) ou apprenez le nom spécifique (\eng{sale}, \eng{payment}, \eng{purchase}). Notez la prononciation (\eng{sale} « sé-il » $\neq$ \eng{sell} « sèl »).
    \item \textbf{Cherchez l'adjectif et l'adverbe} : \eng{payable}, \eng{cheap / cheaply}, \eng{expensive / expensively}. Testez-les dans une phrase type : \eng{The \textbf{cheap} goods sell \textbf{cheaply}.} « Les articles bon marché se vendent à bas prix. »
\end{enumerate}
\end{methode}

\courssub{Consolidation: the irregular verbs of commerce}

\begin{propriete}[Les 5 verbes indispensables du commerce]
\begin{center}
\begin{tabular}{|l|c|c|c|}
\hline
\rowcolor{popblueL}
\textbf{Infinitif} & \textbf{Prétérit} & \textbf{Participe passé} & \textbf{Traduction} \\
\hline
\eng{buy} & \eng{bought} & \eng{bought} & acheter \\
\hline
\eng{sell} & \eng{sold} & \eng{sold} & vendre \\
\hline
\eng{pay} & \eng{paid} & \eng{paid} & payer \\
\hline
\eng{spend} & \eng{spent} & \eng{spent} & dépenser (argent, temps) \\
\hline
\eng{cost} & \eng{cost} & \eng{cost} & coûter (inchangé) \\
\hline
\end{tabular}
\end{center}

\medskip
\noindent \textbf{Remarque :} \eng{spend} suit le même schéma que \eng{sell} (\eng{e} $\rightarrow$ \eng{o} ou \eng{e} brève dans le passé). \eng{Cost} est entièrement invariable : \eng{cost – cost – cost}.
\end{propriete}

\begin{attention}[Faux ami de structure : \eng{spend} et \eng{cost}]
\begin{itemize}
    \item Le sujet de \eng{spend} est une \textbf{personne} : \eng{I \textbf{spent} 2,000 francs on yams.} « J'ai dépensé 2 000 francs en ignames. » Structure : \eng{spend + somme + \textbf{on} + nom}.
    \item Le sujet de \eng{cost} est une \textbf{chose} : \eng{The yams \textbf{cost} 2,000 francs.} « Les ignames coûtent 2 000 francs. »
    \item Ne dites jamais *\eng{I cost 2,000 francs} ni *\eng{The yams spent 2,000 francs}.
\end{itemize}
\end{attention}

\begin{exoresolu}[Compléter le tableau et la phrase]
\textbf{A.} Complétez les formes manquantes.
\begin{center}
\begin{tabularx}{\linewidth}{|X|X|X|X|}
\hline
\rowcolor{popblueL}
\textbf{Verbe (base)} & \textbf{Prétérit} & \textbf{Participe passé} & \textbf{Nom clé} \\
\hline
\eng{buy} & \ldots & \ldots & \eng{buyer} \\
\hline
\eng{sell} & \ldots & \ldots & \eng{sale} \\
\hline
\eng{pay} & \ldots & \ldots & \eng{payment} \\
\hline
\eng{spend} & \ldots & \ldots & \eng{spending} \\
\hline
\end{tabularx}
\end{center}

\textbf{B.} Complétez la phrase avec la forme correcte du verbe entre parenthèses, puis traduisez.

\eng{Last week, Mrs Epote \_\_\_\_\_ (buy) plantains and \_\_\_\_\_ (pay) 3,000 francs. She \_\_\_\_\_ (spend) less than usual because they were \_\_\_\_\_ (cheap).}

\tcblower
\textbf{Solution détaillée :}
\begin{enumerate}
    \item \textbf{Tableau complété :}
    \begin{center}
    \begin{tabularx}{\linewidth}{|X|X|X|X|}
    \hline
    \rowcolor{popblueL}
    \textbf{Verbe} & \textbf{Prétérit} & \textbf{Participe passé} & \textbf{Nom clé} \\
    \hline
    \eng{buy} & \eng{bought} & \eng{bought} & \eng{buyer} \\
    \hline
    \eng{sell} & \eng{sold} & \eng{sold} & \eng{sale} \\
    \hline
    \eng{pay} & \eng{paid} & \eng{paid} & \eng{payment} \\
    \hline
    \eng{spend} & \eng{spent} & \eng{spent} & \eng{spending} \\
    \hline
    \end{tabularx}
    \end{center}
    \item \textbf{Phrase complétée :} \eng{Last week, Mrs Epote \textbf{bought} plantains and \textbf{paid} 3,000 francs. She \textbf{spent} less than usual because they were \textbf{cheap}.}
    \item \textbf{Traduction :} « La semaine dernière, Mme Epote a acheté des plantains et a payé 3 000 francs. Elle a dépensé moins que d'habitude parce qu'ils étaient bon marché. »
    \item \textbf{Justification :} le repère \eng{last week} impose le prétérit simple ; les quatre verbes sont irréguliers (\eng{bought}, \eng{paid}, \eng{spent}) ; \eng{cheap} reste à la forme positive car il n'y a pas de comparaison explicite.
\end{enumerate}
\end{exoresolu}

\courssub{Practice exercises}

\begin{exercice}[Word formation and spelling]
\textbf{A. Complétez le tableau de dérivation.}
\begin{center}
\begin{tabularx}{\linewidth}{|X|X|X|X|}
\hline
\rowcolor{popblueL}
\textbf{Verbe} & \textbf{Nom d'agent (-er)} & \textbf{Nom d'action} & \textbf{Adjectif} \\
\hline
\eng{buy} & \ldots & \ldots & --- \\
\hline
\eng{sell} & \ldots & \eng{sale} & --- \\
\hline
\eng{pay} & --- & \ldots & \eng{payable} \\
\hline
\eng{trade} & \ldots & \ldots & --- \\
\hline
\eng{bargain} & \eng{bargainer} & \ldots & --- \\
\hline
\end{tabularx}
\end{center}

\textbf{B. Choisissez l'orthographe correcte.}

1. The \eng{(seller / saler / sellar)} gave me a discount.

2. I keep the \eng{(receipt / receit / reciept)} for the warranty.

3. She loves \eng{(shoping / shopping / shoppping)} for clothes.

4. These plantains are very \eng{(cheap / sheap / cheep)} today.

5. The car is \eng{(for sell / for sale / for sales)}.

\textbf{C. Prononciation : associez chaque mot à son approximation française.}
\begin{center}
\begin{tabular}{ll}
1. \eng{bargain} & a. « ri-site » \\
2. \eng{receipt} & b. « ba-guine » \\
3. \eng{customer} & c. « keu-sto-meur » \\
4. \eng{cheap} & d. « tchipe » \\
5. \eng{sheep} & e. « chipe » \\
\end{tabular}
\end{center}
\end{exercice}

\begin{corrige}[Exercice Word formation and spelling]
\textbf{A. Tableau complété :}
\begin{center}
\begin{tabularx}{\linewidth}{|X|X|X|X|}
\hline
\rowcolor{popblueL}
\textbf{Verbe} & \textbf{Nom d'agent} & \textbf{Nom d'action} & \textbf{Adjectif} \\
\hline
\eng{buy} & \eng{buyer} & \eng{buying / purchase} & --- \\
\hline
\eng{sell} & \eng{seller} & \eng{sale / selling} & --- \\
\hline
\eng{pay} & --- & \eng{payment} & \eng{payable} \\
\hline
\eng{trade} & \eng{trader} & \eng{trading} & --- \\
\hline
\eng{bargain} & \eng{bargainer} & \eng{bargaining} & --- \\
\hline
\end{tabularx}
\end{center}

\textbf{B. Orthographe correcte :}

1. \textbf{seller} (double \eng{l} pour protéger la voyelle courte).

2. \textbf{receipt} (\eng{p} muet ; \eng{ei} après \eng{c}).

3. \textbf{shopping} (double \eng{p}, règle 1-1-1).

4. \textbf{cheap} (\eng{ch} = « tch », pas « ch »).

5. \textbf{for sale} (nom \eng{sale}, un seul \eng{l}, expression figée).

\textbf{C. Correspondances :}

1 $\rightarrow$ b (« ba-guine ») \quad 2 $\rightarrow$ a (« ri-site ») \quad 3 $\rightarrow$ c (« keu-sto-meur ») \quad 4 $\rightarrow$ d (« tchipe ») \quad 5 $\rightarrow$ e (« chipe »).

\textbf{Attention :} ne pas confondre 4 et 5 — \eng{cheap} « tchipe » (bon marché) et \eng{sheep} « chipe » (mouton) ne diffèrent que par la première consonne.
\end{corrige}

\begin{experience}[Jeu de rôle : « Au marché de Buea » — pronunciation focus]
\textbf{Objectif :} réemployer les familles de mots (\eng{buy / sell / pay / trade}) en soignant la prononciation de \eng{bargain}, \eng{receipt}, \eng{customer}, \eng{cheap}.
\medskip

\textbf{Déroulement (groupes de 3) :}
\begin{enumerate}
    \item \textbf{Rôles :} un \eng{customer} (client), un \eng{seller} (vendeur), un \eng{observer} (observateur).
    \item \textbf{Scénario :} le client veut acheter des denrées (\eng{plantains}, \eng{groundnuts}, \eng{rice}). Il demande le prix, \eng{bargains}, \eng{pays}, puis demande un \eng{receipt}.
    \item \textbf{Contrainte linguistique :} le dialogue doit contenir au minimum une forme de \eng{buy}, une de \eng{sell}, une de \eng{pay}, une de \eng{trade} (ou \eng{trader}), ainsi que \eng{cheap / expensive}, \eng{bargain}, \eng{receipt} et \eng{customer}.
    \item \textbf{Rôle de l'observer :} il note sur une grille si la prononciation des quatre mots pièges est correcte et signale les formes fautives *\eng{payed}, *\eng{selled}, *\eng{shoping}.
    \item \textbf{Restitution :} chaque groupe joue sa scène devant la classe, qui vote pour la « meilleure prononciation commerciale ».
\end{enumerate}
\end{experience}

\coursec{Using the vocabulary in context}

Dans la section précédente, nous avons appris à prononcer et à orthographier correctement les mots de la famille du commerce (\eng{buy — buyer}, \eng{sell — seller — sale}, \eng{bargain}, \eng{discount}...). Mais connaître un mot ne suffit pas : il faut savoir l'\cle{employer en contexte}, c'est-à-dire dans une phrase, un dialogue ou un texte réel. C'est exactement ce que l'examinateur attend de vous au probatoire et au baccalauréat. Cette dernière section est donc consacrée à la \textbf{réutilisation active} de tout le lexique étudié : lecture d'un récit, dialogue modèle, exercices d'application et production guidée.

\courssub{Reading: a short text — “A Saturday at Marché Central”}

Lisons d'abord un récit original qui réemploie naturellement le vocabulaire des sections 1 à 6. Lisez-le deux fois : une première fois pour le sens général, une deuxième fois en soulignant les mots du champ lexical du marché.

\begin{textebox}[A Saturday at Marché Central — reading text]
Last Saturday, I went to Marché Central in Yaoundé with a shopping list. My mother had given me some cash and a basket, and she had warned me: “Never accept the first price!”

At the entrance, the stalls were full of goods: fresh tomatoes, onions, groundnuts, smoked fish and bunches of plantains. Tomatoes were cheap that day, but dried fish was expensive because of scarcity — the fishermen had not brought much from the coast.

I stopped at a stall where a friendly vendor was selling plantains. “How much is this bunch?” I asked. “Two thousand francs,” she replied. I bargained with her for a few minutes, and finally she agreed to reduce the price to one thousand five hundred francs. I was happy: I had got a good discount!

Then I bought a kilo of rice and half a kilo of sugar at a nearby shop. The shopkeeper weighed the rice on his scale and put everything in a plastic bag. I paid in cash, and he gave me my change and a receipt, and I checked the change carefully before leaving.

On my way home, I realised I had spent less than planned. I was proud of myself: bargaining at the market is a real skill, and practice makes perfect!
\end{textebox}

\textbf{Glossaire (mots difficiles) :}

\begin{center}
\begin{tabularx}{\linewidth}{|l|l|X|}
\hline
\rowcolor{popblueL} English word & Nature & Français \\
\hline
shopping list & nom & liste de courses \\
\hline
cash & nom (invariable) & argent liquide, espèces \\
\hline
stall & nom & étal, stand (au marché) \\
\hline
goods & nom (pluriel) & marchandises, denrées \\
\hline
bunch & nom & régime (de plantains, de bananes) \\
\hline
scarcity & nom & rareté, pénurie \\
\hline
vendor & nom & vendeur / vendeuse (au marché) \\
\hline
to bargain & verbe & marchander, négocier le prix \\
\hline
discount & nom & remise, réduction \\
\hline
shopkeeper & nom & commerçant, boutiquier \\
\hline
to weigh & verbe & peser \\
\hline
change & nom (invariable) & monnaie (rendue) \\
\hline
receipt & nom & reçu \\
\hline
\end{tabularx}
\end{center}

\textbf{Comprehension questions:}
\begin{enumerate}
\item Where did the narrator go last Saturday, and with what?
\item Why was dried fish expensive that day?
\item How much did the vendor first ask for the bunch of plantains?
\item What final price did the narrator get after bargaining?
\item What did the shopkeeper give the narrator after payment?
\end{enumerate}

\begin{corrige}[Comprehension questions]
\begin{enumerate}
\item He/She went to Marché Central in Yaoundé, with a shopping list, some cash and a basket.
\item Dried fish was expensive because of scarcity: the fishermen had not brought much from the coast.
\item The vendor first asked two thousand francs.
\item The narrator got the bunch for one thousand five hundred francs (a discount of five hundred francs).
\item The shopkeeper gave him/her the change and a receipt.
\end{enumerate}
\end{corrige}

\begin{exemplebox}[Réemploi du lexique — exemples traduits]
\eng{I always make a shopping list before going to the market.} « Je fais toujours une liste de courses avant d'aller au marché. »

\eng{Prices go up when there is a scarcity of foodstuffs.} « Les prix augmentent quand il y a une pénurie de denrées alimentaires. »

\eng{The vendor reduced the price and gave me a discount.} « La vendeuse a baissé le prix et m'a fait une remise. »

\eng{Don't forget to check your change and ask for a receipt.} « N'oublie pas de vérifier ta monnaie et de demander un reçu. »
\end{exemplebox}

\courssub{Model dialogue: bargaining for a bunch of plantains}

Voici un dialogue complet et naturel entre une cliente et une vendeuse au marché. Observez la \textbf{structure en cinq étapes} : salutation → demande → annonce du prix → négociation → paiement et remerciements.

\begin{textebox}[At the market — a model dialogue]
Customer: Good morning, Madam! How are you today?

Vendor: I'm fine, thank you. Welcome! What can I do for you?

Customer: I would like a bunch of plantains. How much is this one?

Vendor: This big bunch costs two thousand francs, my dear.

Customer: Two thousand? That's too expensive! Can you reduce the price?

Vendor: These plantains are fresh and sweet. But for you, one thousand eight hundred.

Customer: Please, give me a discount. I'll pay one thousand five hundred.

Vendor: Ah, you bargain well! All right, agreed — one thousand five hundred.

Customer: Thank you very much. Here is the money, in cash.

Vendor: And here is your change and your receipt. Have a nice day!

Customer: You too, Madam. Goodbye!
\end{textebox}

\begin{popfigure}
\begin{center}
\begin{tikzpicture}[node distance=0.35cm]
\node[draw=popblue, fill=popblueL, rounded corners, minimum width=2.6cm, minimum height=0.9cm, align=center] (g) {\textbf{1. Greeting}\\ \emph{Good morning!}};
\node[draw=popgreen, fill=popgreenL, rounded corners, minimum width=2.6cm, minimum height=0.9cm, align=center, right=of g] (a) {\textbf{2. Asking}\\ \emph{How much is...?}};
\node[draw=poporange, fill=poporangeL, rounded corners, minimum width=2.6cm, minimum height=0.9cm, align=center, right=of a] (b) {\textbf{3. Bargaining}\\ \emph{Can you reduce...?}};
\node[draw=poppurple, fill=poppurpleL, rounded corners, minimum width=2.6cm, minimum height=0.9cm, align=center, right=of b] (p) {\textbf{4. Paying}\\ \emph{Here is the money.}};
\node[draw=poppink, fill=poppinkL, rounded corners, minimum width=2.6cm, minimum height=0.9cm, align=center, right=of p] (t) {\textbf{5. Thanking}\\ \emph{Thank you! Bye!}};
\draw[-{Stealth}, thick, popdark] (g) -- (a);
\draw[-{Stealth}, thick, popdark] (a) -- (b);
\draw[-{Stealth}, thick, popdark] (b) -- (p);
\draw[-{Stealth}, thick, popdark] (p) -- (t);
\end{tikzpicture}
\end{center}
\legende{Les cinq étapes d'un dialogue d'achat au marché : Greeting → Asking → Bargaining → Paying → Thanking.}
\end{popfigure}

\begin{methode}[How to write a market dialogue in five steps]
\begin{enumerate}
\item \textbf{Greetings} : les deux personnages se saluent poliment (\eng{Good morning! How are you?}).
\item \textbf{Customer's request} : le client demande un article et son prix (\eng{I would like... / How much is...?}).
\item \textbf{Vendor's answer and price} : le vendeur répond et annonce le prix (\eng{It costs... francs.}).
\item \textbf{Bargaining} : le client négocie, le vendeur fait une remise (\eng{That's too expensive! / Can you reduce the price? / Agreed!}).
\item \textbf{Payment and thanks} : paiement, monnaie, reçu, remerciements et au revoir (\eng{Here is your change. / Have a nice day!}).
\end{enumerate}
\end{methode}

\begin{attention}[Vary your speech verbs!]
Dans un dialogue rédigé, évitez de répéter \eng{say} à chaque réplique. Variez les \cle{verbes de parole} selon l'intention : \eng{ask} (demander), \eng{reply} (répondre), \eng{suggest} (suggérer), \eng{agree} (accepter), \eng{insist} (insister), \eng{add} (ajouter), \eng{exclaim} (s'exclamer). Exemple : \eng{“That's too expensive!” the customer exclaimed.} « “C'est trop cher !” s'exclama le client. »
\end{attention}

\courssub{Practice 1: gap-fill exercise}

\begin{exoresolu}[Complete the sentences with the target vocabulary]
Énoncé — Complete with: \emph{change, receipt, bargain, stall, discount, cash, scarcity, shopping list}.

1. Before going to the market, I always write a \dots\ so that I forget nothing.
2. The tomatoes are displayed on that wooden \dots\ over there.
3. Because of the \dots\ of onions, prices have gone up this month.
4. In Cameroonian markets, customers usually \dots\ before paying.
5. The vendor gave me a ten per cent \dots\ because I bought three kilos.
6. I don't have a bank card; I always pay in \dots.
7. After paying, check your \dots\ to be sure the amount is correct.
8. The shopkeeper gave me a \dots\ as proof of payment.

\tcblower
Solution détaillée :
1. \textbf{shopping list} — « liste de courses » ; collocation : \eng{make/write a shopping list}.
2. \textbf{stall} — « étal » ; attention, on dit \eng{at a stall}, pas \eng{on}.
3. \textbf{scarcity} — « rareté, pénurie » ; la montée des prix est la conséquence logique.
4. \textbf{bargain} — « marchander » ; verbe conjugué au présent avec \emph{customers} (pluriel).
5. \textbf{discount} — « remise » ; collocation : \eng{give/get a discount}.
6. \textbf{cash} — « espèces » ; collocation : \eng{pay in cash} (et non \emph{pay by cash}).
7. \textbf{change} — « monnaie rendue » ; collocation : \eng{check your change}.
8. \textbf{receipt} — « reçu » ; attention à l'orthographe : \textbf{ei} après le \emph{c}, et le \emph{p} ne se prononce pas (« ri-ssite »).
\end{exoresolu}

\courssub{Practice 2: rewriting — removing French calques}

Les francophones traduisent souvent mot à mot. Réécrivons des phrases fautives en anglais correct.

\begin{center}
\begin{tabularx}{\linewidth}{|X|X|}
\hline
\rowcolor{popblueL} Calque du français (à éviter) & English correct \\
\hline
I go to the market for \emph{make my expenses}. & I go to the market to do my shopping. \\
\hline
\emph{How much costs} this bunch? & How much does this bunch cost? / How much is this bunch? \\
\hline
The vendor \emph{rendered me the money}. & The vendor gave me my change. \\
\hline
This fish is \emph{very dear}. & This fish is very expensive. \\
\hline
I \emph{negotiated} the price with her. (calque acceptable mais lourd) & I bargained with her over the price. \\
\hline
Give me \emph{a reduction}. & Give me a discount. \\
\hline
\end{tabularx}
\end{center}

\begin{attention}[Pièges de traduction]
\begin{itemize}
\item \eng{dear} signifie « cher » au sens affectif (\eng{my dear friend}) ; pour un prix, utilisez \eng{expensive}.
\item On ne « rend » pas l'argent avec \emph{render} : on dit \eng{to give change} ou \eng{to give the change back}.
\item \eng{to do the shopping} = « faire les courses » ; \eng{to go shopping} = « aller faire du shopping ». Ne dites jamais \emph{make the shopping}.
\end{itemize}
\end{attention}

\begin{exercice}[Rewrite these sentences in correct English]
1. Je fais mes dépenses au marché tous les samedis. (calque : \emph{make my expenses})
2. Combien coûte ce poisson fumé ? (attention à l'ordre des mots)
3. La vendeuse m'a rendu cinq cents francs.
4. Ces tomates sont très chères aujourd'hui.
5. Il m'a donné un reçu après le paiement.
\end{exercice}

\begin{corrige}[Rewriting exercise]
1. \eng{I do my shopping at the market every Saturday.}
2. \eng{How much does this smoked fish cost?} ou \eng{How much is this smoked fish?}
3. \eng{The vendor gave me back five hundred francs (in change).}
4. \eng{These tomatoes are very expensive today.}
5. \eng{He gave me a receipt after the payment.}
\end{corrige}

\courssub{Guided production: write your own mini-dialogue}

\begin{experience}[Write an 8-line dialogue between a customer and a vendor]
En binôme, rédigez un mini-dialogue de \textbf{8 répliques} (4 par personnage) entre un client et un vendeur au marché de votre quartier (Mokolo, Marché Central, Buea market...). Suivez la frise en cinq étapes : Greeting → Asking → Bargaining → Paying → Thanking. Puis jouez le dialogue devant la classe.
\end{experience}

\textbf{Critères de réussite (checklist) :}
\begin{itemize}
\item[$\square$] Au moins \textbf{8 mots du lexique} de la leçon (stall, vendor, bunch, cash, change, receipt, goods, foodstuff...).
\item[$\square$] Au moins \textbf{2 collocations} correctes (\eng{make a shopping list}, \eng{pay in cash}, \eng{give a discount}, \eng{check your change}...).
\item[$\square$] Au moins \textbf{1 formule de négociation} (\eng{Can you reduce the price? / That's too expensive! / I'll give you...}).
\item[$\square$] \textbf{Aucun faux ami} ni calque du français (pas de \emph{dear} pour « cher », pas de \emph{make the shopping}).
\item[$\square$] Des verbes de parole variés si le dialogue est rédigé avec incises (\eng{ask, reply, suggest, agree}).
\end{itemize}

\begin{exemplebox}[Sample opening lines to help you]
\eng{Customer: Good afternoon! I need some foodstuff for the week.}

\eng{Vendor: Welcome! My goods are fresh and my prices are fair.}

\eng{Customer: How much is a kilo of groundnuts?}

\eng{Vendor: Five hundred francs, but I can give you a discount if you take two kilos.}
\end{exemplebox}

\begin{savaistu}[Dialogues and the Baccalauréat]
Au baccalauréat, les épreuves d'expression écrite et orale récompensent la \textbf{richesse de langue}. Un dialogue bien construit — avec du lexique précis (\eng{stall, vendor, bargain, discount, receipt}), des collocations naturelles et des formules idiomatiques (\eng{That's a bargain!}, \eng{It's a deal!}) — peut rapporter des points précieux dans les critères « vocabulary » et « communicative competence ». Entraînez-vous donc à réutiliser ce lexique dans toutes vos productions !
\end{savaistu}

\begin{aretenir}[L'essentiel de la section]
\begin{itemize}
\item Le lexique du marché se réemploie dans trois types de tâches : \textbf{lire} un récit, \textbf{jouer ou rédiger} un dialogue, \textbf{corriger} des calques du français.
\item Un dialogue d'achat suit cinq étapes : Greeting → Asking → Bargaining → Paying → Thanking.
\item Collocations à mobiliser : \eng{make a shopping list}, \eng{pay in cash}, \eng{get a discount}, \eng{give change}, \eng{ask for a receipt}.
\item Pièges à bannir : \emph{dear} pour « cher (coûteux) », \emph{make the shopping}, \emph{render the money}, ordre des mots français dans \emph{How much costs...?}.
\end{itemize}
\end{aretenir}

\newpage

\coursec{Activité d'intégration}

\coursec{Lesson 15 — Vocabulary: Buying and selling at a nearby store/market}

\courssub{Objectifs de l'activité}
Cette activité d'intégration vise à mobiliser l'ensemble du vocabulaire, des collocations, des expressions idiomatiques et des structures grammaticales étudiées dans la leçon 15 (\eng{Buying and selling at a nearby store/market}) au sein d'une tâche de communication orale authentique. À l'issue de cette séance, l'élève doit être capable de :
\begin{itemize}
    \item \textbf{Comprendre} des énoncés oraux relatifs aux prix, aux quantités, aux modalités de paiement et à la négociation dans un contexte de marché camerounais.
    \item \textbf{Produire} un dialogue interactif et cohérent en utilisant correctement le lexique ciblé (articles, contenants, unités de mesure, verbes d'action, adjectifs de qualité).
    \item \textbf{Manipuler} les formes grammaticales clés : noms comptables/non comptables (\eng{countable/uncountable nouns}), quantificateurs (\eng{some, any, much, many, a lot of, a few, a little}), impératif pour le service client, formes interrogatives (\eng{How much...? How many...? Can I have...?}).
    \item \textbf{Identifier} et \textbf{éviter} les faux-amis et calques du français repérés dans la leçon (ex. \eng{receipt} vs \eng{recipe}, \eng{change} vs \eng{exchange}, \eng{expensive} vs \eng{expansive}).
    \item \textbf{Évaluer} la production orale de ses pairs à l'aide d'une grille de critères précis (lexique, prononciation, interaction, réalisation de la tâche).
\end{itemize}

\courssub{Situation}
\begin{textebox}[Contexte : Mokolo Market, Yaoundé — Saturday Morning]
Mokolo Market is the largest and busiest market in Yaoundé. It is early Saturday morning, around 8:00 AM. The alleys are already crowded with buyers and sellers. The air smells of fresh spices, ripe plantains, and smoked fish. You are in the food section, near the stalls selling staples: rice, oil, sugar, tomatoes, and plantains.
\end{textebox}

\begin{textebox}[Document A : Vendor's Price List (Display Board)]
\textbf{MAMA NGONO'S STALL — Fresh \& Quality Goods}
\begin{center}
\begin{tabular}{|l|c|l|}
\hline
\rowcolor{popblueL} \textbf{Article} & \textbf{Unit} & \textbf{Price (FCFA)} \\
\hline
Plantains (ripe) & Bunch (\eng{un régime}) & 2,500 \\
Rice (Thai, 1st choice) & Bag (\eng{un sac} 50 kg) & 28,000 \\
Rice (Local, 2nd choice) & Bag (\eng{un sac} 50 kg) & 22,000 \\
Palm Oil (Red) & Tin (\eng{une boîte} 20 L) & 18,500 \\
Vegetable Oil (Refined) & Tin (\eng{une boîte} 20 L) & 16,000 \\
Tomatoes (Fresh) & Heap (\eng{un tas} / \eng{a pile}) & 1,000 \\
Sugar (White, refined) & Packet (\eng{un paquet} 1 kg) & 800 \\
\hline
\end{tabular}
\end{center}
\textit{Note: Prices are “starting prices” (asking prices). Bargaining is expected.}
\end{textebox}

\begin{textebox}[Document B : Customer's Shopping List \& Budget]
\textbf{Shopper: Pierre (Student in 1ère A, living in Mvog-Ada)}
\begin{itemize}
    \item Budget total: \textbf{10,000 FCFA} (exact cash: one 10,000 note).
    \item Needs:
    \begin{enumerate}
        \item 1 bunch of ripe plantains (for tonight's \eng{ndolé}).
        \item 1 small quantity of rice (approx. 2-3 kg, not a full bag).
        \item 1 tin of vegetable oil (small size if possible, or share a big tin? — but budget is tight).
        \item 2 heaps of tomatoes.
        \item 1 packet of sugar (1 kg).
    \end{enumerate}
\end{itemize}
\end{textebox}

\courssub{Consignes}
\textbf{Phase 1 : Préparation individuelle (10 min)}
\begin{enumerate}
    \item \textbf{Analyse des documents :} Lisez les documents A et B. Calculez le coût total \emph{estimé} aux prix affichés (sans négociation). Pierre a-t-il assez d'argent ? Identifiez les articles pour lesquels il \textbf{doit} négocier (au moins deux) et celui qui pose un problème de conditionnement (huile, riz).
    \item \textbf{Préparation lexicale :} Sur votre brouillon, listez 10 mots/expressions cibles de la leçon que vous comptez utiliser (ex. \eng{a bunch of, a heap of, a tin of, a packet of, to bargain, to give a discount, out of stock, change, receipt, expensive, cheap, fresh, ripe, stale}). Cochez 2 collocations (ex. \eng{drive a hard bargain, tight budget}) et 1 idiome (ex. \eng{to cost an arm and a leg — à éviter car trop fort, préférer "it's a bit pricey"}).
    \item \textbf{Préparation grammaticale :} Écrivez 3 questions que Pierre posera (prix, quantité, remise) et 2 phrases que le vendeur dira (proposition de remise, rupture de stock). Vérifiez : \eng{How much is...? Can you give me a discount? I'll take it. Here is your change.}
\end{enumerate}

\textbf{Phase 2 : Mise en scène par binômes (15 min)}
\begin{enumerate}
    \setcounter{enumi}{3}
    \item Formez des binômes. Désignez l'élève A (\textbf{Vendeur : Mama Ngono ou son apprenti}) et l'élève B (\textbf{Client : Pierre}).
    \item \textbf{Contraintes du jeu de rôle :}
    \begin{itemize}
        \item Le \textbf{Vendeur} : Accueille le client. Annonce qu'un article de la liste (au choix : \eng{rice} ou \eng{oil}) est \eng{out of stock} pour le conditionnement demandé (ex. pas de petit sachet d'huile, pas de riz en vrac aujourd'hui). Propose une alternative ou un \eng{discount} sur un autre article. Donne la monnaie correcte. Remet un \eng{receipt} (un petit papier écrit à la main).
        \item Le \textbf{Client} : Salue. Demande les prix de 3 articles minimum. Négocie \textbf{au moins deux articles} (utiliser \eng{Can you lower the price? / That's too expensive. / What's your last price?}). Vérifie la monnaie rendue. Demande explicitement \eng{a receipt, please}.
        \item Les deux : Utilisez le vocabulaire des quantités (\eng{a bunch, a heap, a tin, a packet, a kilo, half a kilo}). Durée cible : 2 à 3 minutes.
    \end{itemize}
    \item Répétez une fois en inversant les rôles si le temps le permet.
\end{enumerate}

\textbf{Phase 3 : Présentation devant la classe \& Évaluation par les pairs (15 min)}
\begin{enumerate}
    \setcounter{enumi}{6}
    \item Chaque binôme (ou 4-5 binômes volontaires si l'effectif est grand) joue sa scène devant la classe.
    \item \textbf{Grille d'écoute active (pour les observateurs) :} Chaque élève remplit cette grille pendant la représentation.
\end{enumerate}

\begin{center}
\begin{tabularx}{\linewidth}{|X|c|c|c|}
\hline
\rowcolor{popblueL} \textbf{Critère d'écoute (Cochez si entendu)} & \textbf{Binôme 1} & \textbf{Binôme 2} & \textbf{Binôme 3} \\
\hline
\eng{a bunch of plantains} / \eng{a heap of tomatoes} / \eng{a tin of oil} / \eng{a packet of sugar} / \eng{a bag of rice} & & & \\
\hline
\eng{How much is...?} / \eng{How many...?} / \eng{Can I have...?} & & & \\
\hline
\eng{It costs...} / \eng{That's \dots FCFA} / \eng{The price is...} & & & \\
\hline
\eng{Can you give me a discount?} / \eng{Can you lower the price?} / \eng{What's your last price?} & & & \\
\hline
\eng{I'll give you a discount} / \eng{Special price for you} / \eng{Let's say \dots} & & & \\
\hline
\eng{Out of stock} / \eng{We've run out of...} / \eng{Not available today} & & & \\
\hline
\eng{Here is your change} / \eng{Your change is \dots} / \eng{receipt} & & & \\
\hline
\textbf{Collocations entendues (ex: \eng{drive a hard bargain, tight budget, fresh produce, regular customer})} & \_\_\_\_\_\_ & \_\_\_\_\_\_ & \_\_\_\_\_\_ \\
\hline
\textbf{Idiome / Expression fonctionnelle (ex: \eng{it's a deal, a steal, worth every penny})} & \_\_\_\_\_\_ & \_\_\_\_\_\_ & \_\_\_\_\_\_ \\
\hline
\textbf{Aucun faux-ami détecté (\eng{recipe} pour \eng{receipt}, \eng{expansive} pour \eng{expensive}, \eng{exchange} pour \eng{change})} & Oui/Non & Oui/Non & Oui/Non \\
\hline
\end{tabularx}
\end{center}

\textbf{Phase 4 : Vote et Débriefing (5 min)}
\begin{enumerate}
    \setcounter{enumi}{8}
    \item La classe vote pour le \textbf{"Best Bargain"} (le client qui a obtenu le meilleur rapport qualité/prix tout en restant poli) et le \textbf{"Best Seller"} (le vendeur le plus professionnel).
    \item L'enseignant corrige les erreurs majeures entendues (prononciation de \eng{receipt} / « ri-sit »/, \eng{change} / « tchéindj »/, \eng{vegetable} / « vé-djè-te-ble »/, accord \eng{much/many}, faux-amis) et valide les réussites.
\end{enumerate}

\courssub{Ressources à mobiliser}
\begin{propriete}[Rappel : Quantificateurs et Noms (Countable / Uncountable)]
\begin{tabular}{|l|l|l|}
\hline
\rowcolor{popblueL} \textbf{Type} & \textbf{Exemples du marché (noms d'unités / biens)} & \textbf{Quantificateurs / Questions} \\
\hline
\eng{Countable (singulier)} & \eng{a bunch, a heap, a tin, a packet, a bag, a kilo} & \eng{How many...? A few / A number of} \\
\hline
\eng{Countable (pluriel)} & \eng{bunches, heaps, tins, packets, bags, kilos} & \eng{How many...? Many / Several} \\
\hline
\eng{Uncountable} & \eng{rice, oil, sugar, flour, water, money, change} & \eng{How much...? Much / A little / A lot of} \\
\hline
\end{tabular}
\emph{Attention :} Les mots \eng{bunch, heap, tin, packet, bag, kilo} sont des \textbf{noms d'unités} comptables. Les biens qu'ils quantifient peuvent être comptables (\eng{plantains, tomatoes, eggs}) ou non comptables (\eng{rice, oil, sugar}). On ne dit pas \eng{"a rice"} mais \eng{"a bag of rice / a kilo of rice / some rice"}. On ne dit pas \eng{"two sugars"} (sauf pour deux morceaux/sachets) mais \eng{"two packets of sugar"}.
\end{propriete}

\begin{propriete}[Expressions fonctionnelles clés — Achat / Vente / Négociation]
\begin{center}
\begin{tabularx}{\linewidth}{|X|X|}
\hline
\rowcolor{popblueL} \textbf{Fonction} & \textbf{Formules types (Client $\leftrightarrow$ Vendeur)} \\
\hline
Accueillir / Demander & \eng{Good morning. How can I help you? / What are you looking for?} \\
\hline
Demander le prix & \eng{How much is this bunch of plantains? / How much for two heaps?} \\
\hline
Annoncer le prix & \eng{It's 2,500 FCFA. / They are 1,000 FCFA each. / That comes to 5,000 FCFA.} \\
\hline
Trouver cher / Négocier & \eng{That's too expensive / a bit pricey. / Can you give me a discount? / What's your last price? / I only have 10,000 FCFA.} \\
\hline
Proposer remise / Refuser & \eng{Okay, I'll give you 200 FCFA off. / Let's say 2,300. / I can't go lower, it's a fixed price. / Special price for you: 2,200.} \\
\hline
Rupture de stock & \eng{Sorry, we are out of stock. / We've run out of small tins of oil. / Not available today.} \\
\hline
Accepter / Acheter & \eng{Alright, it's a deal. / I'll take it. / I'll take two heaps then.} \\
\hline
Payer / Monnaie / Reçu & \eng{Here is 10,000 FCFA. / Here is your change: 1,500 FCFA. / Can I have a receipt, please? / Here you are.} \\
\hline
\end{tabularx}
\end{center}
\end{propriete}

\begin{attention}[Faux-amis et Pièges pour Francophones — Vigilance active]
\begin{itemize}
    \item \cle{Receipt} / « ri-sit » (reçu) $\neq$ \cle{Recipe} / « ré-si-pi » (recette de cuisine). \textit{Piège :} \eng{"Can I have a recipe?"} $\rightarrow$ Le vendeur vous donne une recette de ndolé !
    \item \cle{Change} / « tchéindj » (monnaie rendue) $\neq$ \cle{Exchange} / « iks-tchéindj » (échange, bureau de change). \textit{Correct :} \eng{"Here is your change."}
    \item \cle{Expensive} / « ik-spen-siv » (cher) $\neq$ \cle{Expansive} / « ik-spen-siv » (vaste, étendu). \textit{Correct :} \eng{"It's too expensive."}
    \item \cle{Article} (article, objet) vs \cle{Article} (article de journal). Au marché, on dit plutôt \eng{item}, \eng{product}, \eng{goods}, \eng{stuff}.
    \item \cle{Quantity} / « kwan-ti-ti » (quantité) $\neq$ \cle{Quality} / « kwa-li-ti » (qualité). Prononciation distincte !
    \item \textbf{Calque :} \eng{"I take it"} (anglicisme oral accepté) mais standard : \eng{"I'll take it"} (futur immédiat / décision sur le moment). Éviter \eng{"I take"} (présent simple = habitude).
    \item \textbf{Calque :} \eng{"It costs 1000"} (correct) mais pas \eng{"It costs 1000 francs"} (redondant si contexte FCFA clair). Préférer \eng{"It's 1,000."}
\end{itemize}
\end{attention}

\begin{aretenir}[Lexique ciblé minimum pour la grille d'évaluation (Mots cibles)]
\textbf{Articles/Contenants :} \cle{a bunch (of plantains/bananas)}, \cle{a heap/pile (of tomatoes/onions)}, \cle{a tin/can (of oil/sardines)}, \cle{a packet/pack (of sugar/biscuits)}, \cle{a bag/sack (of rice/maize)}, \cle{a bottle (of oil/wine)}, \cle{a crate (of eggs/beer)}, \cle{a kilo}, \cle{half a kilo}, \cle{a litre}.\\
\textbf{Verbes/Action :} \cle{to buy}, \cle{to sell}, \cle{to pay (for)}, \cle{to cost}, \cle{to charge}, \cle{to bargain / to haggle}, \cle{to negotiate}, \cle{to give a discount / a reduction}, \cle{to run out of (stock)}, \cle{to hand over}, \cle{to wrap up}.\\
\textbf{Adjectifs/Qualité :} \cle{fresh}, \cle{ripe}, \cle{stale}, \cle{rotten}, \cle{damaged}, \cle{expensive}, \cle{cheap}, \cle{affordable}, \cle{pricey}, \cle{reasonable}, \cle{fixed (price)}.\\
\textbf{Paiement :} \cle{cash}, \cle{note}, \cle{coin}, \cle{change}, \cle{receipt}, \cle{bill (US) / invoice}, \cle{budget}, \cle{total}, \cle{amount}.\\
\textbf{Collocations :} \cle{asking price}, \cle{final price}, \cle{tight budget}, \cle{regular customer}, \cle{fresh produce}, \cle{drive a hard bargain}, \cle{make a deal}.\\
\textbf{Idiomes :} \eng{It's a deal.} / \eng{It's a steal (très bon marché).} / \eng{Worth every penny.} / \eng{To cost an arm and a leg (très cher — à reconnaître, pas forcément à produire).}
\end{aretenir}

\courssub{Proposition de production et corrigé commenté}
\begin{exoresolu}[Modèle de dialogue : Pierre (Client) vs Mama Ngono (Vendeuse)]
\textbf{Contexte :} Pierre arrive au stand. Il a 10 000 FCFA. Le riz en vrac et les petites boîtes d'huile sont en rupture.

\tcblower

\textbf{Mama Ngono:} Good morning, my son! Welcome to Mama Ngono's stall. \eng{How can I help you today?} \\
\textbf{Pierre:} Good morning, Mama. \eng{I'm looking for some plantains, rice, oil, tomatoes and sugar.} \\
\textbf{Mama Ngono:} \eng{You've come to the right place! Everything is fresh today.} Look at these plantains — \eng{ripe and sweet.} \eng{How many bunches do you want?} \\
\textbf{Pierre:} \eng{Just one bunch, please.} \eng{How much is it?} \\
\textbf{Mama Ngono:} \eng{2,500 FCFA for a big bunch.} \\
\textbf{Pierre:} \eng{2,500? That's a bit pricey.} \eng{Can you give me a discount? I'm a regular customer.} \\
\textbf{Mama Ngono:} \textit{[laughs]} \eng{Okay, okay! You drive a hard bargain.} \eng{2,300 FCFA. Final price.} \\
\textbf{Pierre:} \eng{Alright. I'll take it.} \eng{Now, I need some rice. Not a full bag, just two or three kilos.} \\
\textbf{Mama Ngono:} \eng{Ah, my dear,} \eng{we've run out of loose rice today.} \eng{We only have 50 kg bags — 22,000 or 28,000 FCFA.} \eng{Out of stock for small quantities.} \\
\textbf{Pierre:} \eng{Oh no! That's over my budget.} \eng{What about the oil? I need a small tin.} \\
\textbf{Mama Ngono:} \eng{Same problem.} \eng{The small tins are finished.} \eng{Only 20-litre tins left.} \eng{Vegetable oil is 16,000, groundnut oil 18,500.} \\
\textbf{Pierre:} \eng{Too expensive for me.} \eng{Let's focus on tomatoes and sugar then.} \eng{How much for two heaps of tomatoes?} \\
\textbf{Mama Ngono:} \eng{1,000 FCFA per heap. So 2,000 for two.} \\
\textbf{Pierre:} \eng{Can you lower it to 1,800 for the two heaps?} \\
\textbf{Mama Ngono:} \eng{1,900. Not a franc less. They are very fresh.} \\
\textbf{Pierre:} \eng{It's a deal.} \eng{And the packet of sugar?} \\
\textbf{Mama Ngono:} \eng{800 FCFA.} \\
\textbf{Pierre:} \eng{Okay. Let's calculate.} \eng{Plantains 2,300 + Tomatoes 1,900 + Sugar 800 = 5,000 FCFA.} \eng{Here is 10,000 FCFA.} \\
\textbf{Mama Ngono:} \textit{Counts money.} \eng{That's 5,000 FCFA total.} \eng{Here is your change: 5,000 FCFA.} \textit{Hands over a handwritten paper.} \eng{And here is your receipt. Thank you! Come again!} \\
\textbf{Pierre:} \eng{Thank you, Mama. Goodbye!}

\textbf{Commentaire des choix de langue (Niveau A2+/B1) :}
\begin{enumerate}
    \item \textbf{Formules de politesse ancrées :} \eng{"My son", "My dear"} (marqueurs camerounais d'affection/respect), \eng{"Welcome to..."}.
    \item \textbf{Quantificateurs corrects :} \eng{"one bunch", "two or three kilos", "loose rice" (riz en vrac), "two heaps", "a packet"}.
    \item \textbf{Négociation réaliste :} Le client utilise \eng{"a bit pricey"} (registre courant) + argument \eng{"regular customer"}. La vendeuse cède peu (\eng{200 FCFA}) puis tient bon sur tomates (\eng{"Not a franc less"}). Utilisation de \eng{"drive a hard bargain"} (collocation cible).
    \item \textbf{Gestion rupture de stock :} \eng{"We've run out of..."} (phrasal verb vu en leçon) + \eng{"Out of stock for small quantities"} (précision contextuelle).
    \item \textbf{Calcul et paiement :} Récapitulatif chiffré en anglais (\eng{"Let's calculate... Plantains 2,300 +..."}). Demande de monnaie implicite par le paiement \eng{"Here is 10,000"}. Réponse vendeuse standard : \eng{"Here is your change"} + \eng{"receipt"} (mot cible, orthographié/prononcé correctement).
    \item \textbf{Absence de faux-amis :} Pas de \eng{recipe}, \eng{expansive}, \eng{exchange}.
    \item \textbf{Prononciation visée :} \eng{receipt} [« ri-sit »], \eng{change} [« tchéindj »], \eng{vegetable} [« vé-djè-te-ble »], \eng{budget} [« ba-djè-t »].
\end{enumerate}
\end{exoresolu}

\courssub{Bilan de l'activité}
\begin{aretenir}[Auto-évaluation de l'élève — À remplir après l'activité]
Cochez la case correspondant à votre niveau de maîtrise \textbf{après} le jeu de rôle et la correction collective.
\begin{center}
\begin{tabularx}{\linewidth}{|X|c|c|c|}
\hline
\rowcolor{popblueL} \textbf{Compétence / Savoir-faire} & \textbf{Non acquis} & \textbf{En cours} & \textbf{Maîtrisé} \\
\hline
Je connais le vocabulaire des articles, contenants et unités de mesure (bunch, heap, tin, packet, bag, kilo). & & & \\
\hline
J'utilise correctement \eng{How much / How many} et les quantificateurs (\eng{some, any, much, many, a lot of}). & & & \\
\hline
Je sais demander le prix, négocier (\eng{discount, last price, deal}) et réagir à une rupture de stock (\eng{out of stock, run out of}). & & & \\
\hline
Je sais conclure un achat : payer, vérifier la monnaie (\eng{change}), demander un reçu (\eng{receipt}). & & & \\
\hline
J'emploie au moins 2 collocations et 1 expression idiomatique/fonctionnelle de la leçon. & & & \\
\hline
J'évite les faux-amis (\eng{receipt/recipe, change/exchange, expensive/expansive}) et les calques (\eng{"I take"}). & & & \\
\hline
Ma prononciation des mots-clés est intelligible (\eng{receipt, vegetable, change, budget, onion, plantain}). & & & \\
\hline
J'interagis de façon fluide : je réponds à mon partenaire, je ne récite pas un texte appris par cœur. & & & \\
\hline
\end{tabularx}
\end{center}

\end{aretenir}

\begin{methode}[Pour progresser : Prochaines étapes]
\begin{enumerate}
    \item \textbf{Revoir le lexique :} Refaire les exercices de la leçon 15 (ex. matching articles/contenants, gap-fill quantificateurs) en ciblant vos erreurs notées sur la grille.
    \item \textbf{Entraînement oral :} Enregistrez-vous (téléphone) en rejouant le dialogue seul (deux rôles) ou avec un camarade. Écoutez : entendez-vous \eng{receipt} sans le \eng{p} ? \eng{change} avec \eng{tch} ?
    \item \textbf{Écriture :} Rédigez le \eng{receipt} que Mama Ngono a donné à Pierre (format réel : date, articles, prix unitaires, total, monnaie, signature). Cela fixe l'orthographe du vocabulaire technique.
    \item \textbf{Extension :} Préparez une liste de courses pour un autre plat camerounais (\eng{eru, koki, achu, poulet DG}) avec les quantités et un budget. Simulez l'achat mentalement en anglais.
\end{enumerate}
\end{methode}

\begin{savaistu}[Le saviez-vous ? — Le marché de Mokolo et l'économie informelle]
Mokolo Market (\eng{Mokolo Market}), situé au cœur de Yaoundé, est l'un des plus grands marchés d'Afrique centrale. Il incarne l'\eng{informal sector} (secteur informel) qui emploie plus de 80\% de la population active au Cameroun. Les vendeuses comme "Mama Ngono" sont souvent des \eng{wholesalers} (grossistes) et \eng{retailers} (détaillantes) à la fois. La négociation (\eng{haggling / bargaining}) n'est pas une simple coutume : c'est un mécanisme de \eng{price discovery} (découverte du prix) où l'information est asymétrique. Le \eng{receipt} manuscrit sur un carnet à souche (\eng{duplicate book}) reste la norme, bien que le \eng{mobile money} (Mobile Money / MoMo) transforme peu à peu les paiements \eng{cash}. Savoir "faire le marché" en anglais est une compétence professionnelle réelle pour le commerce transfrontalier avec le Nigeria ou pour le tourisme.
\end{savaistu}

\newpage

\coursec{Méthodes et exercices résolus}

\coursec{Lesson 15 — Vocabulary: Buying and Selling at the Market}

\courssub{Fiche méthode 1 — Identifier et employer le vocabulaire du thème}

\begin{methode}[Apprendre et réutiliser le lexique du marché]
Pour maîtriser le vocabulaire de l'achat et de la vente (\eng{buying and selling}), suis cette démarche :

\begin{enumerate}
\item \textbf{Classe le vocabulaire par catégories} : les lieux (\eng{market, stall, shop, store, supermarket}), les personnes (\eng{customer, seller, trader, vendor, shopkeeper}), les produits (\eng{foodstuff, groceries, articles, goods}) et l'argent (\eng{price, change, receipt, discount}).
\item \textbf{Apprends chaque mot avec son article et sa préposition} : on dit \eng{at the market}, \eng{in a shop}, \eng{to buy something from somebody}, \eng{to sell something to somebody}.
\item \textbf{Mémorise les collocations} (associations de mots) : \eng{to pay in cash}, \eng{to bargain over the price}, \eng{to reduce the price}, \eng{to weigh the tomatoes}.
\item \textbf{Construis une phrase personnelle} avec chaque mot nouveau : \eng{My mother buys vegetables at Mokolo market every Saturday.}
\item \textbf{Teste-toi par familles de mots} : \eng{to buy} (verbe) → \eng{a buyer} (personne) → \eng{a purchase} (nom de l'action).
\end{enumerate}

Réflexe d'examen : dans un exercice de vocabulaire, relis toujours la phrase entière avant de choisir le mot — c'est le \cle{contexte} qui impose la réponse, pas la ressemblance avec le français.
\end{methode}

\begin{exoresolu}[Compléter un texte à trous — type examen]
\textbf{Énoncé.} Complete the following passage with words from the box. Use each word once only.

\begin{center}
\begin{tabular}{|l|l|l|l|l|}
\hline
\rowcolor{popblueL} stall & change & bargain & groceries & vendor \\
\hline
\end{tabular}
\end{center}

Every Saturday, Mrs Etoundi goes to the central market in Yaoundé to buy her \ldots\ldots\ldots for the week. She stops at a fruit \ldots\ldots\ldots where the \ldots\ldots\ldots sells mangoes and bananas. In Cameroon, customers often \ldots\ldots\ldots over the price before paying. When she pays with a 5,000 FCFA note, the seller gives her the \ldots\ldots\ldots .

\tcblower
\textbf{Solution détaillée.}

Raisonnement phrase par phrase :

\begin{enumerate}
\item \eng{to buy her \ldots\ for the week} : on cherche un nom de produits de consommation courante. Le seul mot possible est \cle{groceries} (provisions, denrées). \eng{She buys her groceries for the week.} « Elle achète ses provisions pour la semaine. »
\item \eng{a fruit \ldots} : on cherche un lieu de vente au marché. Un étal de fruits se dit \cle{stall} : \eng{a fruit stall}. « un étal de fruits ». Attention : on ne dit pas \eng{a fruit shop} au marché, mais bien \eng{stall}.
\item \eng{where the \ldots\ sells mangoes} : le sujet du verbe \eng{sells} est une personne qui vend : \cle{vendor} (vendeur/vendeuse ambulante). \eng{The vendor sells mangoes.} « Le vendeur vend des mangues. »
\item \eng{customers often \ldots\ over the price} : il faut un verbe ; la préposition \eng{over} impose \cle{bargain} : \eng{to bargain over the price}. « marchander le prix ».
\item \eng{gives her the \ldots} : après un paiement avec un billet, le vendeur rend la monnaie : \cle{change}. \eng{He gives her the change.} « Il lui rend la monnaie. »
\end{enumerate}

\textbf{Texte complété (réponse rédigée) :}

\eng{Every Saturday, Mrs Etoundi goes to the central market in Yaoundé to buy her groceries for the week. She stops at a fruit stall where the vendor sells mangoes and bananas. In Cameroon, customers often bargain over the price before paying. When she pays with a 5,000 FCFA note, the seller gives her the change.}

« Chaque samedi, Mme Etoundi va au marché central de Yaoundé pour acheter ses provisions de la semaine. Elle s'arrête à un étal de fruits où le vendeur vend des mangues et des bananes. Au Cameroun, les clients marchandent souvent le prix avant de payer. Quand elle paie avec un billet de 5 000 FCFA, le vendeur lui rend la monnaie. »

\textbf{Conclusion.} Chaque réponse est justifiée par la grammaire (nom, verbe, sujet) et par la collocation (\eng{fruit stall}, \eng{bargain over}, \eng{give the change}).
\end{exoresolu}

\courssub{Fiche méthode 2 — Réemployer le lexique à l'oral : le dialogue au marché}

\begin{methode}[Conduire un dialogue d'achat en anglais]
Pour un jeu de rôle ou une production orale au marché, structure toujours ton échange en cinq étapes :

\begin{enumerate}
\item \textbf{Saluer et attirer l'attention} : \eng{Good morning! How much are the tomatoes?} / \eng{Can I help you, madam?}
\item \textbf{Demander le prix} : \eng{How much is a kilo of onions?} / \eng{What is the price of this bag of rice?}
\item \textbf{Marchander poliment} : \eng{That is too expensive! Can you reduce the price?} / \eng{I will give you 2,000 francs.}
\item \textbf{Conclure l'achat} : \eng{All right, I will take two kilos.} / \eng{Here is the money.}
\item \textbf{Terminer poliment} : \eng{Thank you. Here is your change. Have a nice day!}
\end{enumerate}

Réflexes linguistiques :
\begin{itemize}
\item \eng{How much} + nom indénombrable ou verbe au singulier : \eng{How much is the rice?} ; \eng{How many} + nom pluriel : \eng{How many oranges do you want?}
\item Pour la quantité : \eng{a kilo of}, \eng{a bag of}, \eng{a bunch of plantains}, \eng{a litre of oil}.
\item Reste poli : \eng{please}, \eng{Can you...?}, \eng{I would like...} (et non l'impératif sec, impoli en anglais).
\end{itemize}
\end{methode}

\begin{exoresolu}[Compléter un dialogue — type examen oral et écrit]
\textbf{Énoncé.} Complete the dialogue between a customer and a trader at Buea market with suitable expressions.

\textbf{Trader:} Good morning, sir! (1) \ldots\ldots\ldots\ldots ?

\textbf{Customer:} Yes, please. (2) \ldots\ldots\ldots\ldots is a bunch of plantains?

\textbf{Trader:} It is 3,000 francs.

\textbf{Customer:} That is expensive! (3) \ldots\ldots\ldots\ldots reduce the price?

\textbf{Trader:} For you, sir, I can sell it at 2,500 francs.

\textbf{Customer:} All right. (4) \ldots\ldots\ldots\ldots take it. Here is a 5,000-franc note.

\textbf{Trader:} Thank you. (5) \ldots\ldots\ldots\ldots is your change: 2,500 francs.

\tcblower
\textbf{Solution détaillée.}

\begin{enumerate}
\item Le vendeur ouvre le dialogue : formule d'accueil classique \eng{Can I help you?} « Puis-je vous aider ? / Vous désirez ? ». Le point d'interrogation et la réponse \eng{Yes, please} confirment cette formule.
\item Le client demande un prix : \eng{How much is a bunch of plantains?} « Combien coûte un régime de plantains ? ». On utilise \eng{How much} car le prix est indénombrable ; la réponse \eng{It is 3,000 francs} le confirme.
\item Le client marchande poliment : \eng{Can you reduce the price?} « Pouvez-vous baisser le prix ? ». L'auxiliaire \eng{Can} + sujet \eng{you} construit la question polie.
\item Le client accepte : \eng{I will take it.} « Je le prends. ». \eng{will} exprime ici une décision immédiate ; \eng{it} reprend \eng{a bunch of plantains} (singulier).
\item Le vendeur rend la monnaie : \eng{Here is your change.} « Voici votre monnaie. ». Inversion fréquente après \eng{Here} : \eng{Here is...}.
\end{enumerate}

\textbf{Dialogue complet (réponse rédigée) :}

\eng{Trader: Good morning, sir! Can I help you?}

\eng{Customer: Yes, please. How much is a bunch of plantains?}

\eng{Trader: It is 3,000 francs.}

\eng{Customer: That is expensive! Can you reduce the price?}

\eng{Trader: For you, sir, I can sell it at 2,500 francs.}

\eng{Customer: All right. I will take it. Here is a 5,000-franc note.}

\eng{Trader: Thank you. Here is your change: 2,500 francs.}

\textbf{Conclusion.} Un bon dialogue d'achat suit toujours la progression : salutation → question sur le prix → marchandage → décision → paiement et remerciements.
\end{exoresolu}

\courssub{Fiche méthode 3 — Éviter les faux amis et les calques du français}

\begin{methode}[Repérer et corriger un faux ami]
\begin{enumerate}
\item \textbf{Méfie-toi des mots qui ressemblent au français} : \eng{to assist} ne veut pas dire « assister à » (c'est \eng{to attend}) ; \eng{a magasin} ne se traduit pas par \eng{a magazine} (un magasin = \eng{a shop / a store} ; un magazine = \eng{a magazine}).
\item \textbf{Vérifie la préposition} : on dit \eng{to buy something from somebody} (et non \eng{to buy something to somebody}) ; \eng{to sell something to somebody} ; \eng{to pay for the goods} (payer \emph{pour} quelque chose).
\item \textbf{Attention aux calques de structure} : « Combien ça coûte ? » se dit \eng{How much does it cost?} ou \eng{How much is it?}, jamais \eng{How much it costs?} dans une question directe.
\item \textbf{Singulier/pluriel} : \eng{money} est indénombrable (\eng{the money is...}) ; \eng{news} est singulier ; mais \eng{goods} (marchandises) est toujours pluriel : \eng{the goods are expensive}.
\item \textbf{En cas de doute, reformule} avec un mot sûr plutôt que de risquer un faux ami.
\end{enumerate}

\begin{center}
\begin{tabularx}{\linewidth}{|X|X|X|}
\hline
\rowcolor{popblueL} Français & English correct & Piège à éviter \\
\hline
un magasin & a shop / a store & \eng{a magazine} = un magazine (revue) \\
\hline
la monnaie & the change & \eng{money} = l'argent en général \\
\hline
marchander & to bargain & \eng{to negotiate} = négocier (contrat) \\
\hline
une économie (épargne) & savings & \eng{economy} = l'économie (du pays) \\
\hline
acheter à quelqu'un & to buy from somebody & pas de \eng{to} après \eng{buy} \\
\hline
\end{tabularx}
\end{center}
\end{methode}

\begin{exoresolu}[Corriger les erreurs — type examen]
\textbf{Énoncé.} Each of the following sentences contains one mistake (false friend, wrong preposition or calque from French). Correct the mistake and rewrite the sentence.

\begin{enumerate}
\item I buy my books to the shopkeeper near the school.
\item She went to the magazine to buy some bread.
\item How much costs this bag of rice?
\item The goods is very expensive in this supermarket.
\item The trader gave me my money after I paid with a big note.
\end{enumerate}

\tcblower
\textbf{Solution détaillée.}

\begin{enumerate}
\item \textbf{Erreur :} \eng{buy ... to}. En anglais, on achète \emph{à} quelqu'un avec la préposition \eng{from}. Correction : \eng{I buy my books from the shopkeeper near the school.} « J'achète mes livres au libraire près de l'école. »
\item \textbf{Erreur :} faux ami \eng{magazine}. Un « magasin » se dit \eng{shop} ou \eng{store} ; \eng{magazine} désigne une revue. Correction : \eng{She went to the shop to buy some bread.} « Elle est allée au magasin acheter du pain. »
\item \textbf{Erreur :} calque de l'ordre des mots français. Une question directe exige l'auxiliaire \eng{does}. Correction : \eng{How much does this bag of rice cost?} « Combien coûte ce sac de riz ? » (ou \eng{How much is this bag of rice?}).
\item \textbf{Erreur :} accord. \eng{Goods} est toujours pluriel en anglais. Correction : \eng{The goods are very expensive in this supermarket.} « Les marchandises sont très chères dans ce supermarché. »
\item \textbf{Erreur :} faux ami de sens. Après un paiement, le vendeur rend \emph{la monnaie} : \eng{change}, et non \eng{money}. Correction : \eng{The trader gave me my change after I paid with a big note.} « Le commerçant m'a rendu la monnaie après que j'ai payé avec un gros billet. »
\end{enumerate}

\textbf{Conclusion.} La plupart des erreurs de lexique viennent d'une traduction mot à mot du français : pense toujours en anglais, avec la préposition et l'accord anglais.
\end{exoresolu}

\courssub{Fiche méthode 4 — Construire les familles de mots (Word Building)}

\begin{methode}[Former les mots de la même famille]
Pour les exercices de formation de mots (\eng{word building}), applique ces suffixes :

\begin{enumerate}
\item \textbf{Personne qui fait l'action} : verbe + \eng{-er} / \eng{-or} : \eng{buy → buyer}, \eng{sell → seller}, \eng{trade → trader}, \eng{consume → consumer}.
\item \textbf{Nom de l'action ou de la chose} : \eng{buy → purchase} (achat), \eng{sell → sale} (vente), \eng{pay → payment}, \eng{choose → choice}.
\item \textbf{Adjectifs} : \eng{expense → expensive}, \eng{value → valuable}, \eng{cheap} (bon marché, sans suffixe).
\item \textbf{Vérifie la nature grammaticale demandée} par la phrase : après un article (\eng{a/the}), il faut un nom ; avant un nom, un adjectif ; après un auxiliaire, un verbe.
\item \textbf{Attention aux modifications orthographiques} : \eng{choose → choice} (changement de voyelle), \eng{sell → sale} (et non \eng{sell → sell}).
\end{enumerate}

\begin{center}
\begin{tabular}{|l|l|l|}
\hline
\rowcolor{popblueL} Verbe & Nom (action/chose) & Personne \\
\hline
to buy & a purchase & a buyer \\
\hline
to sell & a sale & a seller \\
\hline
to pay & a payment & a payer \\
\hline
to trade & trade & a trader \\
\hline
to consume & consumption & a consumer \\
\hline
\end{tabular}
\end{center}
\end{methode}

\begin{exoresolu}[Word building — type examen]
\textbf{Énoncé.} Complete each sentence with the correct form of the word in brackets.

\begin{enumerate}
\item The \ldots\ldots\ldots of this shop are always polite. (sell)
\item Mobile money is a common method of \ldots\ldots\ldots in Douala. (pay)
\item My father made a good \ldots\ldots\ldots : he bought a bag of maize at a low price. (purchase → use the noun)
\item The \ldots\ldots\ldots of foodstuff has increased this year. (consume)
\item A \ldots\ldots\ldots sells goods to customers. (trade)
\end{enumerate}

\tcblower
\textbf{Solution détaillée.}

\begin{enumerate}
\item Après l'article \eng{the} et devant \eng{are} (pluriel), il faut un nom de personne au pluriel : \cle{sellers}. \eng{The sellers of this shop are always polite.} « Les vendeurs de ce magasin sont toujours polis. »
\item Après la préposition \eng{of}, il faut un nom : \cle{payment}. \eng{Mobile money is a common method of payment in Douala.} « L'argent mobile est un moyen de paiement courant à Douala. »
\item Après l'adjectif \eng{good}, il faut un nom : \cle{purchase}. \eng{My father made a good purchase: he bought a bag of maize at a low price.} « Mon père a fait un bon achat : il a acheté un sac de maïs à bas prix. »
\item Après \eng{the} et devant \eng{of}, il faut le nom de l'action : \cle{consumption}. \eng{The consumption of foodstuff has increased this year.} « La consommation de denrées alimentaires a augmenté cette année. »
\item Après l'article \eng{a}, il faut un nom de personne : \cle{trader}. \eng{A trader sells goods to customers.} « Un commerçant vend des marchandises aux clients. »
\end{enumerate}

\textbf{Conclusion.} Pour réussir ces exercices, identifie d'abord la nature grammaticale attendue (nom, verbe, adjectif), puis applique le bon suffixe en vérifiant l'orthographe.
\end{exoresolu}

\courssub{Fiche méthode 5 — Prononciation et orthographe : les pièges du thème}

\begin{propriete}[Mots difficiles à prononcer et à écrire]
\begin{center}
\begin{tabularx}{\linewidth}{|l|X|X|}
\hline
\rowcolor{popblueL} Mot & Prononciation (approx. fr.) & Piège orthographique \\
\hline
\eng{receipt} & « ri-si-t » & \eng{p} muet ; \eng{ei} se prononce \eng{/i:/} \\
\hline
\eng{vegetable} & « vedj-te-beul » & 3 syllabes seulement, pas 4 \\
\hline
\eng{customer} & « cos-to-meu » & \eng{u} après \eng{st} ; pas \eng{costumer} \\
\hline
\eng{business} & « bi-zeu-nis » & \eng{u} après \eng{b} ; \eng{i} se prononce \eng{/i/} \\
\hline
\eng{bargain} & « ba-ga-ne » & \eng{ai} = \eng{/a:/} ; \eng{g} dur \\
\hline
\eng{change} & « tchéindj » & \eng{ch} = \eng{/tʃ/} ; \eng{ge} = \eng{/dʒ/} \\
\hline
\eng{expensive} & « ik-spen-siv » & \eng{ex-} ; pas \eng{expansive} \\
\hline
\eng{stall} & « stol » & \eng{a} = \eng{/ɔ:/} (britannique) \\
\hline
\end{tabularx}
\end{center}
\end{propriete}

\begin{exemplebox}[Entraînement à la dictée]
Écris les phrases suivantes (dictée ou auto-dictée) :
\begin{enumerate}
\item \eng{The customer asked for a receipt at the vegetable stall.}
\item \eng{Business is good today; the goods are not expensive.}
\item \eng{Can you give me change for this 10,000 franc note?}
\end{enumerate}
\textbf{Correction :} Vérifie \eng{receipt}, \eng{vegetable}, \eng{customer}, \eng{business}, \eng{goods}, \eng{expensive}, \eng{change}, \eng{stall}.
\end{exemplebox}

\courssub{Fiche méthode 6 — Rédiger un court texte avec le lexique du thème}

\begin{methode}[Écrire un paragraphe sur les achats au marché]
Pour une production écrite courte (80–100 mots) sur le thème \eng{buying and selling} :

\begin{enumerate}
\item \textbf{Introduction} : situe la scène (qui, où, quand) : \eng{Every weekend, my aunt goes to the market in Bamenda.}
\item \textbf{Développement} : utilise au moins cinq mots du champ lexical : \eng{stall, customer, to bargain, price, change, foodstuff, vendor, cash, receipt, discount}.
\item \textbf{Ordre logique} : décrire le lieu → nommer les produits → raconter l'achat (prix, marchandage, paiement).
\item \textbf{Conclusion} : une phrase d'opinion : \eng{I think bargaining is part of our culture.}
\item \textbf{Relecture} : vérifie les prépositions (\eng{at the market}, \eng{buy from}), l'accord singulier/pluriel (\eng{goods are}) et l'orthographe (\eng{business}, \eng{receipt}, \eng{customer}).
\end{enumerate}
\end{methode}

\begin{exoresolu}[Composition guidée — type examen]
\textbf{Énoncé.} Write a paragraph of about 80–100 words describing how people buy and sell foodstuff at your local market. Use at least five of the following words: \eng{stall, vendor, to bargain, price, change, foodstuff, customer, cash}.

\tcblower
\textbf{Solution détaillée.}

\textbf{Étape 1 — Plan.} Introduction (lieu et moment) → les vendeurs et les produits → le marchandage → le paiement → opinion finale.

\textbf{Étape 2 — Choix du vocabulaire.} Je retiens six mots imposés : \eng{stall, vendor, to bargain, price, change, foodstuff}, plus \eng{customer} et \eng{cash}.

\textbf{Étape 3 — Rédaction (modèle de réponse) :}

\eng{Every morning, the market in my town is full of life. Vendors arrange their foodstuff on wooden stalls: tomatoes, plantains, fish and vegetables. Customers move from one stall to another, comparing prices. In Cameroon, people usually bargain before they buy anything. For example, if a vendor asks 1,000 francs for a bowl of tomatoes, the customer may offer 700 francs, and they finally agree on a fair price. Most people pay in cash, and the vendor gives them their change. I like going to the market because bargaining makes shopping lively and friendly.}

« Chaque matin, le marché de ma ville est plein de vie. Les vendeurs disposent leurs denrées sur des étals en bois : tomates, plantains, poisson et légumes. Les clients passent d'un étal à l'autre, comparant les prix. Au Cameroun, les gens marchandent généralement avant d'acheter quoi que ce soit. Par exemple, si un vendeur demande 1 000 francs pour une cuvette de tomates, le client peut proposer 700 francs, et ils se mettent finalement d'accord sur un prix juste. La plupart des gens paient en espèces, et le vendeur leur rend la monnaie. J'aime aller au marché parce que le marchandage rend les courses animées et conviviales. »

\textbf{Étape 4 — Vérification.} Le texte compte environ 95 mots ; il contient huit mots du champ lexical imposé ; les prépositions sont correctes (\eng{full of}, \eng{from one stall to another}, \eng{agree on}, \eng{pay in cash}) ; l'accord est respecté (\eng{vendors arrange}, \eng{the customer may offer}).

\textbf{Conclusion.} Un bon paragraphe de vocabulaire n'est pas une liste de mots : c'est un texte logique où chaque mot du thème est employé naturellement et correctement.
\end{exoresolu}

\begin{attention}[Les 5 erreurs qui coûtent des points au Bac]
\begin{enumerate}
\item \textbf{Faux ami \eng{magazine}.} Un « magasin » se dit \eng{a shop} ou \eng{a store} ; \eng{a magazine} est une revue. Écrire \eng{I went to the magazine to buy rice} est une erreur éliminatoire en expression écrite.
\item \textbf{Préposition après \eng{buy} et \eng{sell}.} On achète \emph{à} quelqu'un : \eng{to buy from somebody} ; on vend \emph{à} quelqu'un : \eng{to sell to somebody}. Ne mélange jamais les deux.
\item \textbf{Accord de \eng{goods} et \eng{money}.} \eng{Goods} est toujours pluriel (\eng{the goods are expensive}) ; \eng{money} est indénombrable et singulier (\eng{the money is on the table}). Interdit : \eng{a money}, \eng{the goods is}.
\item \textbf{Ordre des mots dans les questions de prix.} « Combien coûte ce tissu ? » = \eng{How much does this cloth cost?} ou \eng{How much is this cloth?}. Le calque \eng{How much costs this cloth?} est fautif : il faut l'auxiliaire \eng{does}.
\item \textbf{Orthographe des mots-clés du thème.} \eng{business} (et non \eng{buisness}), \eng{receipt} (le \eng{p} ne se prononce pas mais s'écrit), \eng{customer} (et non \eng{costumer}), \eng{expensive} (et non \eng{expansive}), \eng{change} pour « la monnaie » (et non \eng{money}).
\end{enumerate}
\end{attention}

\begin{savaistu}[Le marché au Cameroun : vocabulaire culturel]
Au Cameroun, les marchés sont des lieux de vie essentiels.
\begin{itemize}
\item \eng{Mokolo Market} (Yaoundé) et \eng{Marché Central} (Douala) sont les plus grands.
\item On y trouve des \eng{bendskins} (moto-taxis) pour transporter les \eng{goods}.
\item Le \eng{Mobile Money} (MTN MoMo, Orange Money) est devenu un \eng{common method of payment} aux côtés du \eng{cash}.
\item Les \eng{foodstuffs} typiques : \eng{plantains}, \eng{cassava} (manioc), \eng{cocoyams} (macabo), \eng{ndolé} (feuilles de bitterleaf), \eng{egusi} (graines de courge), \eng{bush pepper} (poivre de brousse).
\item Le \eng{bargaining} (marchandage) est une norme sociale, pas une impolitesse.
\end{itemize}
\end{savaistu}

\begin{experience}[Jeu de rôle : « Marché de Buea »]
\textbf{Consigne.} En binôme : l'un est \eng{trader} (vendeur d'articles variés : \eng{plantains, tomatoes, rice, oil}), l'autre est \eng{customer} (avec une liste de courses et 10 000 FCFA).
\begin{enumerate}
\item Le \eng{customer} demande le prix de 3 articles.
\item Le \eng{trader} annonce des prix élevés.
\item Le \eng{customer} \eng{bargains} pour chaque article.
\item Ils s'accordent sur un \eng{total price}.
\item Le \eng{customer} paie en \eng{cash} (billet de 10 000) ; le \eng{trader} rend le \eng{change} et donne un \eng{receipt} (optionnel).
\end{enumerate}
\textbf{Contrainte linguistique :} Utiliser obligatoirement : \eng{How much is...?}, \eng{Can you reduce...?}, \eng{I will take...}, \eng{Here is your change}. L'enseignant évalue la fluidité, la politesse (\eng{please, madam/sir}) et l'exactitude du vocabulaire.
\end{experience}

\begin{exercice}[Entraînement personnel — À rendre au professeur]
\begin{enumerate}
\item \textbf{Vocabulaire :} Traduire en anglais : 1. un étal de légumes 2. la monnaie 3. marchander 4. un reçu 5. payer en espèces 6. un client 7. des denrées alimentaires 8. un commerçant.
\item \textbf{Grammaire :} Choisir la bonne forme. 1. The goods (is / are) fresh today. 2. I bought this bread (from / to) the baker. 3. How much (does / is) this cost? 4. She gave me the (money / change). 5. This (magazine / shop) sells newspapers.
\item \textbf{Word building :} Compléter le tableau.
\begin{center}
\begin{tabular}{|l|l|l|}
\hline
\rowcolor{popblueL} Verbe & Nom (action) & Personne \\
\hline
to purchase & \ldots\ldots\ldots & \ldots\ldots\ldots \\
\hline
\ldots\ldots\ldots & a sale & \ldots\ldots\ldots \\
\hline
to pay & \ldots\ldots\ldots & \ldots\ldots\ldots \\
\hline
\end{tabular}
\end{center}
\item \textbf{Expression écrite :} Écris un dialogue de 8 répliques entre un vendeur de \eng{plantains} au marché de Bamenda et un client qui veut acheter deux régimes. Inclus : salutation, prix, marchandage, paiement (billet de 5 000), monnaie, remerciements.
\end{enumerate}
\end{exercice}

\begin{corrige}[Exercice n° 1 à 4]
\begin{enumerate}
\item 1. \eng{a vegetable stall} 2. \eng{change} 3. \eng{to bargain} 4. \eng{a receipt} 5. \eng{to pay in cash} 6. \eng{a customer} 7. \eng{foodstuff} (ou \eng{foodstuffs}) 8. \eng{a trader} / \eng{a vendor}.
\item 1. \eng{are} 2. \eng{from} 3. \eng{does} 4. \eng{change} 5. \eng{shop}.
\item \eng{to purchase} → \eng{a purchase} → \eng{a purchaser} ; \eng{to sell} → \eng{a sale} → \eng{a seller} ; \eng{to pay} → \eng{a payment} → \eng{a payer}.
\item \textbf{Dialogue modèle :}
\eng{Vendor: Good morning, madam! Can I help you?}
\eng{Customer: Yes. How much is a bunch of plantains?}
\eng{Vendor: 3,000 francs per bunch.}
\eng{Customer: That's too expensive. Can you reduce the price?}
\eng{Vendor: For you, 2,500 francs.}
\eng{Customer: All right. I will take two bunches. Here is 5,000 francs.}
\eng{Vendor: Thank you. Here is your change. Have a nice day!}
\eng{Customer: Thank you. Goodbye!}
\end{enumerate}
\end{corrige}

\newpage

\coursec{Exercices}

\courssub{Je vérifie mes connaissances}

\begin{exercice}[QCM : le bon mot]
Choisis la réponse correcte (a, b ou c) pour compléter chaque phrase.
\begin{enumerate}
\item Could I have a \ldots{} of rice, please?
\begin{itemize}
\item[a)] magazine \quad b) bag \quad c) bill
\end{itemize}
\item The tomatoes are cheap today: only 200 francs a \ldots{}.
\begin{itemize}
\item[a)] kilo \quad b) coin \quad c) wallet
\end{itemize}
\item ``How much does this shirt \ldots{}?'' --- ``Five thousand francs.''
\begin{itemize}
\item[a)] value \quad b) cost \quad c) price
\end{itemize}
\item I am sorry, madam, sugar is \ldots{} at the moment. Come back tomorrow.
\begin{itemize}
\item[a)] out of stock \quad b) in cash \quad c) on sale
\end{itemize}
\item Keep your \ldots{}: you will need it if you want to exchange the shoes.
\begin{itemize}
\item[a)] change \quad b) receipt \quad c) discount
\end{itemize}
\item The shopkeeper gave me a 10 \% \ldots{} because I bought three bags of flour.
\begin{itemize}
\item[a)] bargain \quad b) debt \quad c) discount
\end{itemize}
\end{enumerate}
\end{exercice}

\begin{textebox}[A morning at Mokolo market]
Every Saturday, Mrs Etame wakes up at five o'clock and goes to Mokolo market in Yaoundé. She sells fresh vegetables: tomatoes, onions, okra and bitterleaf. Her stall is near the entrance, so many customers stop there first. Mrs Etame buys her goods from farmers in the villages around the city. She always checks the quality before paying. In the morning, prices are high, but in the afternoon she reduces them because she does not want to carry the vegetables back home. Customers like to bargain with her, and she often gives a small discount to regular clients. Last Saturday, a young man tried to pay with a torn banknote, but she refused it politely and asked for another one. By six o'clock in the evening, she had sold everything and counted her money with a smile.
\end{textebox}

\begin{exercice}[Vrai ou faux ?]
Lis le texte \emph{A morning at Mokolo market} ci-dessus, puis indique si chaque affirmation est \textbf{true} ou \textbf{false}. Justifie ta réponse par une courte citation du texte.
\begin{enumerate}
\item Mrs Etame goes to the market every day.
\item Her stall is far from the entrance of the market.
\item She grows all her vegetables herself.
\item Prices are lower in the afternoon.
\item She accepted the torn banknote from the young man.
\end{enumerate}
\end{exercice}

\begin{exercice}[Vocabulaire : devine le mot]
Trouve le mot anglais du lexique du marché qui correspond à chaque définition.
\begin{enumerate}
\item The money a seller gives back when you pay with a note that is too big: \ldots{}
\item A person who sells goods in a shop: a shop\ldots{}
\item A long, narrow piece of paper that shows what you bought and how much you paid: a \ldots{}
\item To discuss the price in order to pay less: to \ldots{}
\item A small table or open structure where goods are sold in a market: a \ldots{}
\item Food and other products that a shop has available for sale: \ldots{}
\end{enumerate}
\end{exercice}

\begin{exercice}[Complète avec le mot de la bonne famille]
Complète chaque phrase avec la forme correcte du mot entre parenthèses (verbe, nom de personne, nom d'action ou adjectif).
\begin{enumerate}
\item My father is a \ldots{}: he sells phone credit near the post office. (sell)
\item The \ldots{} of the land took place last week. (sell)
\item Mobile money is a very common method of \ldots{} in Cameroon. (pay)
\item This second-hand bag is very \ldots{}: it costs only 1,000 francs. (afford)
\item Cameroon \ldots{} cocoa and coffee with many countries. (trade)
\item The \ldots{} of a new schoolbag is about 8,000 francs. (cost)
\end{enumerate}
\end{exercice}

\courssub{Je m'entraîne}

\begin{exercice}[Traduction]
Traduis les phrases suivantes en anglais, en utilisant le lexique étudié dans la leçon.
\begin{enumerate}
\item Combien coûte un sac de riz dans ce magasin ?
\item Le vendeur m'a rendu la monnaie sur un billet de dix mille francs.
\item Les arachides sont en rupture de stock aujourd'hui.
\item J'ai marchandé et j'ai obtenu une réduction de cinq cents francs.
\item Gardez votre reçu : il vous servira de preuve d'achat.
\item Ma mère fait ses courses au marché tous les samedis matin.
\end{enumerate}
\end{exercice}

\begin{textebox}[At the market]
Customer: Good morning, madam. How much does a bucket of garri (1) \ldots{}?

Seller: Good morning. It is 2,500 francs, my dear (2) \ldots{}.

Customer: That is expensive! Can we (3) \ldots{} a little?

Seller: Ah, you like to negotiate! Because you are a student, I will give you a (4) \ldots{} of 300 francs.

Customer: Thank you! Do you also sell smoked fish?

Seller: I am sorry, smoked fish is (5) \ldots{} today. The fishermen did not come.

Customer: No problem. Then I will take the garri. Here are 5,000 francs.

Seller: Do you have smaller notes? I do not have enough (6) \ldots{} this morning.

Customer: Yes, here are two notes of 1,000 and two coins of 100.

Seller: Perfect. I always prefer (7) \ldots{} to mobile money. Wait, let me write your (8) \ldots{}.

Customer: Thank you. Your (9) \ldots{} is very well organised, madam.

Seller: Come back next week! I promise you a good (10) \ldots{} on plantains.
\end{textebox}

\begin{exercice}[Texte à trous : dialogue au marché]
Complète le dialogue \emph{At the market} ci-dessus avec les dix mots de la banque suivante (un mot par espace) :

\begin{center}
\textbf{discount -- change -- receipt -- out of stock -- bargain -- cost -- stall -- cash -- customer -- price}
\end{center}
\end{exercice}

\begin{exercice}[Corrige les erreurs]
Chaque phrase contient une erreur typique des francophones (calque, faux ami ou ordre des mots). Réécris chaque phrase correctement.
\begin{enumerate}
\item How much costs this bag of onions?
\item Yesterday, I made the market with my aunt.
\item The merchant rended me the money.
\item I want to buy a magazine of rice for my mother.
\item She sells her products in a magazine near the school.
\item The price of the shoes are very expensive.
\end{enumerate}
\end{exercice}

\begin{exercice}[Appariement : attention aux faux amis]
Relie chaque mot français (1--8) à sa traduction anglaise correcte (a--h). Attention : certains mots anglais ressemblent à des mots français mais ont un sens différent.

\begin{center}
\begin{tabularx}{\linewidth}{|X|X|}
\hline
\rowcolor{popblueL} \textbf{Français} & \textbf{English} \\
\hline
1. un magasin & a) change \\
\hline
2. la monnaie (rendue) & b) a shop \\
\hline
3. un marchand & c) a bill \\
\hline
4. une pièce de monnaie & d) a trader \\
\hline
5. une facture & e) a coin \\
\hline
6. une librairie & f) a bookshop \\
\hline
7. le prix & g) the price \\
\hline
8. un cabas / un sac de courses & h) a shopping bag \\
\hline
\end{tabularx}
\end{center}
\end{exercice}

\courssub{Je me prépare au Bac}

\begin{textebox}[Mama Ngo Bell's provision store]
In the neighbourhood of Bonabéri, in Douala, everybody knows Mama Ngo Bell's provision store. It is a small wooden shop painted blue and white, just opposite the motorbike taxi stage. She opened it fifteen years ago with the money she had saved from selling puff-puff and beans at the roadside.

The shop sells almost everything a family needs: rice, groundnut oil, cubes, soap, matches, candles, exercise books and even phone credit. Mama Ngo Bell knows the price of every single article by heart. When a customer asks, ``How much is a tin of sardines?'', she answers immediately, without looking at any notebook.

Business is not always easy. Prices at the wholesale market in Sandaga change every month, and transporters have increased their fares because of fuel costs. ``If I buy dear, I must sell dear,'' she explains, ``otherwise I make a loss.'' But she is careful: if her prices are too high, customers will simply go to the big supermarket down the road.

Mama Ngo Bell also gives credit to a few trusted customers. She writes their names and debts in a small red notebook, and they pay at the end of the month, when salaries arrive. ``A shopkeeper must have a good memory and a good heart,'' she says with a laugh. ``Without customers, there is no business; without trust, there are no customers.''
\end{textebox}

\begin{exercice}[Compréhension et langue — type Bac]
Lis attentivement le texte \emph{Mama Ngo Bell's provision store} ci-dessus, puis réponds aux questions.

\textbf{A. Comprehension questions.} Réponds en anglais, par des phrases complètes.
\begin{enumerate}
\item Where is Mama Ngo Bell's shop situated? (2 marks)
\item How did she get the money to open her shop? (2 marks)
\item Mention four articles that customers can buy in her shop. (2 marks)
\item Give two reasons why business is difficult for her. (2 marks)
\item What does the red notebook contain? (2 marks)
\end{enumerate}

\textbf{B. Vocabulary in context.} Trouve dans le texte un mot ou une expression qui signifie :
\begin{enumerate}
\item a place where motorbike taxis wait for passengers (paragraph 1)
\item to remember something perfectly, without reading it (paragraph 2)
\item a large shop where traders buy goods in big quantities to resell them (paragraph 3)
\item money that you must pay back to someone (paragraph 4)
\end{enumerate}

\textbf{C. Language work.}
\begin{enumerate}
\item Turn into the passive voice: ``She sells almost everything a family needs.''
\item Turn into reported speech: ``If I buy dear, I must sell dear,'' she explains.
\item Complete with the correct preposition: ``She gives credit \ldots{} a few trusted customers.''
\end{enumerate}
\end{exercice}

\begin{exercice}[Traduction — type Bac]
Traduis le passage suivant en anglais correct et naturel, en soignant le vocabulaire de l'achat et de la vente.

\begin{textebox}[Texte à traduire]
Hier matin, mon frère et moi sommes allés au marché pour acheter des provisions. Les prix avaient augmenté à cause des pluies, alors nous avons marchandé avec chaque vendeur. Chez la marchande de légumes, j'ai payé avec un billet de cinq mille francs et elle m'a rendu la monnaie. Malheureusement, l'huile était en rupture de stock, donc nous devrons y retourner demain. Avant de partir, le vendeur de poisson nous a offert une petite réduction parce que nous sommes des clients réguliers.
\end{textebox}
\end{exercice}

\begin{exercice}[Composition — type Bac]
\textbf{Sujet :} You went to a nearby store to buy a schoolbag, but you had a problem: the price was too high and the shopkeeper did not want to bargain at first.

Write a dialogue of \textbf{at least ten exchanges} (about \textbf{150 to 180 words}) between you and the shopkeeper. In your dialogue, you must:
\begin{itemize}
\item greet the shopkeeper and ask for the article you need;
\item ask about the price and try to bargain politely;
\item ask for a discount or a special offer;
\item pay and ask for your change and a receipt;
\item use \textbf{at least eight words or expressions} from the vocabulary of buying and selling studied in this lesson (for example: \eng{How much does it cost?}, \eng{out of stock}, \eng{discount}, \eng{change}, \eng{receipt}, \eng{bargain}, \eng{in cash}, \eng{affordable}).
\end{itemize}

\textbf{Barème indicatif :} respect de la situation de communication et du format dialogue (4 marks), exactitude grammaticale (4 marks), richesse et pertinence du vocabulaire du thème (4 marks), cohérence et organisation (4 marks), orthographe et ponctuation (4 marks).
\end{exercice}



\newpage

\coursec{Corrigés des exercices}

\begin{corrige}[Exercice 1 — QCM : le bon mot]
\begin{enumerate}
\item \textbf{b) bag} --- \eng{a bag of rice} « un sac de riz ». \eng{Magazine} est un faux ami : il signifie « magazine, revue », jamais « magasin ». \eng{Bill} = « facture » ou « billet de banque » (US).
\item \textbf{a) kilo} --- \eng{200 francs a kilo} « 200 francs le kilo » ; la préposition \eng{a} exprime ici « le, la » (prix à l'unité).
\item \textbf{b) cost} --- \eng{How much does this shirt cost?} « Combien coûte cette chemise ? ». \eng{Price} est un nom, pas un verbe ; \eng{value} signifie « évaluer » ou « valeur ».
\item \textbf{a) out of stock} --- « en rupture de stock ». \eng{In cash} = « en espèces » ; \eng{on sale} = « en solde / en vente ».
\item \textbf{b) receipt} --- « reçu, ticket de caisse ». \eng{Change} = « la monnaie (rendue) » ; \eng{discount} = « réduction ».
\item \textbf{c) discount} --- \eng{a 10 \% discount} « une réduction de 10 \% ». \eng{Bargain} = « marchander » (verbe) ou « bonne affaire » (nom) ; \eng{debt} = « dette ».
\end{enumerate}
\end{corrige}

\begin{corrige}[Exercice 2 — Vrai ou faux ?]
\begin{enumerate}
\item \textbf{False.} ``Every Saturday, Mrs Etame wakes up at five o'clock and goes to Mokolo market.'' (Elle y va chaque samedi, pas chaque jour.)
\item \textbf{False.} ``Her stall is near the entrance.'' (\eng{near} « près de » est le contraire de \eng{far from} « loin de ».)
\item \textbf{False.} ``Mrs Etame buys her goods from farmers in the villages around the city.'' (Elle achète ses marchandises aux agriculteurs ; elle ne les cultive pas elle-même.)
\item \textbf{True.} ``In the morning, prices are high, but in the afternoon she reduces them.'' (\eng{to reduce} « réduire, baisser ».)
\item \textbf{False.} ``A young man tried to pay with a torn banknote, but she refused it politely.'' (\eng{torn} « déchiré », participe passé irrégulier de \eng{tear -- tore -- torn}.)
\end{enumerate}
\end{corrige}

\begin{corrige}[Exercice 3 — Vocabulaire : devine le mot]
\begin{enumerate}
\item \textbf{change} (la monnaie rendue ; invariable dans ce sens).
\item \textbf{shopkeeper} (le/la commerçant(e), gérant(e) d'un magasin).
\item \textbf{receipt} (le reçu ; attention au \emph{p} muet : prononciation « ri-ssite »).
\item \textbf{bargain} (marchander ; on dit aussi \eng{to haggle}).
\item \textbf{stall} (l'étal, l'étalage au marché).
\item \textbf{stock} (le stock, les marchandises disponibles ; d'où \eng{out of stock} « en rupture de stock » et \eng{in stock} « en stock »).
\end{enumerate}
\end{corrige}

\begin{corrige}[Exercice 4 — Complète avec le mot de la bonne famille]
\begin{enumerate}
\item \textbf{seller} --- nom de personne : \eng{sell} (verbe) $\rightarrow$ \eng{seller} « vendeur » (suffixe \emph{-er}).
\item \textbf{sale} --- nom d'action : \eng{sell} $\rightarrow$ \eng{sale} « vente » (forme irrégulière, comme \eng{sell -- sold -- sold}).
\item \textbf{payment} --- nom d'action : \eng{pay} $\rightarrow$ \eng{payment} « paiement » (suffixe \emph{-ment}).
\item \textbf{affordable} --- adjectif : \eng{afford} $\rightarrow$ \eng{affordable} « abordable » (suffixe \emph{-able}).
\item \textbf{trades} --- verbe conjugué à la 3\up{e} personne du singulier : \eng{Cameroon trades cocoa\ldots} « Le Cameroun fait du commerce de cacao\ldots ».
\item \textbf{cost} --- nom : \eng{the cost of} « le coût de ». Rappel : \eng{cost} est à la fois verbe et nom (\eng{cost -- cost -- cost}).
\end{enumerate}
\end{corrige}

\begin{corrige}[Exercice 5 — Traduction]
\begin{enumerate}
\item \eng{How much does a bag of rice cost in this shop?} (Attention à l'ordre des mots : auxiliaire \eng{does} avant le sujet ; jamais \emph{How much costs\ldots})
\item \eng{The seller gave me my change on a ten-thousand-franc note.} (On dit aussi \eng{gave me change for a 10,000-franc note}.)
\item \eng{Groundnuts are out of stock today.} (\eng{groundnuts} « arachides » ; on trouve aussi \eng{peanuts}.)
\item \eng{I bargained and got a discount of five hundred francs.} (\eng{get -- got -- got(ten)}.)
\item \eng{Keep your receipt: it will serve as proof of purchase.} (\eng{proof of purchase} « preuve d'achat ».)
\item \eng{My mother does her shopping at the market every Saturday morning.} (On dit \eng{to do one's shopping} « faire ses courses », jamais \emph{to make the market}.)
\end{enumerate}
\end{corrige}

\begin{corrige}[Exercice 6 — Texte à trous : dialogue au marché]
\begin{enumerate}
\item \textbf{cost} --- \eng{How much does a bucket of garri cost?} (verbe après l'auxiliaire \eng{does}).
\item \textbf{customer} --- la vendeuse s'adresse à sa « cliente ».
\item \textbf{bargain} --- « marchander » (verbe après \eng{can}).
\item \textbf{discount} --- \eng{a discount of 300 francs} « une réduction de 300 francs ».
\item \textbf{out of stock} --- « en rupture de stock » (les pêcheurs ne sont pas venus).
\item \textbf{change} --- la vendeuse n'a pas assez de monnaie à rendre.
\item \textbf{cash} --- \eng{to prefer cash to mobile money} « préférer les espèces au mobile money » (\eng{prefer \ldots{} to \ldots}).
\item \textbf{receipt} --- \eng{to write a receipt} « écrire un reçu ».
\item \textbf{stall} --- « votre étal est très bien organisé ».
\item \textbf{price} --- \eng{a good price on plantains} « un bon prix sur les plantains ».
\end{enumerate}
\end{corrige}

\begin{corrige}[Exercice 7 — Corrige les erreurs]
\begin{enumerate}
\item \eng{How much does this bag of onions cost?} --- En anglais, la question exige l'auxiliaire \eng{do/does} ; on ne peut pas inverser directement le verbe \eng{cost} comme en français (« combien coûte\ldots ? »).
\item \eng{Yesterday, I went to the market with my aunt.} ou \eng{I did the shopping with my aunt.} --- « Faire le marché » ne se traduit pas par \emph{to make the market} ; \eng{to go to the market} = « aller au marché ».
\item \eng{The shopkeeper gave me my change.} --- « Rendre la monnaie » = \eng{to give (back) the change} ; le verbe français « rendre » ne se calque pas (\emph{rended} n'existe pas). \eng{Merchant} existe mais est rare et soutenu ; on dit \eng{shopkeeper} ou \eng{trader}.
\item \eng{I want to buy a bag of rice for my mother.} --- Faux ami : \eng{magazine} = « revue, magazine » ; « un sac de riz » = \eng{a bag of rice}.
\item \eng{She sells her products in a shop near the school.} --- Même faux ami : « un magasin » = \eng{a shop} (GB) ou \eng{a store} (US).
\item \eng{The price of the shoes is very high.} --- Deux erreurs : le sujet \eng{the price} est singulier (\eng{is}), et un prix est \eng{high} ou \eng{low}, jamais \eng{expensive} ; c'est l'article qui est \eng{expensive} ou \eng{cheap}.
\end{enumerate}
\end{corrige}

\begin{corrige}[Exercice 8 — Appariement : attention aux faux amis]
1--b ; 2--a ; 3--d ; 4--e ; 5--c ; 6--f ; 7--g ; 8--h.

Pièges à retenir : \eng{a library} = « une bibliothèque » et non « une librairie » (\eng{a bookshop}) ; \eng{a coin} = « une pièce de monnaie » tandis que \eng{change} = « la monnaie rendue » ; \eng{a bill} = « une facture » (GB) ou « un billet de banque » (US) ; \eng{the price} se prononce « praïce » et s'écrit avec \emph{-ce}, à ne pas confondre avec \eng{a prize} « un prix, une récompense ».
\end{corrige}

\begin{corrige}[Exercice 9 — Compréhension et langue — type Bac]
\textbf{A. Comprehension questions.}
\begin{enumerate}
\item Her shop is situated in the neighbourhood of Bonabéri, in Douala, just opposite the motorbike taxi stage.
\item She got the money from her savings: she had saved money by selling puff-puff and beans at the roadside.
\item Customers can buy, for example, rice, groundnut oil, soap and exercise books (autres réponses possibles : cubes, matches, candles, phone credit).
\item Business is difficult because prices at the wholesale market change every month, and because transporters have increased their fares because of fuel costs.
\item The red notebook contains the names of her trusted customers and their debts.
\end{enumerate}

\textbf{B. Vocabulary in context.}
\begin{enumerate}
\item \eng{the motorbike taxi stage}
\item \eng{to know something by heart}
\item \eng{the wholesale market}
\item \eng{debts} (ou \eng{a debt})
\end{enumerate}

\textbf{C. Language work.}
\begin{enumerate}
\item \eng{Almost everything a family needs is sold (by her).} --- Passif : le complément d'objet devient sujet ; \eng{sells} (présent) devient \eng{is sold} (auxiliaire \eng{be} au présent + participe passé ; \eng{sell -- sold -- sold}).
\item \eng{She explains that if she buys dear, she must sell dear.} --- Discours rapporté au présent (\eng{explains}) : les temps ne reculent pas ; seuls les pronoms s'adaptent (\eng{I} $\rightarrow$ \eng{she}).
\item \eng{She gives credit \textbf{to} a few trusted customers.} --- \eng{to give something to somebody} « donner quelque chose à quelqu'un ».
\end{enumerate}

\textbf{Points de vigilance :} en compréhension, réponds toujours par une phrase complète qui reprend les mots de la question ; en vocabulaire, recopie le mot exact du texte sans le modifier ; au passif, vérifie l'accord de l'auxiliaire \eng{be} avec le nouveau sujet.
\end{corrige}

\begin{corrige}[Exercice 10 — Traduction — type Bac]
\textbf{Proposition de traduction :}

\eng{Yesterday morning, my brother and I went to the market to buy provisions. Prices had gone up because of the rains, so we bargained with every seller. At the vegetable seller's stall, I paid with a five-thousand-franc note and she gave me my change. Unfortunately, oil was out of stock, so we will have to go back tomorrow. Before we left, the fish seller gave us a small discount because we are regular customers.}

\textbf{Points de vigilance :}
\begin{itemize}
\item « mon frère et moi » = \eng{my brother and I} (jamais \emph{me and my brother} en tête de phrase) ;
\item « les prix avaient augmenté » = plus-que-parfait \eng{had gone up} (action antérieure à l'action principale) ;
\item « chez la marchande » = \eng{at the vegetable seller's (stall)} ;
\item « rendre la monnaie » = \eng{to give (back) the change}, pas de calque de « rendre » ;
\item « nous devrons » = futur \eng{we will have to} ;
\item « clients réguliers » = \eng{regular customers}.
\end{itemize}

\textbf{Barème indicatif :} exactitude grammaticale (temps, auxiliaires, ordre des mots) 8 marks ; pertinence du lexique du thème 6 marks ; orthographe et ponctuation 3 marks ; fluidité et naturel 3 marks.
\end{corrige}

\begin{corrige}[Exercice 11 — Composition — type Bac]
\textbf{Modèle de dialogue (environ 165 mots) :}

\eng{Me: Good morning, sir. Do you sell schoolbags?}

\eng{Shopkeeper: Good morning. Yes, we do. These ones just arrived from Douala.}

\eng{Me: They are beautiful. How much does the blue one cost?}

\eng{Shopkeeper: It is 12,000 francs.}

\eng{Me: Twelve thousand! That is not affordable for a student. Can we bargain a little?}

\eng{Shopkeeper: I am sorry, my prices are fixed.}

\eng{Me: Please, sir. My mother gave me only 9,000 francs. Can you give me a discount?}

\eng{Shopkeeper: Hmm\ldots{} Because schools are reopening, I can offer you a special discount of 2,000 francs.}

\eng{Me: Thank you very much! Do you also sell exercise books?}

\eng{Shopkeeper: Unfortunately, they are out of stock today. Come back on Friday.}

\eng{Me: No problem. Here are 10,000 francs, in cash.}

\eng{Shopkeeper: And here is your change. Let me write your receipt.}

\eng{Me: Thank you, sir. Goodbye!}

\eng{Shopkeeper: Goodbye, and tell your friends about my shop!}

\textbf{Barème indicatif :} respect de la situation de communication et du format dialogue (4 marks), exactitude grammaticale (4 marks), richesse et pertinence du vocabulaire du thème (4 marks), cohérence et organisation (4 marks), orthographe et ponctuation (4 marks).

\textbf{Points de vigilance :} annonce clairement les deux personnages ; alterne des répliques courtes ; place naturellement au moins huit mots du thème (\eng{cost}, \eng{bargain}, \eng{discount}, \eng{affordable}, \eng{out of stock}, \eng{in cash}, \eng{change}, \eng{receipt}) ; vérifie les questions avec \eng{do/does} et l'orthographe de \eng{receipt} (« ri-ssite », \emph{p} muet) et de \eng{bargain} ; termine par une salutation polie.
\end{corrige}

\newpage

\coursec{Fiche bilan}

\begin{popfigure}
\begin{tikzpicture}[mindmap, grow cyclic, every node/.style={concept, text=white, align=center, font=\small},
root concept/.append style={concept color=popdark, font=\bfseries\small},
level 1/.append style={level distance=3.6cm, sibling angle=51.4},
level 2/.append style={level distance=2.1cm, sibling angle=40, font=\scriptsize}]
\node[root concept]{Buying and Selling}
child[concept color=popblue]{ node {Places and people}
  child { node {market, stall, shop, store} }
  child { node {customer, seller, trader, hawker} }
}
child[concept color=poporange]{ node {Articles and foodstuff}
  child { node {goods, items, products} }
  child { node {foodstuff: yams, garri, palm oil} }
}
child[concept color=popgreen]{ node {Prices and payment}
  child { node {price, cost, change, receipt} }
  child { node {cash, pay, afford, owe} }
}
child[concept color=poppurple]{ node {Bargaining}
  child { node {bargain, haggle, discount} }
  child { node {“How much is it?”} }
}
child[concept color=poppink]{ node {Collocations and idioms}
  child { node {a good bargain, value for money} }
  child { node {to rip off, to cost a fortune} }
}
child[concept color=popgold]{ node {False friends}
  child { node {magasin = shop (not magazine)} }
  child { node {monnaie = change} }
}
child[concept color=popmuted]{ node {Word families}
  child { node {buy / buyer / purchase} }
  child { node {sell / seller / sale} }
};
\end{tikzpicture}
\legende{Carte mentale : le vocabulaire de l'achat et de la vente au marché et à la boutique}
\end{popfigure}

\begin{bilan}
\textbf{1. Lexique clé (à connaître par cœur)}
\begin{itemize}
\item \textbf{Lieux :} \eng{market} (marché), \eng{stall} (étal), \eng{shop/store} (boutique, magasin), \eng{supermarket} (supermarché), \eng{roadside vendor} (vendeur au bord de la route).
\item \textbf{Personnes :} \eng{customer} (client), \eng{seller/vendor/trader} (vendeur, commerçant), \eng{shopkeeper} (boutiquier), \eng{hawker} (colporteur), \eng{wholesaler/retailer} (grossiste/détaillant).
\item \textbf{Marchandises :} \eng{goods} (marchandises), \eng{items/articles} (articles), \eng{foodstuff} (denrées alimentaires), \eng{groceries} (provisions), \eng{fresh produce} (produits frais).
\item \textbf{Prix et paiement :} \eng{price} (prix), \eng{cheap/expensive} (bon marché/cher), \eng{discount} (remise), \eng{change} (monnaie), \eng{receipt} (reçu), \eng{to pay in cash} (payer en espèces), \eng{to afford} (avoir les moyens de), \eng{to owe} (devoir de l'argent).
\item \textbf{Négociation :} \eng{to bargain/to haggle} (marchander), \eng{a bargain} (une bonne affaire), \eng{to reduce/lower the price} (baisser le prix), \eng{fixed price} (prix fixe).
\end{itemize}

\textbf{2. Collocations et expressions idiomatiques}
\begin{itemize}
\item \eng{value for money} (bon rapport qualité-prix) ; \eng{to cost a fortune} (coûter une fortune) ; \eng{to rip somebody off} (arnaquer quelqu'un) ; \eng{to shop around} (comparer les prix) ; \eng{to sell like hot cakes} (se vendre comme des petits pains).
\item Questions utiles : \eng{How much is it?} / \eng{How much does it cost?} (Combien ça coûte ?) ; \eng{Can you give me a discount?} (Pouvez-vous me faire une remise ?) ; \eng{That is my last price.} (C'est mon dernier prix.)
\end{itemize}

\textbf{3. Faux amis et pièges}
\begin{center}
\begin{tabularx}{\linewidth}{|X|X|X|}
\hline
\rowcolor{popblueL} Français & English correct & Piège à éviter \\
\hline
magasin & shop, store & \eng{magazine} = revue \\
\hline
monnaie & change & \eng{money} = argent \\
\hline
une affaire & a bargain, a deal & \eng{business} = commerce \\
\hline
marchander & to bargain, to haggle & pas \eng{to marchand} \\
\hline
denrées & foodstuff, food items & pas \eng{teeth things} \\
\hline
\end{tabularx}
\end{center}

\textbf{4. Familles de mots et prononciation}
\begin{itemize}
\item \eng{to buy -- bought -- bought} (acheter ; « baï -- bot ») ; \eng{buyer} (acheteur) ; \eng{to sell -- sold -- sold} (vendre ; « sèl -- sold ») ; \eng{sale} (vente, soldes) ; \eng{seller} (vendeur).
\item \eng{cheap} (« tchipe ») ne se confond pas avec \eng{chip} (« tchip », frite) ; \eng{receipt} se prononce « ri-site » (le \eng{p} est muet).
\end{itemize}

\textbf{5. Méthode express}
\begin{enumerate}
\item Repérer le lieu et les personnes dans le texte ou le dialogue.
\item Classer le vocabulaire en colonnes : lieux / personnes / produits / prix / négociation.
\item Vérifier chaque mot suspect contre la liste des faux amis.
\item Réemployer au moins cinq mots du thème dans une phrase ou un dialogue personnel.
\end{enumerate}
\end{bilan}

\begin{aretenir}[Checklist avant le Bac]
\begin{itemize}
\item[$\square$] Je sais nommer au moins cinq lieux où l'on achète et cinq personnes qui vendent.
\item[$\square$] Je connais le vocabulaire des prix : \eng{price, cost, cheap, expensive, discount, change, receipt}.
\item[$\square$] Je sais conjuguer \eng{buy -- bought -- bought} et \eng{sell -- sold -- sold} sans erreur.
\item[$\square$] Je sais poser la question du prix : \eng{How much is it?} et \eng{How much does it cost?}.
\item[$\square$] Je sais marchander en anglais : \eng{Can you reduce the price? That is too expensive!}.
\item[$\square$] J'évite le faux ami \eng{magazine} pour « magasin » : je dis \eng{shop} ou \eng{store}.
\item[$\square$] Je distingue \eng{money} (argent) et \eng{change} (monnaie rendue).
\item[$\square$] Je prononce correctement \eng{receipt} (« ri-site », \eng{p} muet) et \eng{cheap} (« tchipe »).
\item[$\square$] Je connais trois expressions idiomatiques : \eng{a good bargain}, \eng{value for money}, \eng{to cost a fortune}.
\item[$\square$] Je sais rédiger un court dialogue au marché avec au moins huit mots du thème.
\end{itemize}
\end{aretenir}

\findecours

\end{document}
